\documentclass[a4paper]{ctexbook}
\usepackage{imakeidx}
\makeindex[title=索引,intoc]
\usepackage{tocbibind}
\usepackage[]{subcaption}
\usepackage[]{bicaption}
\usepackage{graphicx}
\usepackage[]{amsmath}
\usepackage[]{amssymb}
\usepackage[]{bigstrut}
\usepackage[dvipsnames,svgnames]{xcolor}
\usepackage[]{multicol}
\usepackage{mhchem}
\usepackage[tone,extra,safe]{tipa}
\usepackage[normalem]{ulem}
% \usepackage[]{ninecolors}
% \NineColors{saturation=low}
% \usepackage[figuresright]{rotating}
\usepackage[]{lscape}
\usepackage{tabularray}
\UseTblrLibrary{booktabs}
% \usepackage[]{paracol}
\usepackage[]{lipsum}
\usepackage[]{zhlipsum}
\usepackage[]{sidenotes}
\usepackage[]{subfiles}
\usepackage[]{siunitx}
\usepackage[]{metalogo}
\usepackage[]{tcolorbox}
\tcbuselibrary{skins}
\tcbuselibrary{breakable}
\tcbuselibrary{documentation}
\usepackage{geometry} 
\geometry{left=2.5cm,right=2.5cm,top=2.5cm,bottom=2.5cm}

\newtcolorbox[auto counter,number within=section]{sample}[2][]{%
arc is angular,breakable,parbox=false,before upper=\par,before lower=\par,drop shadow,skin=bicolor,colbacklower=white,colback=blue!8!white,colframe=blue!46!black, arc=1mm,fonttitle=\bfseries, title={\heiti Sample}~\thetcbcounter: #2,#1}

\newtcolorbox[auto counter,number within=section]{note}[2][]{%
breakable,parbox=false,before upper=\par,before lower=\par,drop shadow,skin={spartan,bicolor},colbacklower=white,colback=red!3!white,colframe=brown!46!black,fonttitle=\bfseries, title={\heiti Note}~\thetcbcounter: #2,#1}

\newtcolorbox[auto counter,number within=section]{task}[2][]{%
colback=Emerald!10,colframe=cyan!40!black,fonttitle=\bfseries,title={\heiti Task}~\thetcbcounter: #2,#1}

\newtcolorbox[auto counter,number within=subsection]{eg}[2][]{%
attach title to upper,coltitle=black,arc=0mm,top=0mm,boxrule=0pt,beforeafter skip balanced=0pt,colback=Thistle!80!red!20!white,colframe=white,fonttitle=\upshape\bfseries,fontupper=\itshape,title={\heiti E.g.}~\thetcbcounter:\ #2,#1}

  
\usepackage[]{adforn}
% \renewcommand{\chaptername}{第 \thechapter{} 章}
% \newcommand{\exercisename}{Exercice}
\newenvironment{problemset}[1][本章练习]{
  \begin{center}
    \phantomsection\addcontentsline{toc}{section}{\texorpdfstring{本章练习}{Exercice}}
    % \markboth{#1}{\rightmark}
    \markright{#1}
    \textcolor{black}{\Large\bfseries\adftripleflourishleft~#1~\adftripleflourishright}
  \end{center}
  \begin{enumerate}}{
  \end{enumerate}}


\usepackage[]{hyperref}
\hypersetup{colorlinks, linkcolor=blue}

\begin{document}


\begin{titlepage}
    \vspace*{\stretch{1}}
    \begin{center}
      {\huge\bfseries Abstract Writing for Medical \\\vspace{3pt} Research Papers}\\[3ex]
      {\huge\bfseries 医学论著英语摘要写作}\\[6.5ex]
      {\large\bfseries Qi Hui, Chen Feina, Guo Haiyan}           \\
      \vspace{4ex}
    %   Thesis  submitted to                    \\[5pt]
      \textit{Fujian Medical University}                \\[2cm]

    %   \includegraphics[width=0.4\textwidth]{fig/fjmulogo.jpg}

    %   in partial fulfilment for the award of the degree of \\[2cm]
    %   \textsc{\Large Doctor of Philosophy}    \\[2ex]
    %   \textsc{\large Mathematics}             \\[12ex]
      \vfill
      Department of Arts and Sciences\\
    %   Address                                 \\
      \vfill
      \date{}
    \end{center}
    \vspace{\stretch{2}}
\end{titlepage}

\frontmatter
\chapter*{前言}\addcontentsline{toc}{chapter}{前言}

近年来，针对医学人才培养和大学外语教育教学，国家层面提出新的指导思想。2012年教育部和原卫生部颁布《关于实施卓越医生教育培养计划的意见》，2018年教育部、国家卫生健康委员会、国家中医药管理局发布《关于加强医教协同实施卓越医生教育培养计划2.0的意见》，指出卓越医师培养在语言层面的要求为“拓展医学生的国际视野，为培养高层次、国际化的医学拔尖创新人才奠定基础”。2017年教育部高等学校大学英语教学指导委员会发布的《大学英语教学指南》提出“学生可以通过学习与专业或未来工作有关的学术英语或职业英语，获得在学术或职业领域进行交流的相关能力”的要求。

根据国家政策的新要求，我们将医学论著英语摘要写作纳入本科生医学英语教学，但是教学过程中发现目前市面上几乎没有合适的基于语料库的医学英语论文摘要写作教材。为解决此问题，我们组建编写团队，依托省级教改项目，进行基于语料库的教材编写。

本教材属于专门用途英语课程教材，可供医学英语摘要写作教学使用，是依托医学论著英语摘要语料库编写的教材。

编写团队成员首先基于语料库语言学理论，建设 115万个词的医学论著英语摘要语料库，为学生医学英语学习提供海量真实的语言数据。为了确保语料的代表性，根据教育部 《学位授予和人才培养学科目录》医学专业一级学科“临床医学”下设的18个二级学科，收集了2016-2018年来自各二级学科影响因子位列前10名的 $2 \sim 3$ 本国际期刊的论著摘要，每个二级学科的摘要篇数基本相同。基于自建语料库，运用体裁分析理论，从宏观层面分析医学论著摘要的篇章结构，并采用定性与定量分析相结合的方法，从微观层面归纳每个语步和语阶的典型语言实现方式。

教材分为六章。第一章作为导入章节，介绍医学论著摘要的定义、重要性及分类；第二章讲述医学论著摘要的语步和语阶的划分；第三至六章分别具体讲授医学论著英语摘要四个语步中每个语阶的语言实现方式。依据在语料库中的出现频率，对每个语步中的常用词、词块等予以排序。教材中所有案例及练习均为该库中的真实语料。

本教材突出体现以下特点： 

1. 教学资源真实，具有科学性 

我们将基于语料库的定量方法与定性观察话语及语境特征的方法相结合。基于海量真实的语料进行研究，强调了经验的重要性，降低了编写者直觉对数据选择的影响，使结果更具说服力。

2. 内容编排合理，具有逻辑性 

教材内容各部分结构平衡协调、浑然一体，能帮助学生了解并掌握医学论著英语摘要的词汇、句法和篇章层面的基本架构。第一章介绍医学论著摘要写作的基本概念；第二章明确摘要中的语步及语阶，四个语步分别为：语步1 (M1) 确定研究领域，语步2 (M2) 描述研究过程，语步3 (M3) 总结研究结果，语步4 (M4) 讨论研究结果；第四章至第六章分别从语用功能、语言实现及案例分析等方面，对每个语步和语阶进行详细讲解。课后练习形式多样，能帮助学生巩固所学知识。

\chapter*{译者序}\addcontentsline{toc}{chapter}{译者序}


本书是由福建医科大学文理艺术学院的齐晖、陈菲娜、郭海燕老师为主编，陈晶为学术秘书，交由复旦大学出版社出版，专供福医大学生使用的英语摘要写作教科书。我本人也是三位老师的学生，纵然课堂生动有趣、干货满满，但苦于同校前辈制作的扫描件观感不佳，笔记整理不便，译者决心要进行文字重排处理。其中自觉原书排版不善之处，皆进行重新编排，以符合译者审美。

原书通本以英文编写，编者似乎意图借此提升我等英语阅读水平，奈何文本中穿插语言学专有名词，初学时疲于翻译、苦不堪言。此外，期末复习期间，全英文本并不利于提升复习效率，故译者对主要文本进行翻译，对照复习。本套重置本将基于该译本进行整理，包含三种排版样式——原文重排版、双语对照版、译文版。

原文重排版和译文版保留了较大的边注区域，便于读者阅读时进行笔记标注，同时译者在此区域添加了部分专有名词的解释。同时，根据排版需要设计了多种类型的盒子对文本区块进行包裹，提升视觉效果的同时，排版的连续性改善了部分区域注意力被引导分散的问题。

本书初稿为个人翻译作品，在此感谢后期参与校对工作的朋友们。若发现纰漏，请在本项目 \href{https://github.com/Neurocane/Abstract-Writing-for-Medical-Research-Papers}{GitHub} 上提交\verb|issue|~\sout{提交了也不一定马上改，咕咕咕}。在此声明，此套图书仅供学习交流使用，请勿用于商业用途或其他领域。本项目工作人员均为爱发电，若发现其他平台账号引流或要求付费购买，都是要被钉在耻辱柱上的。原书内容的著作权归原书作者所有，请多支持正版原著。编辑此书的模板，您可以基于 \href{https://www.latex-project.org/lppl/lppl-1-3c/}{\LaTeX 项目公共许可证 v1.3c} 进行修改和使用。

以上三个版本请学弟学妹们按需取用，20届的老前辈们祝各位期末考试顺利。

\vfill

\begin{flushright}
    \emph{Neurocane}\\
    2023/7/7
\end{flushright}

\newpage

\section*{编译环境}

\begin{itemize}
    \item \textbf{操作系统}： Windows
    \item \textbf{语言}： \LaTeX
    \item \textbf{编译环境}： \XeLaTeX
    \item \textbf{TeX Live版本}： TeX Live 2022
\end{itemize}

\section*{命令环境注释}

\paragraph*{e.g.环境} 标签命名规则：eg:\verb|<语步/语阶缩写>|-\verb|本章节e.g.环境序号|，如：\verb|label={eg:1s1-1}|

\tableofcontents

\mainmatter

\newgeometry{includeall,left=2cm,right=1cm,top=2.5cm,bottom=2.5cm,textwidth=12cm,marginparwidth=5cm}

\chapter{摘要概述}\label{chapter1}
\section{摘要定义}

美国心理学会（APA）规范（2010）指出，摘要是一篇文章内容的简短、全面的总结。根据美国国家标准协会（1979）的规定，摘要是对文件内容的简略、准确的表述，最好是由作者亲自编写，与论文一起出版。一般来说，摘要是文章内容的简明、准确和全面的陈述。它应当是原创的文章，而非论文的节选。

\section{摘要的功能和地位}

摘要是一种独特的体裁，在某种程度上在学术阅读和写作中起着关键的作用。在信息爆炸的时代，学术界每天都有大量的新出版物产生。每位读者都没有切实可行的办法来获取每篇新文章，或者即使可以获取，也无法阅读每篇新出版物。网上发表的摘要简洁而全面，可以方便快捷地获得。因此，阅读摘要可能是任何学术阅读和写作的一个有用的出发点。在这个意义上，摘要是一篇文章中阅读量最大的部分。

摘要至少有三个功能（Huckin，2001）。首先，它可以作为一个独立的小文本，让读者快速了解一项研究的目标、方法、结果和结论，这些都是摘要的主要内容。其次，它可以作为一种筛选工具，让读者充分了解长篇文章是否对他们的需求有很大价值，是否值得进一步阅读。一个好的摘要，在某种程度上可以增加被引用或参考的机会。第三，对于那些选择阅读文章全文的读者来说，摘要可以作为一个预览，创造一个可以指导阅读的解释框架。


\section{摘要的类型}

一般来说，摘要分为两类，指示性和信息性，取决于它们所传达的信息类型。两者之间的一个非典型区别是，指示性摘要被视为论文的纲要，通常较短且简单，而信息性摘要被视为论文的摘要，通常较长且较复杂。

这两种类型的摘要在所包含的内容上也有所不同。指示性摘要通常包括报告或研究的目的、范围和方法，但很少包括结果或结论。阅读指示性摘要不能代替阅读论文，因为没有涵盖所有的关键部分。它在社会科学论文中使用得比较广泛。另一方面，信息性摘要通常包括研究的所有关键部分，如背景、目的、方法、结果和结论。这是广泛用于医学领域的摘要类型。在本书中，我们主要讨论信息性摘要的写作。在以下各章中提到的摘要都是信息性摘要。


\section{信息性摘要的类型}

有两种类型的信息摘要：结构化摘要和非结构化摘要。

如果在每一节文字的开头都使用了标题或标签，那么它就是一个结构化摘要。每一节通常写在一个单独的段落中，但有时各节也写在一个单独或连续的段落中。标题可以是背景、目的、方法、结果、结论等。它们根据不同期刊规定的标准而不同。结构化摘要似乎更受医学相关出版物的青睐。

如果没有使用标题或标签来表示摘要的不同部分，则为非结构化摘要。它总是一个单独的段落。这两种类型的摘要的主要区别在于是否有标题。在非结构化摘要中，项目的内容和顺序与结构化摘要中一样。

期刊规定应使用哪种文体，如果你不确定，请查看作者指南。如果没有提到，请留意你愿意提交和发表论文的期刊所青睐的摘要类型。用占主导地位的风格来写你的摘要。

\begin{sample}[label={myautocounter}]{\heiti}
  \textbf{BACKGROUND} 
  
  In patients with acute heart failure, early intervention with an intravenous vasodilator has been proposed as a therapeutic goal to reduce cardiac-wall stress and, potentially, myocardialinjury, thereby favorably affecting patients' long-term prognosis.

  \textbf{METHODS} 
  
  In this double-blind trial, we randomly assigned 2, 157 patients with acute heart failure to receive a continuous intravenous infusion of either ularitide at a dose of 15 ng per kilogram of body weight per minute or matching placebo for 48 hours, in addition to accepted therapy.
  Treatment was initiated a median of 6 hours after the initial clinical evaluation. The coprimaryd outcomes were death from cardiovascular causes during a median follow-up of 15 months and a hierarchical composite end point that evaluated the initial 48-hour clinical course.
  
  \textbf{RESULTS} 
  
  Death from cardiovascular causes occurred in 236 patients in the ularitide group and 225 patients in the placebo group (21.7\% vs. 21.0\%; hazard ratio, 1.03;96\% confidence interval, 0.85 to 1.25; P=0.75). In the intention-to-treat analysis, there was no significant between-group difference with respect to the hierarchical composite outcome. The ularitide group had greater reductions in systolic blood pressure and in levels of N-terminal pro-brain natriuretic peptide than the placebo group. However, changes in cardiac troponin T levels during the infusion did not differ between the two groups in the 55\% of patients with paired data.

  \textbf{CONCLUSIONS} 
  
  In patients with acute heart failure, ularitide exerted favorable physiological effects (without affecting cardiac troponin levels), but short-term treatment did not affect a clinical composite end point or reduce long-term cardiovascular mortality.

  \begin{flushright}
    ---Effect of Ularitide on Cardiovascular Mortality in Acute Heart Failure.

    \emph{New England Journal of Medicine (2017)}
  \end{flushright}

  
\end{sample}

\begin{sample}[label={myautocounter}]{\heiti}
  \textbf{OBJECTIVE} 
  
  To evaluate the association between the parameters of 24-hour multichannel intraluminal impedance (MII)-pH monitoring and the symptoms or quality of life (QoL) in laryngopharyngeal reflux (LPR) patients.

  \textbf{DESIGN} 
  
  Prospective cohort study without controls.

  \textbf{SETTING} 
  
  University teaching hospital.

  \textbf{METHODS} 
  
  Forty-five LPR patients were selected from subjects who underwent 24-hour MII-pH monitoring and were diagnosed with LPR from September 2014 to May 2015. Reflux Symptom Index (RSI), Health-related Quality of Life (HRQoL), Short Form 12 (SF-12) Survey questionnaires were surveyed. Spearman's correlation was used to analyse the association between the symptoms or QoL and 24-hour MII-pH monitoring.

  \textbf{RESULTS}

  Most parameters in 24-hour MII-pH monitoring showed weak or no correlation with RSI, HRQoL and SF-12. Only number of non-acid reflux events that reached the larynx and pharynx (LPR-non-acid) and number of total reflux events that reached the larynx and pharynx (LPR-total) parameters showed strong correlation with heartburn in RSI (R=0.520, $P<0.001$, R=0.478, $P=0.001$, respectively). Multiple regression analysis showed that there was only one significant regression coefficient between LPR-non-acid and voice/hoarseness portion of HRQoL (b=1.719, $P=0.022$).

  \textbf{CONCLUSION} 

  Most parameters of 24-hour MII-pH monitoring did not reflect subjective symptoms or QoL in patients with LPR.

  \begin{flushright}
    ---Association between 24-hour combined multichannel intraluminal impedance-pH monitoring and symptoms or quality of life in patients with laryngopharyngeal reflux. 
    
    \emph{Clinical Otolaryngology (2017)}
  \end{flushright}

  
\end{sample}

\begin{sample}[label={myautocounter}]{\heiti}
  Due to the high incidence of recurrent squamous cell carcinoma of the head and neck andd the toxicity profile of current salvage regimens, there is a need for tolerable and effective treatment options. We performed a retrospective matched case series to report our experience with recurrent high-risk patients who received capecitabine (CAP) therapy in the adjuvant setting after salvage therapy. The 5-year recurrence-free survival rates for the CAP and control cohorts were 54\% (95\% CI, 0.27\%--0.75\%) and 27\% (95\% CI, 0.09\%--0.50\%), respectively.Multivariable Cox modeling showed a significant improvement in recurrence-free survival in thed CAP cohort (hazard ratio, 0.19; 95\% CI, 0.04--0.92; $P=.0 392$). While this was a respective analysis that could not control for all variables, these exploratory findings offer insights that may inform a prospective study to determine CAP efficacy.
  
  \begin{flushright}
    ---Capecitabine after Surgical Salvage in Recurrent Squamous Cell Carcinoma of Head and Neck. 
    
    \emph{Otolaryngology---Head \& Neck Surgery (2017)}
  \end{flushright}


\end{sample}

\begin{note}[label={myautocounter}]{\heiti 本书使用的语料库}
  本书使用和分析的数据来自一个定制的语料库，其中包括115万个医学研究文章（RA）摘要。医学学科被划分为18个子学科，每个子学科中随机抽取2-3个主要期刊的RA摘要。挑选的期刊都具有相对较高的影响因子，并且每个子学科的文本数量都比较接近 (\autoref{tab:Sub-disciplines and journals in each sub-discipline})。

\end{note}

{\small
  \begin{longtblr}[
      caption = {医学子学科及各子学科期刊},
      label = {tab:Sub-disciplines and journals in each sub-discipline},
  ]{
      width = \textwidth,
      colspec = {X[1,l,h]  X[2,l,h]},
      rowhead = 1, rowfoot = 0, % 每个分页里表头表尾的数量
      % row{odd} = {blue8}, 
      row{even} = {LemonChiffon},
      column{2} = {font=\itshape}
  }
    
    \toprule
    Sub-discipline & Journal \\
    \midrule

    Anesthesiology & {British Journal of Anaesthesia \\ Anesthesiology \\ Anesthesia and Analgesia}\\
    Dermatology & {Journal of American Academy of Dermatology \\ Giornale Italiano di Dermatologia e Venereologia}\\
    Emergency Medicine & {Annals of Emergency Medicine \\ Internal and Emergency Medicine \\ Academic Emergency Medicine}\\
    Geriatrics & {Neurobiology of Aging \\ Aging Cell \\ Age and Ageing}\\
    Internal Medicine & {The New England Journal of Medicine \\ The Lancet \\ JAMA-Journal of the American Medical Association}\\
    Medical Imaging & {The Journal of Nuclear Medicine \\ Investigative Radiology \\ Radiology}\\
    Medical Laboratory & {Clinical Chemistry and Laboratory Medicine \\ Clinical Biochemistry}\\
    Neurology & {The Lancet Neurology \\ Annals of Neurology}\\
    Obstetrics and Gynecology & {Obstetrics \& Gynecology \\ American Journal of Obstetrics \& Gynecology \\ An International Journal of Obstetrics \& Gynecology}\\
    Oncology & {Journal of Clinical Oncology \\ The lancet Oncology}\\
    Ophthalmology & {Ophthalmology \\ American Journal of Ophthalmology \\ Archives of Ophthalmology}\\
    Otolaryngology (ENT) & {Head \& Neck \\ Clinical Otolaryngology \\ Otolaryngology---Head \& Neck Surgery}\\
    Pain Medicine & {The Clinical Journal of Pain \\ Pain Medicine \\ Regional Anesthesia and Pain Medicine}\\
    Pediatrics & {Journal of the American academy of child \& Adolescent psychiatry \\ Pediatrics \\ JAMA pediatrics}\\
    Physical medicine and rehabilitation & {Neurorehabilitation and neural repair \\ Journal of fluency disorders}\\
    Psychiatry & {Molecular psychiatry \\ The American journal of psychiatry \\ JAMA psychiatry}\\
    Sports medicine & {Medicine and Science in Sports and Exercise \\ Sports Medicine \\ The American Journal of Sports Medicine}\\
    Surgery & {Annals of Surgery \\ American Journal of Transplantation \\ Journal of Neurology, Neurosurgery \& Psychiatry}\\

    \bottomrule

  \end{longtblr}
  }

  \begin{task}[label={myautocounter}]{\heiti 语料库任务}
    你能为自己构建一个包含 100, 000 字符的语料库吗？
  \end{task}

  \begin{note}[label={myautocounter}]{\heiti 本书使用的语料库}
    在本书中，“语料库任务”旨在提高您基于语料库法探索医学学术论文摘要语言构成的能力。这些任务中的大多数可能需要使用诸如AntConc、WordSmith在内的专业软件。
  \end{note}

  \section{词汇表}

  {\small
  \begin{longtblr}[
      caption = {第一章词汇表},
      label = {tab:Glossary of Chapter 1},
  ]{
      width = \textwidth,
      colspec = {X[1,l,h]  X[1,l,h]  X[3,l,h]},
      rowhead = 1, rowfoot = 0, % 每个分页里表头表尾的数量
      % row{odd} = {blue8}, 
      row{even} = {azure9},
  }
      
  \toprule
  \textbf{WORDS} & \textbf{MEANING} & \textbf{MEANING or EXAMPLE}\\
  \midrule

  \textbf{excerpt}/\textipa{"eks3:pt}/ & \emph{v.} 摘录；引用  & If a long piece of writing or music is excerpted, short pieces from it are printed or played on their own. \\
  \textbf{genre}/\textipa{"ZA:nr@}/ & \emph{n.} 体裁 & a particular type of art, writing, music etc, which has certain features that all examples of this type share.\\
  {\textbf{mandate} \\ /\textipa{"m\ae ndeIt}/} & \emph{v.} 授权；强制执行；委托办理 & to tell someone that they must do a particular thing. \\
  {\textbf{methodology} \\ /\textipa{""meT@"d6l@\textdyoghlig i}/} & \emph{n.} 方法学 & a set of methods and principles used to perform a particular activity. \\
  \textbf{opt}/\textipa{6pt}/ & \emph{v.} 选择；挑选 & to choose one thing or do one thing instead of another \\
  \textbf{pivotal}/\textipa{"pIv@tl}/ & \emph{adj.} 关键性的；核心的 & more important than anything else in a situation or system. \\
  
  \bottomrule

  \end{longtblr}
  }

  \begin{problemset}
    \item \textbf{Identify whether the following abstracts are structured or unstructured and tell the reasons.}
    
    \textbf{Abstract 1}

    \hspace*{2em}Objectives: The aim of this study was to analyze changes in health care utilization and cost among a sample of highly impaired children and adolescents who sought a 3-week intensive interdisciplinary pain treatment (IPT).

    \hspace*{2em}Materials and Methods: Claims data from 7 statutory health insurance companies were analyzed for 65 children and adolescents who sought IIPT at the German Paediatric Pain Centre. The annual health care utilization and cost were determined for the following 4 areas: outpatient care, inpatient care, medications, and remedies and aids. We analyzed the changes in resource utilization in the year before (pre\_1 y) IPT and in the subsequent year (post\_1 y).
    
    \hspace*{2em}Results: Within the first year after IPT, overall health care costs did not decrease significantly. However, the pattern of health care utilization changed. First, significantly more children and adolescents started outpatient psychotherapy ($P=0.001$). Second, the number of hospitalized children decreased significantly from 1-year pre to 1-year post ($P=0.001$). Accordingly, there were significantly fewer hospitalizations for primary chronic pain disorders at 1-year post ($P<0.001$). The prescription of nonopioids, co-analgesics and opioids was significantly reduced from 1-year pre to 1-year post (all $P<0. 013$).
    
    \hspace*{2em}Discussion: The present results indicate that the health care costs of children and adolescents with severe chronic pain disorders do not significantly decrease 1 year after IPT; however, the treatment becomes more goal-focused. Differential diagnosis measures and nonindicated therapeutic interventions decreased, and more indicated interventions, such as psychotherapy, were used. Future research is needed to investigate the economic long-term changes after IPT.
    
    \begin{flushright}
      ---Health Care Utilization and Cost in Children and Adolescents with Chronic Pain: Analysis of Health Care Claims Data 1 Year Before and After Intensive Interdisciplinary Pain Treatment. 
      
      \emph{The Clinical Journal of Pain (2017)}
    \end{flushright}

    \textbf{Abstract 2}

    \hspace*{2em}Previous studies of brain structure in Tourette syndrome (TS) have produced mixed results, and most had modest sample sizes. In the present multicenter study, we used structural magnetic resonance imaging (MRI) to compare 103 children and adolescents with TS to a well-matched group of 103 children without tics. We applied voxel-based morphometry methods to test gray matter (GM) and white matter (WM) volume differences between diagnostic groups, accounting for MRI scanner and sequence, age, sex and total GM$+$WM volume. The TS group demonstrated lower WM volume bilaterally in orbital and medial prefrontal cortex, and greater GM volume in posterior thalamus, hypothalamus and midbrain. These results demonstrate evidence for abnormal brain structure in children and youth with TS, consistent with and extending previous findings, and they point to new target regions and avenues of study in TS. For example, as orbital cortex is reciprocally connected with hypothalamus, structural abnormalities in these regions may relate to abnormal decision making, reinforcement learning or somatic processing in TS.

    \begin{flushright}
      ---Brain structure in pediatric Tourette syndrome. 

      \emph{Molecular Psychiatry (2017)}
    \end{flushright}

    \textbf{Abstract 3}

    \hspace*{2em}Objective: To assess the feasibility of detecting signature volatile organic compounds in the breath of patients with oral squamous cell carcinoma.
    
    \hspace*{2em}Study Design: Prospective cohort pilot study.
    
    \hspace*{2em}Setting: University hospital.
    
    \hspace*{2em}Subjects and Methods: Using gas chromatography and mass spectrometry, emitted volatile organic compounds in the breath of patients before and after curative surgery (n=10) were compared with those of healthy subjects (n=4). It was hypothesized that certain volatile organic compounds disappear after surgical therapy. A characteristic signature of these compounds for diseased patients was compiled and validated.
    
    \hspace*{2em}Results: Breath analyses revealed 125 volatile organic compounds in patients with oral cancer. A signature of 8 compounds that were characteristic for patients with oral cancer could be detected: 3 from this group presented were absent after surgery.
    
    \hspace*{2em}Conclusion: The presented results confirmed the hypothesis of an absence of cancer-associated volatile organic compounds in the breath after therapy. In this pilot study, we proved the feasibility of this test approach. Further studies should be initiated to establish protocols for usage in a clinical setting.
    
    \begin{flushright}
      ---Volatile Organic Compounds in the Breath of Oral Squamous Cell Carcinoma Patients: A Pilot Study. 

      \emph{Otolaryngology Head and Neck Surgery (2017)}
    \end{flushright}
   
  \end{problemset}


\chapter{语步和语阶辨认}\label{chapter2}

\section{语步辨认}

医学RA的摘要是由一些语步组成的，这些语步共同实现其交流的目的。在这个意义上，语步可定义为 “文本中执行特定交流功能的部分” (Kanoksilapatham, 2007, p.23)。

大多数医学论文摘要都采用了四步走的方案，第一步（M1）确定研究领域，第二步（M2）描述研究过程，第三步（M3）总结主要成果，第四步（M4）得出结论。这些都是医学研究报告摘要中的常规语步。

在结构化的摘要中，M2、M3和M4可以通过方法(method)、结果(result)和结论(conclusion)的标题轻松识别。虽然有些摘要中没有“方法”的具体标题，但M2被细分为几个语阶，可以清楚地识别。它们可以被标记为 “实验设计(design)”、“实验环境设置(setting)”、“参与者(participants)”、“干预(interventions)”、“主要观察指标(main outcome measures)”，等等。

当涉及到M1时，标有“目标(objective)”或“背景(background)”或两者皆有的章节都属于M1部分的内容。在一些结构化的摘要中，“背景” 和“目标”都被标示出来，有些只有“目标”，而有些只有“背景”。那些标有“目标”或“背景”的部分，在实际内容上通常由这两个部分组成。此外，在实践过程中，目标和背景密切相关，通常被视为一个整体，以实现确定研究领域这一特定交流目的。

对于结构化的摘要，可以通过标题这个信号词之一来识别四个语步。对于没有标题或标签的非结构化摘要，需要人工识别四种语步。其他的信号词也可以帮助识别这些语步， 这些信号词将在\hyperref[chapter3]{第三章}到\hyperref[chapter6]{第六章}作相应的说明。

\begin{sample}[label={myautocounter}]{\heiti}

\textbf{(M1)}

\textbf{BACKGROUND }

it is still equivocal whether there is a potential role of late-life physical activity in ameliorating the challenges of increasing healthcare expenditure due to the consequence of global population ageing.

\textbf{OBJECTIVE} 

this study aimed to examine the prospective association between physical activity and subsequent hospital care utilisation in older adults and to explore the optimal dose of physical activity required to reduce hospital care utilisation.

\textbf{(M2)}

\textbf{DESIGN} 

this was a prospective cohort study based on the data from the Taiwan 2005 National Health Interview Survey, which were linked to the 2005-12 claims data from the National Health Insurance system.

\textbf{PARTICIPANTS}

1,760 older adults aged 65 or more.

\textbf{METHODS} 

the frequency, duration and intensity for physical activity were assessed, and total physical activity energy expenditure was estimated. The average annualised hospital care utilisation for the period 2006 through 2012, including number of hospitalisations, number of days in hospital and the costs of hospitalisation, were calculated.

\textbf{(M3)}

\textbf{RESULTS }

older adults engaging in at least moderate volume of physical activity ($\geqslant 1,000$ kcal/week) experienced fewer subsequent hospital admissions and fewer days in hospital than did sedentary individuals, after adjusting for covariates. Trends for reduced hospitalisation costs were also found. These associations persisted in sensitivity analyses, including tests of reverse causation.

\textbf{(M4)}

\textbf{CONCLUSION}

this study has provided evidence that older adults who are at least moderately active may minimise utilisation of hospital care services. The findings highlight the importance of maintaining a physically active lifestyle in later life.

\begin{flushright}
  ---Prospective association between late-life physical activity and hospital care utilisation: a 7-year nationwide follow-up study. \emph{Age and Aging (2017)}
\end{flushright}

\end{sample}

\begin{sample}[label={myautocounter}]{\heiti}

\textbf{(M1)}

\textbf{PURPOSE }

Stuttering can trigger anxiety and other psychological and emotional reactions, and limit participation in society. It is possible that psychological counseling could enhance stuttering treatment outcomes; however, little is known about how clients view such counseling. The purpose of this study was to gain an understanding of clients' experiences with, and perceptions of, a psychological counseling service that was offered as an optional adjunct to speech therapy for stuttering.

\textbf{(M2)}

\textbf{METHOD }

Nine individuals who stutter (13--38 years old) participated in semi-structured interviews. Six participants had taken part in psychological counseling; three participants did not do so. Interview data were analyzed using grounded theory as a guiding framework.

\textbf{(M3)}

\textbf{RESULTS }

Four thematic clusters emerged from participants' accounts: insights into personal decision-making, why others may not participate in counseling, psychological counseling as a worthwhile part of therapy, and counseling as a necessary component in a stuttering treatment program.

\textbf{(M4)}

\textbf{CONCLUSION}

In addition to experiencing barriers and facilitators to help-seeking that are reported in related fields, participants accounts also revealed novel facilitators (i.e., a `why not' mentality and the importance of having a pre-existing relationship with the clinician who offered the service) and barriers (i. e., viewing the service as a `limited resource,' and, the overwhelming nature of intensive stuttering treatment programs). Findings suggest that clients value the option to access psychological counseling with trained mental health professionals to support the stuttering treatment provided by speech-language pathologists. Participants made recommendations for the integration of psychological counseling into stuttering treatment programs.

\begin{flushright}
  ---Psychological counseling as an adjunct to stuttering treatment: Clients' experiences and perceptions. 
  
  \emph{Journal of Fluency Disorder (2017)}
\end{flushright}


\end{sample}


\begin{sample}[label={myautocounter}]{\heiti}
\textbf{(M1)}

Genetic and neuroimaging research has identified neurobiological correlates of obesity. However, evidence for an integrated model of genetic risk and brain structural alterations in the pathophysiology of obesity is still absent.

\textbf{(M2)}

Here we investigated the relationship between polygenic risk for obesity, gray matter structure and body mass index (BMI) by the use of univariate and multivariate analyses in two large, independent cohorts ($n=330$ and $n=347$).

\textbf{(M3)}

Higher BMI and higher polygenic risk for obesity were significantly associated with medial prefrontal gray matter decrease, and prefrontal gray matter was further shown to significantly mediate the effect of polygenic risk for obesity on BMI in both samples.

\textbf{(M4)}

Building on this, the successful individualized prediction of BMI by means of multivariate pattern classification algorithms trained on whole-brain imaging data and external validations in the second cohort points to potential clinical applications of this imaging trait marker.

\begin{flushright}
  ---Prefrontal gray matter volume mediates genetic risks for obesity. \emph{Molecular Psychiatry (2017)}
\end{flushright}

\end{sample}

\section{语阶辨认}

在有明显层次结构的体裁中，语步通常由语阶或子语步组成，它们是语步的下级单位。在医学文摘中，有些语阶属常规语阶，有些是可选语阶。语阶的类型和频率也显示了作者的修辞目的。 信号词对识别这些语阶很有帮助，本书将在\hyperref[chapter3]{第三章}到\hyperref[chapter6]{第六章}中作相应说明。

语料库中不同期刊的摘要是随机抽取的，每个子学科有6篇文章，共108篇文章。通过人工识别，构成每个语步的语阶及其对应功能被确定下来（\autoref{tab:Move/Step Scheme of Medical RA Abstracts}）。如果该语阶在80\%以上的文章中使用，则被认为是常规的，否则是可选的。

{\small
\begin{longtblr}[
    caption = {Move/Step Scheme of Medical RA Abstracts},
    label = {tab:Move/Step Scheme of Medical RA Abstracts},
]{
    width = \textwidth,
    colspec = {X[1,l,h]  X[1,c,h]  X[3,l,h] X[1,c,h]},
    rowhead = 1, rowfoot = 0, % 每个分页里表头表尾的数量
    % row{odd} = {blue8}, 
    % row{even} = {azure9},
    row{2-6} = {azure9},
    row{7-14} = {white},
    row{15-17} = {azure9},
    row{18-22} = {white},
}
    
\toprule
\textbf{Move/Step} & \textbf{Move/Step Abbr.} & \textbf{Communicative functions} & \textbf{Percentage}\\
\midrule
Move1 & M1 & Creating a research territory/space &  \\
\hspace*{1ex}Move1Step1 & 1S1 & Presenting current knowledge or relevant information established by previous studies & 56.5\% \\
\hspace*{1ex}Move1Step2 & 1S2 & Establishing a niche/problem\sidenotemark[1] & 41.7\% \\
\hspace*{1ex}Move1Step3a & 1S3a & Indicating main purposes & 83.3\% \\
\hspace*{1ex}Move1Step3b & 1S3b & Raising hypotheses & 2.8\% \\
Move2 & M2 & Describing research process &  \\
\hspace*{1ex}Move2Step1 & 2S1 & Reporting on medical ethics review & 2.8\% \\
\hspace*{1ex}Move2Step2 & 2S2 & Explaining briefly research design & 54.6\% \\
\hspace*{1ex}Move2Step3 & 2S3 & Describing subjects or data and their selection criteria & 92.6\% \\
\hspace*{1ex}Move2Step4 & 2S4 & Describing experimemal procedure, such as interventions, examinations, etc. & 60.2\% \\
\hspace*{1ex}Move2Step5 & 2S5 & Describing main outcomes and their measures & 93.5\% \\
\hspace*{1ex}Move2Step6 & 2S6 & Describing data analysis methods & 10.2\% \\
\hspace*{1ex}Move2Step7 & 2S7 & Reporting on registration information & 9.3\% \\
Move3 & M3 & Summarizing results &  \\
\hspace*{1ex}Move3Step1 & 3S1 & Providing information on valid samples & 43.5\% \\
\hspace*{1ex}Move3Step2 & 3S2 & Illustrating overall observation or main results & 100\% \\
Move4 & M4 & Drawing conclusions &  \\
\hspace*{1ex}Move4Step1 & 4S1 & Reiterating pivotal results & 16.7\% \\
\hspace*{1ex}Move4Step2 & 4S2 & Indicating limitations & 3.7\% \\
\hspace*{1ex}Move4Step3 & 4S3 & Stating the significance of the results & 99.1\% \\
\hspace*{1ex}Move4Step4 & 4S4 & Predicting future studies & 19.4\% \\

\bottomrule

\end{longtblr}
}

\sidenotetext[1]{{\small niche 原本指法国天主教徒房屋墙壁上预留用于放置圣母玛利亚的神龛。20世纪80年代，该词以“利基市场 (Niche Market) ”被美国商学家引入市场营销领域，意指那些被市场中的统治者/有绝对优势的企业忽略的某些细分市场。英文论文写作中，niche引申为被主流研究忽视的问题，即现有研究的缺陷、不足，指向该论文意图解决的问题。}}

\begin{sample}[label={myautocounter}]{\heiti}
  \textbf{BACKGROUND}

  \textbf{(1S1)} Graded exercise therapy is an effective and safe treatment for chronic fatigue syndrome, \textbf{(1S2)} but it is therapist intensive and availability is limited. \textbf{(1S3a)} We aimed to test the efficacy and safety of graded exercise delivered as guided self-help.
  
  \textbf{METHODS}

  \textbf{(2S2)} In this pragmatic randomised controlled trial, \textbf{(2S3)} we recruited adult patients (18 years and older) who met the UK National Institute for Health and Care Excellence criteria for chronic fatigue syndrome from two secondary-care clinics in the UK. \textbf{(2S4)} Patients were randomly assigned to receive specialist medical care (SMC) alone (control group) or SMC with additional guided graded exercise self-help (GES). Block randomisation (randomly varying block sizes) was done at the level of the individual with a computer-generated sequence and was stratified by centre, depression score, and severity of physical disability. Patients and physiotherapists were necessarily unmasked from intervention assignment; the statistician was masked from intervention assignment. SMC was delivered by specialist doctors but was not standardised; GES consisted of a self-help booklet describing a six-step graded exercise programme that would take roughly 12 weeks to complete, and up to four guidance sessions with a physiotherapist over 8 weeks (maximum 90 min in total). \textbf{(2S5)} Primary outcomes were fatigue (measured by the Chalder Fatigue Questionnaire) and physical function (assessed by the Short Form-36 physical function subscale); both were self-rated by patients at 12 weeks after randomisation and analysed in all randomised patients with outcome data at follow-up (ie, by modified intention to treat). We recorded adverse events, including serious adverse reactions to trial interventions.\textbf{ (2S6) }We used multiple linear regression analysis to compare SMC with GES, adjusting for baseline and stratification factors.\textbf{ (2S7) }This trial is registered at ISRCTN, number ISRCTN22975026.
  
  \textbf{FINDINGS}

  \textbf{(3S1)} Between May 15,2012, and Dec 24,2014, we recruited 211 eligible patients, of whom 107 were assigned to the GES group and 104 to the control group. \textbf{(3S2) }At 12 weeks, compared with the control group, mean fatigue score was 19.1 (SD-7.6) in the GES group and
  22.9(6.9) in the control group (adjusted difference-4.2 points, 95\% CI -6.1 to -2.3, $P<0.0 001$; effect size 0.53) and mean physical function score was 55.7(23.3) in the GES groupand 50. 8(25.3) in the control group (adjusted difference 6.3 points, 1.8 to 10.8, P=0.006;
  0.20). No serious adverse reactions were recorded and other safety measures did not differ between the groups, after allowing for missing data.

  \textbf{INTERPRETATION}

  \textbf{(4S2)} GES is a safe intervention that might reduce fatigue and, to a lesser extent, physical disability for patients with chronic fatigue syndrome. \textbf{(4S4)} These findings need confirmation and extension to other health-care settings.
  
  \begin{flushright}
    ---Guided graded exercise self-help plus specialist medical care versus specialist medical care alone for chronic fatigue syndrome (GETSET): a pragmatic randomised controlled trial. 
    
    \emph{The Lancet (2017)}
  \end{flushright}
  
\end{sample}

\section{词汇表}

{\small
\begin{longtblr}[
    caption = {第二章词汇表},
    label = {tab:Glossary of Chapter 2},
    note{a} = {英文论文中指代当前文献中的差距、问题或缺陷。即现有研究尚未解决的部分。},
]{
    width = \textwidth,
    colspec = {X[1,l,h]  X[1,l,h]  X[3,l,h]},
    rowhead = 1, rowfoot = 0, % 每个分页里表头表尾的数量
    % row{odd} = {blue8}, 
    row{even} = {azure9},
}
    
\toprule
\textbf{WORDS} & \textbf{MEANING} & \textbf{MEANING or EXAMPLE}\\
\midrule

\textbf{conventional}/\textipa{k@n"venS@nl}/ & \emph{adj.} 传统的；常规的 & a conventional method, product, practice, etc. has been used for a long time and is considered the usual type \\
\textbf{ethics}/\textipa{"eTiks}/ & \emph{n.} 伦理标准 & [plural] moral rules or principles of behaviour for deciding what is right and wrong \\
\textbf{intervention}/\textipa{""Int@"venSn}/ & \emph{n.} 干预；介入 & an action or ministration that produces an effect or is intended to alter the course of a pathologic process \\
\textbf{niche}/\textipa{ni:S}/ & \emph{n.} 利基\TblrNote{a}；生态位；微环境  & a gap in the previous research \\
\textbf{reiterate}/\textipa{ri:"It@reIt}/ & \emph{v.} 重申  & to say something again, usually in order to emphasize it \\

\bottomrule

\end{longtblr}
}

\begin{problemset}
  \item \textbf{Identify the four moves in the abstracts.}
  
  \textbf{Abstract 1}

  \hspace*{2em}Objective: To evaluate the association between the parameters of 24-hour multichannel intraluminal impedance (MII)-pH monitoring and the symptoms or quality of life (QoL) in laryngopharyngeal reflux (LPR) patients.
  
  \hspace*{2em}Design: Prospective cohort study without controls.
  
  \hspace*{2em}Setting: University teaching hospital.
  
  \hspace*{2em}Methods: Forty-five LPR patients were selected from subjects who underwent 24-hour MII-
  pH monitoring and were diagnosed with LPR from September 2014 to May 2015. Reflux Symptom Index (RSI), Health-related Quality of Life (HRQoL), Short Form 12 (SF-12) Survey questionnaires were surveyed. Spearman's correlation was used to analyse the association between the symptoms or QoL and 24-hour MII-pH monitoring.
  
  \hspace*{2em}Results: Most parameters in 24-hour MII-pH monitoring showed weak or no correlation with RSI, HRQoL and SF-12. Only number of non-acid reflux events that reached the larynxand pharynx (LPR-non-acid) and number of total reflux events that reached the larynx and pharynx (LPR-total) parameters showed strong correlation with heartburn in RSI (R=0.520, $P<0$.001, R=0.478, P=0.001, respectively). Multiple regression analysis showed that there was only one significant regression coefficient between LPR-non-acid and voice/hoarseness portion of HRQoL (b=1.719, P=0.022).
  
  \hspace*{2em}Conclusion: Most parameters of 24-hour MII-pH monitoring did not reflect subjective symptoms or QoL in patients with LPR.
  
  \begin{flushright}
    ---Association between 24-hour combined multichannel intraluminal impedance-pH monitoring and symptoms or quality of life in patients with laryngopharyngeal reflux. \emph{Clinical Otolaryngology (2017)}
  \end{flushright}
  
  \textbf{Abstract 2}
  
  \hspace*{2em}Skeletal muscle mitochondrial oxidative capacity declines with age and negatively affects walking performance, but the mechanism for this association is not fully clear. We tested the hypothesis that impaired oxidative capacity affects muscle performance and, through this mechanism, has a negative effect on walking speed. Muscle mitochondrial oxidative capacity was measured by in vivo phosphorus magnetic resonance spectroscopy as the postexercise phosphocreatine resynthesis rate, Kpc, in 326 participants (154 men), aged 24-97 years (mean 71), in the Baltimore Longitudinal Study of Aging. Muscle strength and quality were determined by knee extension isokinetic strength, and the ratio of knee extension strength to thigh muscle cross-sectional area derived from computed topography, respectively. In multivariate linear regression analyses, kpc, was associated with muscle strength ($\beta=0. 140, P=0.007$) and muscle quality ($\beta=0.127, P=0.022$), independent of age, sex, height, and weight; muscle strength was also a significant independent correlate of walking speed ($P<0. 02$ for all tasks) and in a formal mediation analysis significantly attenuated the association between kpc and three of four walking tasks (18\%-29\% reduction in $\beta$ for kpcr). This is the first demonstration in human adults that mitochondrial function affects muscle strength and that inefficiency in muscle bioenergetics partially accounts for differences in mobility through this mechanism.

  \begin{flushright}
    ---Muscle strength mediates the relationship between mitochondrial energetics and walking performance. \emph{Aging Cell (2017)}
  \end{flushright}

  \textbf{Abstract 3}

  \hspace*{2em}Propacetamol, a water-soluble prodrug form of paracetamol, is hydrolyzed by esterase to generate paracetamol in the blood. Each gram of propacetamol is equal to 0.5 g of paracetamol. It has been reported to cause hypotension in critically ill patients with a fever. We aimed to investigate the hemodynamic effects of propacetamol for the control of fever in patients with diverse severities of illness who were managed in the emergency department (ED). We also aimed to identify clinical factors related to significant hemodynamic alterations in ED patients. This was a retrospective study of 1507 ED patients who received propacetamol. Significant hemodynamic alterations were defined as systolic blood pressure (SBP) <90 mmHg or diastolic blood pressure (DBP) <60 mmHg, or a drop in SBP >30 mmHg, which required treatments with a bolus of fluid or vasopressor administration. Postinfusion SBP and DBP were significantly lower than the preinfusion SBP and DBP. A clinically significant drop in BP occurred in 162 (10.7\%) patients, and interventions were necessary. Among the predictors assessed, congestive heart failure (OR 6.21,95\% CI 2.67--14.45) and chills (OR 3.10,95\% CI 2.04--4.70) were independent factors for a significant hemodynamic change. Administrationof propacetamol can provoke a reduction in BP in ED patients. This reduction was clinically significant for 10\% of infusions. Clinicians should be aware of this potential deleterious effect, especially in patients with congestive heart failure or who experience chills prior to the administration of propacetamol.

  \begin{flushright}
    ---Clinically significant hemodynamic alterations after propacetamol injection in the emergency department: prevalence and risk factors. \emph{Internal and Emergency Medicine (2016)}
  \end{flushright}

  \item \textbf{Reorder the four moves in the abstracts.}
  
  \textbf{Abstract 1}

  \hspace*{2em}Methods: A dual-virus tracing strategy combining retroviral birthdating with rabies virus-
  mediated putative retrograde trans-synaptic tracing was used to identify and compare presynaptic inputs onto adult-born and early-born DGCs in the rat pilocarpine model of mTLE.

  \hspace*{2em}---Move \uline{\hspace*{3em}}
  
  \hspace*{2em}Objective: To understand how monosynaptic inputs onto adult-born dentate granule cells (DGCs) are altered in experimental mesial temporal lobe epilepsy (mTLE) and whether their integration differs from early-born DGCs that are mature at the time of epileptogenesis.

  \hspace*{2em}---Move \uline{\hspace*{3em}}
  
  \hspace*{2em}Interpretation: These data support the presence of substantial hippocampal circuit remodeling after an epileptogenic insult that generates prominent excitatory monosynaptic inputs, both local recurrent and widespread feedback loops, onto DGCs. Both adult-born and early-bom DGCs are targets of new inputs from other DGCs as well as from CA3 and CA1 pyramidal cells after pilocarpine treatment, changes that likely contribute to epileptogenesis in experimental mTLE.
  
  \hspace*{2em}---Move \uline{\hspace*{3em}}

  \hspace*{2em}Results: Our results demonstrate that hilar ectopic DGCs preferentially synapse onto adult-born DGCs after pilocarpine-induced status epilepticus (SE), whereas normotopic DGCs synapse onto both adult-born and early-born DGCs. We also find that parvalbumin- and somatostatin-interneuron inputs are greatly diminished onto early-born DGCs after SE. However, somatostatin-interneuron inputs onto adult-born DGCs are maintained, likely due topreferential sprouting. Intriguingly, CA3 pyramidal cell backprojections that specifically target adult-born DGCs arise in the epileptic brain, whereas axons of interneurons and pyramidal cells in CA1 appear to sprout across the hippocampal fssure to preferentilly synapse onto early-bomDGCs.

  \hspace*{2em}---Move \uline{\hspace*{3em}}

  \begin{flushright}
    ---Rabiestracing of birthdated dentae granule cells in rat temporal lobe epilepsy. \emph{Annals of Neurology(2017)}
  \end{flushright}

  \textbf{Abstract 2}
  
  \hspace*{2em}We identified a mssense Asn396Ser mutation (rs77960347) in the endothelial lipase (LIG) gene, occurring with an allele frequency of 1\% in the general population, which was significantly associated with depressive symptoms (P-value $= 5.2 \times 10--08, \beta=7.2$). Replication in three independent data sets ($N= 3612$) confirmed the association of Asn396Ser (P-value = $7.1 \times 10-03$, $\beta = 2.55$) with depressive symptoms.

  \hspace*{2em}---Move \uline{\hspace*{3em}}

  \hspace*{2em}Despite a substantial genetic component, efforts to identify common genetic variation underlying depression have largely been unsuccessful. In the current study we aimed to identify rare genetic variants that might have large effects on depression in the general population.

  \hspace*{2em}---Move \uline{\hspace*{3em}}

  \hspace*{2em}Using high-coverage exome-sequencing, we studied the exonic variants in 1 265 individuals from the Rotterdam study (RS), who were assessed for depressive symptoms.

  \hspace*{2em}---Move \uline{\hspace*{3em}}

  \hspace*{2em}LIPG is predicted to have enzymatic function in steroid biosynthesis, cholesterol biosynthesis and thyroid hormone metabolic processes. The Asn396Ser variant is predicted tod have a damaging effect on the function of LIPG. Within the discovery population, carriers also showed an increased burden of white matter lesions (P-value=$3.3 \times 10-02$) and a higher risk of Alzheimer's disease (odds ratio=2.01; P-value = 2.8x 10-02) compared with the noncarriers. Together, these findings implicate the Asn396Ser variant of LIPG in the pathogenesis of depressive symptoms in the general population.

  \hspace*{2em}---Move \uline{\hspace*{3em}}

  \begin{flushright}
    ---Exome-sequencing in a large population-based study reveals a rare Asn396Ser variant in the LIPG gene associated with depressive symptoms. \emph{Molecular Psychiatry (2017)}
  \end{flushright}

  \textbf{Abstract 3}

  \hspace*{2em}Conclusion: This study suggests that hearing problems in later life could increase the risk of having difficulties performing IADLs, which include more complex everyday tasks such as shopping and light housework. However, further studies are needed to determine the associations observed including the underlying pathways.
  
  \hspace*{2em}---Move \uline{\hspace*{3em}} 
  
  \hspace*{2em}Methods: Data were collected on self-reported hearing impairment including hearing aid use, and disability assessed as mobility limitations (problems walking/taking stairs), difficulties with activities of daily living (ADL) and instrumental ADL (IADL). Mortality data were obtained from the National Health Service register.
  
  \hspace*{2em}---Move \uline{\hspace*{3em}}
  
  \hspace*{2em}Background and objective: Hearing impairment is common in older adults and has been implicated in the risk of disability and mortality. We examined the association between hearing impairment and risk of incident disability and all-cause mortality.

  \hspace*{2em}---Move \uline{\hspace*{3em}}
  
  \hspace*{2em}Results: Among 3,981 men, 1,074(27\%) reported hearing impairment. Compared with men with no hearing impairment, men who could hear and used a hearing aid, and men who could not hear despite a hearing aid had increased risks of IADL difficulties (age-adjusted OR 1.86, 95\% CI 1.29--2.70; OR 2.74, 95\% CI 1.53--4.93, respectively). The associations remained after further adjustment for covariates including social class, lifestyle factors, comorbidities and social engagement. Associations of hearing impairment with incident mobility limitations, incident ADL difficulties and all-cause mortality were attenuated on adjustment for covariates.
  
  \hspace*{2em}---Move \uline{\hspace*{3em}}
  
  \begin{flushright}
    ---Hearing impairment and incident disability and all-cause mortality in older British community-dwelling men. \emph{Age and Ageing (2016)}
  \end{flushright}
  
\end{problemset}


\chapter{语步一}\label{chapter3}

语步一 (M1) 是摘要的第一部分。它为研究确定了一个研究领域或空间。

\section{M1中的语阶}

M1通常包括三个语阶，每个语阶都有明确的交流目的。

1S1 介绍当前的知识或由以前的研究确立的相关信息。

1S2 确立利基/问题

1S3a 指明主要目的

1S3b 提出假设

通常在最后一步，要么提出目标，要么提出假设，而只有少数摘要同时包括这两个方面。

\begin{sample}[label={myautocounter}]{\heiti}
  Senescent cells are present in premalignant lesions and sites of tissue damage and accumulate in tissues with age. In vivo identification, quantification and characterization of senescent cells are challenging tasks that limit our understanding of the role of senescent cells in diseases and aging. Here, we present a new way to precisely quantify and identify senescent cells in tissues on a single-cell basis.

  \begin{flushright}
    ---Quantitative identification of senescent cells in aging and disease. \emph{Aging Cell (2017)}
  \end{flushright}

  \tcblower

  \noindent \textbf{STEP IDENTIFICATION}

  \vspace*{10pt}
  {\small\noindent
  \begin{tblr}{colspec={X[1,c]X[5,l]},cell{even}{1,2} = {azure9}}
    \toprule
    \textbf{Step} & \textbf{Sample} \\ 
    \midrule
    
    1S1 & Senescent cells are present in premalignant lesions and sites of tissue damage and accumulate in tissues with age.  \\
    1S2 & In vivo identification, quantification and characterization of senescent cells are challenging tasks that limit our understanding of the role of senescent cells in diseases and aging. \\
    1S3a & Here, we present a new way to precisely quantify and identify senescent cells in tissues on a single-cell basis. \\
      
    \bottomrule
  \end{tblr}
  }

  \noindent \textbf{ANALYSIS} 
  
  This is a typical example of M1 involving three steps. The first sentence informs the readers of the background information related to the research. The second sentence, where the word ``no'' can be seen as a signal for 1S2, identifies the problem-the lack of a standard. The last sentence, with the subject ``aim'' and the to-infinitive, shows the research objective.
\end{sample}



\begin{sample}[label={myautocounter}]{\heiti}

  \textbf{OBJECTIVE} 
  
  Seizures are more frequent in patients with Alzheimer's disease (AD) and can hasten cognitive decline. However, the incidence of subclinical epileptiform activity in AD and its consequences are unknown. Motivated by results from animal studies, we hypothesized higherthan expected rates of subclinical epileptiform activity in AD with deleterious effects on cognition.

  
  \begin{flushright}
    ---Incidence and impact of subclinical epileptiform activity in Alzheimer's disease. \emph{Annals of Neurology (2016)}
  \end{flushright}

  \tcblower

  \noindent \textbf{STEP IDENTIFICATION}

  \vspace*{10pt}
  {\small\noindent
  \begin{tblr}{colspec={X[1,c]X[5,l]},cell{even}{1,2} = {azure9}}
    \toprule
    \textbf{Step} & \textbf{Sample} \\ 
    \midrule
    
    1S1 & Seizures are more frequent in patients with Alzheimer's disease (AD) and can hasten cognitive decline. \\
    1S2 & However, the incidence of subclinical epileptiform activity in AD and its consequences are unknown. \\
    1S3b & Motivated by results from animal studies, we hypothesized higherthan expected rates of subclinical epileptiform activity in AD with deleterious effects on cognition. \\
      
    \bottomrule
  \end{tblr}
  }
  
\end{sample}

1S3a是这一语步中相对常规的语阶，而其他语阶是可选的。这些语阶遵循上述顺序，以实现作者的交流目的，即澄清什么是已知的，什么是未知的，什么是要研究的。

\begin{sample}[label={myautocounter}]{\heiti}

  \textbf{BACKGROUND} 
  
  CT-P6 is a proposed biosimilar to reference trastuzumab. In this study, we aimed to establish equivalence of CT-P6 to reference trastuzumab in neoadjuvant treatment of HER2-positive early-stage breast cancer.

  
  \begin{flushright}
    ---CT-P6 compared with reference trastuzumab for HER2-positive breast cancer: a randomised, double-blind, active-controlled, phase 3 equivalence trial. \emph{Lancet Oncology (2017)}
  \end{flushright}

  \tcblower

  \noindent \textbf{STEP IDENTIFICATION}

  \vspace*{10pt}
  {\small\noindent
  \begin{tblr}{colspec={X[1,c]X[5,l]},cell{even}{1,2} = {azure9}}
    \toprule
    \textbf{Step} & \textbf{Sample} \\ 
    \midrule
    
    1S1 & BACKGROUND: CT-P6 is a proposed biosimilar to reference trastuzumab. \\
    1S3a & In this study, we aimed to establish equivalence of CT-P6 to reference trastuzumab in neoadjuvant treatment of HER2-positive early-stage breast cancer. \\
      
    \bottomrule
  \end{tblr}
  }

\end{sample}

\begin{sample}[label={myautocounter}]{\heiti}

  At a population level, dietary consumption of fish rich in docosahexaenoic acid (DHA) is associated with prevention of cognitive decline but this association is not clear in carriers of the apolipoprotein E epsilon 4 allele (E4). Plasma and liver DHA concentrations show significant alterations in EA carriers, in part corrected by DHA supplementation. However, whether DHA sufficiency in E4 carriers has consequences on cognition is unknown.

  
  \begin{flushright}
    ---Docosahexaenoic acid prevents cognitive deficits in human apolipoprotein E epsilon 4-targeted replacement mice. \emph{Neurobiology of Aging (2017)}
  \end{flushright}

  \tcblower

  \noindent \textbf{STEP IDENTIFICATION}

  \vspace*{10pt}
  {\small\noindent
  \begin{tblr}{colspec={X[1,c]X[5,l]},cell{even}{1,2} = {azure9}}
    \toprule
    \textbf{Step} & \textbf{Sample} \\ 
    \midrule
    
    1S1 & At a population level, dietary consumption of fish rich in docosahexaenoic acid (DHA) is associated with prevention of cognitive decline but this association is not clear in carriers of the apolipoprotein E epsilon 4 allele (E4). Plasma and liver DHA concentrations show significant alterations in EA carriers, in part corrected by DHA supplementation. \\
    1S2 & However, whether DHA sufficiency in E4 carriers has consequences on cognition is unknown. \\
      
    \bottomrule
  \end{tblr}
  }

\end{sample}

\begin{sample}[label={myautocounter}]{\heiti}

  The aim of this cohort study is to compare the symptom burden of patients who have an planned admission. unplanned admission to an acute palliative care unit (APCU) with patients who have a regular planned admission.

  
  \begin{flushright}
    ---Characteristics of patients with an unplanned admission to an acute palliative care unit. \emph{Internal and Emergency Medicine (2017)}
  \end{flushright}

  \tcblower

  \noindent \textbf{STEP IDENTIFICATION}

  \vspace*{10pt}
  {\small\noindent
  \begin{tblr}{colspec={X[1,c]X[5,l]},cell{even}{1,2} = {azure9}}
    \toprule
    \textbf{Step} & \textbf{Sample} \\ 
    \midrule
    
    1S3a & The aim of this cohort study is to compare the symptom burden of patients who have an planned admission. unplanned admission to an acute palliative care unit (APCU) with patients who have a regular planned admission. \\
      
    \bottomrule
  \end{tblr}
  }

\end{sample}

然而，需要特别提醒的是，在少数情况下，这些语阶也可以用不同的顺序呈现。

\begin{sample}[label={myautocounter}]{\heiti}
  
  \textbf{OBJECTIVE}

  To compare the rate of positive resection margins between radioactive seed localization (RSL) and wire-guided localization (WGL) after breast conserving surgery (BCS).
 
  \textbf{BACKGROUND} 
  
  WGL is the current standard for localization of nonpalpable breast lesions in BCS, but there are several difficulties related to the method.

  
  \begin{flushright}
    ---Radioactive Seed Localization or Wire-guided Localization of Nonpalpable Invasive and InSitu Breast Cancer: A Randomized, Multicenter, Open-label Trial. \emph{Annals of Surgery (2017)}
  \end{flushright}

  \tcblower

  \noindent \textbf{STEP IDENTIFICATION}

  \vspace*{10pt}
  {\small\noindent
  \begin{tblr}{colspec={X[1,c]X[5,l]},cell{even}{1,2} = {azure9}}
    \toprule
    \textbf{Step} & \textbf{Sample} \\ 
    \midrule
    
    1S3a & OBJECTIVE: To compare the rate of positive resection margins between radioactive seed localization (RSL) and wire-guided localization (WGL) after breast conserving surgery (BCS). \\
    1S1 & BACKGROUND: WGL is the current standard for localization of nonpalpable breast lesions in BCS, \\
    1S2 & but there are several difficulties related to the method. \\
      
    \bottomrule
  \end{tblr}
  }

\end{sample}

\begin{sample}[label={myautocounter}]{\heiti}

  \textbf{OBJECTIVE} 
  
  The aim of this study was to investigate the efficacy of intraperitoneal local anesthetic (IPLA) on pain after acute laparoscopic appendectomy in children.

  \textbf{SUMMARY OF BACKGROUND} 
  
  IPLA reduces pain in adult elective surgery. It has not been well studied in acute peritoneal inflammatory conditions. We hypothesized that IPLA would improve recovery in pediatric acute laparoscopic appendectomy.

  
  \begin{flushright}
    ---Intraperitoneal Local Anesthetic for Laparoscopic Appendectomy in Children: A Randomized Controlled Trial. \emph{Annals of Surgery (2017)}
  \end{flushright}

  \tcblower

  \noindent \textbf{STEP IDENTIFICATION}

  \vspace*{10pt}
  {\small\noindent
  \begin{tblr}{colspec={X[1,c]X[5,l]},cell{even}{1,2} = {azure9}}
    \toprule
    \textbf{Step} & \textbf{Sample} \\ 
    \midrule
    
    1S3a & \textbf{OBJECTIVE}: The aim of this study was to investigate the efficacy of intraperitoneal local anesthetic (IPLA) on pain after acute laparoscopic appendectomy in children. \\
    1S1 & \textbf{SUMMARY OF BACKGROUND}: IPLA reduces pain in adult elective surgery.  \\
    1S2 & It has not been well studied in acute peritoneal inflammatory conditions. \\
    1S3b & We hypothesized that IPLA would improve recovery in pediatric acute laparoscopic appendectomy. \\
      
    \bottomrule
  \end{tblr}
  }

\end{sample}

\section{每个语阶的语言特征}

  \subsection{语阶1(1S1) 介绍当前的知识或由以前的研究确立的相关信息。}

    \subsubsection{语阶分析}

    通过这一语阶，作者自然而然地介绍了所进行的研究。在这语阶中，介绍或解释了当前的知识或由以前的研究确立的相关信息，其中可能包括相关的机制或定义、某一现象的可能性、某项研究的意义等。

    有两种方式可以引出这一语阶。作者可以在同一范围内介绍知识或信息（如E.g.~\ref{eg:1s1-1}）。另外，作者还可以介绍以前相关研究的结论（如E.g.~\ref{eg:1s1-2}）。在后一种情况下，作者总是采用诸如“preclinical studies”这样的词或短语来表明信息的来源（如E.g.~\ref{eg:1s1-2}）。
    
    \begin{eg}[label={eg:1s1-1}]{}
      Topical immunomodulators(Tl)-including corticosteroids, calcineurin inhibitors, and vitamin D analogues-are commonly prescribed in multiple specialties,$\dots$
    \end{eg}

    \begin{eg}[label={eg:1s1-2}]{}
      \uline{Preclinical studies} have found radiotherapy enhances antitumour immune responses.
    \end{eg}

    \subsubsection{语言实现}

    一般现在时在这一语阶中经常被用来描述当前对该主题的理解。

    \begin{eg}[label={myautocounter}]{}
      In chronic hemodialysis, physical functioning (PF) \uline{is known} to be poor.
    \end{eg}

    当“previous studies”或其他具有类似含义的短语被用作句子的主语，现在完成时就成了首选。不过在这种情况下，一般现在时仍被用来指称当前的知识或由以前的研究建立的相关信息。

    \begin{eg}[label={myautocounter}]{}
      Previous studies have shown that more active older adults have better cognition and brain health based on a variety of structural neuroimaging measures.
    \end{eg}

    \begin{task}[label={myautocounter}]{\heiti 语料库任务}
      在M1中，哪种动词常用现在完成时？
    \end{task}

    在这一语阶，主动语态和被动语态都可以使用。

    \begin{eg}[label={myautocounter}]{}
      Physicians \uline{are} often \uline{asked} to prognosticate soon after a patient presents with stroke.
    \end{eg}

    \begin{eg}[label={myautocounter}]{}
      Youth baseball frequently \uline{results} in repetitive strain injuries.
    \end{eg}

    在这一语阶中通常会使用情态助词，如“may”、“can”，以及其他表示可能性的词，如“possible”。

    \begin{eg}[label={myautocounter}]{}
      Urban design \uline{may} affect children's habitual physical activity by influencing active commuting and neighborhood play.
    \end{eg}

    \begin{eg}[label={myautocounter}]{}
      Hirsutism in females \uline{can} be a source of considerable psychological distress and a threat to female identity.
    \end{eg}

    \begin{eg}[label={myautocounter}]{}
      It is \uline{possible} that psychological counseling could enhance stuttering treatment outcomes.
    \end{eg}

    \subsubsection{词组}

    \begin{enumerate}
      \item has been associated with
      \begin{eg}[label={myautocounter}]{}
        Ethnicity \uline{has been associated with} clinical and experimental pain responses.
      \end{eg}

      \item has been shown to
      \begin{eg}[label={myautocounter}]{}
        Estrogen administration following menopause \uline{has been shown to} support hippocampally mediated cognitive processes.
      \end{eg}
    \end{enumerate}
  
  \subsection{语阶2(1S2) 确立利基/问题}

    \subsubsection{语阶分析}

    在这一语阶中，通常对语阶1提供的信息给予负面评价，包括以前研究的局限性或缺陷，这就留下了一个有待填补的空白，一个有待解决的问题或一个有待检验的想法。换句话说，这一步暗示了作者研究的价值和意义。

    \begin{eg}{}
      However, to date, underlying neuronal mechanisms of these WM load-dependent activation changes in aging remain \uline{poorly understood}.
    \end{eg}

    \subsubsection{语言实现}

    一般现在时通常出现在这个语阶中。但当“study”这个词作为主动句的主语，或在被动句中作为省略代词时，更多的是采用现在完成时。

    \begin{eg}{}
      Despite the availability of objective tests, gastroesophageal reflux disease (GERD) diagnosis and management in infants \uline{remains} controversial and highly variable.
    \end{eg}

    \begin{eg}{}
      However, no study \uline{has} directly \uline{compared} these outcomes between sports.
    \end{eg}

    诸如“however”和“but”这样的词可以用来引导读者从语阶1过渡到语阶2，将注意力放到要指出的问题或差距上。

    \begin{eg}{}
      Delirium is associated with adverse postoperative outcomes, \uline{but} controversy exists regarding whether delirium is an independent predictor of mortality.
    \end{eg}

    \begin{eg}{}
      Excellent outcomes have been reported for anterior cruciate ligament (ACL) reconstruction (ACLR) in professional athletes in a number of different sports. \uline{However}, no study has directly compared these outcomes between sports.
    \end{eg}

    除了上面提到的两个词，其他表示对比或否定意义的词，如“not”或“unknown”，也可以看作是1S2的信号词。

    \begin{eg}{}
      Validated models to predict risk for complications are \uline{not} available, and the effect of treatment on risk is \uline{unknown}.
    \end{eg}

    \subsubsection{词组}

    \begin{enumerate}
      \item little is known about the
      \begin{eg}{}
        However, \uline{little is known about the} effect of pharmacological PHD inhibition on tumor expansion, and on liver regeneration after surgical resection.
      \end{eg}
    \end{enumerate}

  \subsection{语阶3(1S3a) 指明主要目的}

    \subsubsection{语阶分析}

    语阶1和语阶2显示了有关该研究的已知情况和尚待完成的工作。然后，语阶3a说明了所提出的研究的目的或期望。它概述了进行该研究的原因。

    \subsubsection{语言实现}

    在大多数情况下，动词的不定式(to-infinitive)被用来表达这一语阶的目的。它既可以放在句子的开头，也可以放在中间，这一特点可以从下面的例子看出。

    \begin{eg}{}
      \uline{To achieve} the efficient usage of contrast material (CM) in high-pitch CT aortography, an appropriate duration of the CM injection is crucial.
    \end{eg}

    \begin{eg}{}
      The purpose of this study was \uline{to evaluate} how travel burden and hospital volume influence treatment and outcomes of patients with locally advanced esophageal cancer.
    \end{eg}

    “The aim/purpose/objective of this study”或“we”经常被用作句子的主语。一般过去时在这语阶中经常出现，特别是当“was to”、“sought to”、“tested”等作为谓语时。

    \begin{eg}{}
      \uuline{The objective of this study} \uline{was to} measure the success of intubation of a simulated patient in an upright versus supine position by novice intubators after brief training.
    \end{eg}
    \begin{eg}{}
      \uuline{We} \uline{sought to} determine if time from emergency department (ED) physician evaluation until operative intervention is independently associated with appendiceal perforation (AP) in children.
    \end{eg}
    \begin{eg}{}
      \uuline{We} \uline{tested} whether any association with APOE e4 status on cognitive ability was larger in older ages or in those with cardiometabolic diseases.
    \end{eg}
    
    
    需要提及的是，当 "report "作为谓语时，将使用一般现在时。

    \begin{eg}{}
      We \uline{report} the effect of intravenous alteplase on long-term survival after ischaemic stroke of participants in the Third International Stroke Trial (IST-3).
    \end{eg}

    \subsubsection{词组}

    \begin{enumerate}
      \item the aim of this study was to evaluate the
      \begin{eg}{}
        \uline{The aim of this study was to evaluate the} application of Mindray BC-6800 body fluid (BF) mode in cytometric analysis of CSF compared to light microscopy (LM).
      \end{eg}

      \item the aim of this study was to investigate
      \begin{eg}{}
        \uline{The aim of this study was to investigate} the S100B utility for the determination of concussion in a professional 15-players-a-side rugby team.
      \end{eg}

      \item the aim of this study was to assess
      \begin{eg}{}
        \uline{The aim of this study was to assess} the effect of HDF on serum sST2 and NT-proBNP concentrations in End-stage Renal Disease (ESRD) patients.
      \end{eg}

      \item the purpose/objective of this study was to
      \begin{eg}{}
        \uline{The purpose of this study was to} present our investigation of the utility of a novel flexible robotic system for transoral supraglottic laryngectomy and total laryngectomy.
      \end{eg}

      \item of this study was to determine the
      \begin{eg}{}
        The overall aim \uline{of this study was to determine the} effect of introducing a smartphone pain application (app), for both Android and iPhone devices that enables chronic pain patients to assess, monitor, and communicate their status to their providers.
      \end{eg}

      \item of this study was to investigate/evaluate/compare/examine
      \begin{eg}{}
        The aim \uline{of this study was to evaluate} the diagnostic performance of susceptibility-weighted magnetic resonance imaging (SW-MRI) for the differentiation of osteophytes and disc herniations of the spine compared with that of conventional spine MR sequences and radiography.
      \end{eg}

      \item we aimed to assess the
      \begin{eg}{}
        \uline{We aimed to assess the} efficacy and safety of the MEK inhibitor binimetinib versus that of dacarbazine in patients with advanced NRAS-mutant melanoma.
      \end{eg}

    \end{enumerate}

  \subsection{语阶3(1S3b) 提出假设}
    
    \subsubsection{语阶分析}

    虽然大多数作者在M1中提出了目标或目的，以介绍要解决的问题，但也有少数作者可能选择提出一个假设，由研究来检验或证明。

    \begin{eg}{}
      We hypothesized that the use of 2-chloroprocaine would be associated with a faster recovery from sensorimotor block.
    \end{eg}

    \subsubsection{语言实现}

    在这个语阶中，动词“hypothesize”作为谓语，主语可以是“we”或“the authors”，前者比后者更常见。一般过去时比一般现在时更常被使用。而情态动词“would”在句子中经常出现，以表示概率或可能性。

    \begin{eg}{}
      In the current study, \uline{we hypothesized} that this method \uline{would} be equal to or better than the transversus abdominis plane block with regard to pain relief and its duration of action after cesarean delivery.
    \end{eg}

    \begin{eg}{}
      \uline{The authors hypothesized} that the outcomes of meniscal repair associated with concomitant multiligament reconstruction \uline{would} significantly improve from preoperatively to postoperatively at a minimum of 2 years after the index surgery.
    \end{eg}


    \subsubsection{词组}

    \begin{enumerate}
      \item tested the hypothesis that
      \begin{eg}
        We \uline{tested the hypothesis that} extrafascial placement of the catheter tip reduces the rate of hemidiaphragmatic paresis compared with intrafascial tip placement for CISB while providing effective analgesia.
      \end{eg}
    \end{enumerate}

\section{材料阅读}

这一小节提供了3段M1材料，并对其中语阶的语言特点进行了详细分析。

\begin{sample}[label={myautocounter}]{\heiti}
  
  \textbf{BACKGROUND} 
  
  Head impacts and resulting head accelerations cause concussive injuries. There is no standard for reporting head impact data in sports to enable comparison between studies.

  \textbf{OBJECTIVE }
  
  The aim was to outline methods for reporting head impact acceleration data in sport and the effect of the acceleration thresholds on the number of impacts reported.


  \begin{flushright}
    ---The Influence of Head Impact Threshold for Reporting Data in Contact and Collision Sports: Systematic Review and Original Data Analysis. \emph{Sports Medicine (2015)}
  \end{flushright}

  \tcblower

  \noindent \textbf{STEP IDENTIFICATION}

  \vspace*{10pt}
  {\small\noindent
  \begin{tblr}{colspec={X[1,c]X[5,l]},cell{even}{1,2} = {azure9}}
    \toprule
    \textbf{Step} & \textbf{Sample} \\ 
    \midrule
    
    1S1 & \textbf{BACKGROUND}: Head impacts and resulting head accelerations cause concussive injuries. \\
    1S2 & There is no standard for reporting head impact data in sports to enable comparison between studies. \\
    1S3a & \textbf{OBJECTIVE }: The aim was to outline methods for reporting head impact acceleration data in sport and the effect of the acceleration thresholds on the number of impacts reported. \\
      
    \bottomrule
  \end{tblr}
  }

  \noindent \textbf{ANALYSIS}

  This is a typical example of M1 involving three steps. The first sentence informs the readers of the background information related to the research. The second sentence, where the word ``no'' can be seen as a signal for 1S2, identifies the problem---the lack of a standard. The last sentence, with the subject ``aim'' and the to-infinitive, shows the research objective.

\end{sample}

\begin{sample}[label={myautocounter}]{\heiti}
  \textbf{BACKGROUND} 
  
  Ergometrine is a uterotonic agent that is recommended in the prevention and management of postpartum hemorrhage. Despite its long-standing use, the mechanism by which it acts in humans has never been elucidated fully. The objective of this study was to investigate the role of adrenoreceptors in ergometrine's mechanism of action in human myometrium. The study examined the hypothesis that a-adrenoreceptor antagonism would result in the reversal of the uterotonic effects of ergometrine.

  \begin{flushright}
    ---A Role for Adrenergic Receptors in the Uterotonic Effects of Ergometrine in Isolated Human Term Nonlaboring Myometrium. \emph{Anesthesia and-Analgesia (2017)}
  \end{flushright}

  \tcblower

  \noindent \textbf{STEP IDENTIFICATION}

  \vspace*{10pt}
  {\small\noindent
  \begin{tblr}{colspec={X[1,c]X[5,l]},cell{even}{1,2} = {azure9}}
    \toprule
    \textbf{Step} & \textbf{Sample} \\ 
    \midrule
    
    1S1  & BACKGROUND: Ergometrine is a uterotonic agent that is recommended in the prevention and management of postpartum hemorrhage.  \\
    1S2  & Despite its long-standing use, the mechanism by which it acts in humans has never been elucidated fully. \\
    1S3a  & The objective of this study was to investigate the role of adrenoreceptors in ergometrine's mechanism of action in human myometrium. \\
    1S3b  & The study examined the hypothesis that a-adrenoreceptor antagonism would result in the reversal of the uterotonic effects of ergometrine. \\
      
    \bottomrule
  \end{tblr}
  }

  \noindent \textbf{ANALYSIS} 

  This paragraph involves four sentences. The first one provides the readers with the necessary knowledge of Ergometrine, a drug used to promote the contractions of the muscle ofthe womb (uterus). The second sentence, with the word "never"serving as the signal for 1S2, points out the gap that needs to be filled -the mechanism has not "been elucidated fully". Here, "elucidate"means to make clear. Present perfect tense is used, because the omitted agent in this passive sentence is "the previous study". The third sentence introduces the objective and the last one raises a hypothesis. 
  
  What needs to be noticed is that 1S3a and 1S3b are not very often seen with each other. This sample is one of the few exceptions.

\end{sample}

\begin{sample}[label={myautocounter}]{\heiti}
  
  \textbf{OBJECTIVE} 
  
  To investigate the prognostic significance of p16 in patients with hypopharyngeal squamous cell carcinoma (HPSCC) and to evaluate the relationship between p16 and human papillomavirus (HPV). Unlike in oropharyngeal SCC (OPSCC), the prognostic significance of p16 in HPSCC and its association with HPV is unclear.


  \begin{flushright}
    ---p16 not a prognostic marker for hypopharyngeal squamous cell carcinoma.
    
    \emph{Archives of Otorhinolaryngology-Head \& Neck Surgery (2012)}
  \end{flushright}

  \tcblower

  \noindent \textbf{STEP IDENTIFICATION}

  \vspace*{10pt}
  {\small\noindent
  \begin{tblr}{colspec={X[1,c]X[5,l]},cell{even}{1,2} = {azure9}}
    \toprule
    \textbf{Step} & \textbf{Sample} \\ 
    \midrule
    
    1S3a & \textbf{OBJECTIVE}: To investigate the prognostic significance of p16 in patients with hypopharyngeal squamous cell carcinoma (HPSCC) and to evaluate the relationship between p16 and human papillomavirus (HPV). \\
    1S2 & Unlike in oropharyngeal SCC (OPSCC), the prognostic significance of p16 in HPSCC and its association with HPV is unclear. \\
      
    \bottomrule
  \end{tblr}
  }

  \noindent \textbf{ANALYSIS} 
  
  This sample presents the steps with orders different from Sample 1 and 2. It starts with to-infinitives showing the objectives, and proceeds with the research gap, "unclear"being a signal.

\end{sample}

\section{词汇表}

{\small
\begin{longtblr}[
    caption = {第三章词汇表},
    label = {tab:Glossary of Chapter 3},
]{
    width = \textwidth,
    colspec = {X[1,l,h]  X[1,l,h]  X[3,l,h]},
    rowhead = 1, rowfoot = 0, % 每个分页里表头表尾的数量
    % row{odd} = {blue8}, 
    row{even} = {azure9},
}
    
\toprule
\textbf{WORDS} & \textbf{MEANING} & \textbf{MEANING or EXAMPLE}\\
\midrule

\textbf{chunk}/\textipa{tS2Nk}/ & \emph{n.} 组块；区块；数据块 & a considerable amount \\
{\textbf{mechanism} \\ /\textipa{"mek@""nIz@m}/} & \emph{n.} 机制；机械装置；方法 & a special way of getting something done within a particular system \\
\textbf{modal auxiliary verb} & \emph{n.} 情态动词 & A modal auxiliary verb, often simply called a modal verb or even just a modal, is used to change the meaning of other verbs (commonly known as main verbs) by expressing modality-that is, asserting (or denying) possibility, likelihood, ability, permission, obligation, or future intention. \\
{\textbf{pertinent} \\ /\textipa{"p3:rtIn@nt}/} & \emph{adj.} 有关的；恰当的；相宜的 & relevant to a particular subject \\
{\textbf{predicate} \\ /\textipa{"predIk@t}/} & \emph{n.} 谓语 & a part of a sentence containing a verb that makes a statement about the subject of the verb, such as ``went home''in ``John went home. d'' \\
{\textbf{territory} \\ /\textipa{"ter@t0:ri}/} & \emph{n.} 领域；地域； & refer to an area of knowledge or experience \\

\bottomrule

\end{longtblr}
}

\begin{problemset}
  \item \textbf{Find out how many steps are involved in the following examples.}
  1S1 Presenting current knowledge or relevant information established by previous studies

  1S2 Establishing a niche/problem

  1S3a Indicating main purposes

  1S3b Raising hypotheses

  \begin{enumerate}
    \item At a population level, dietary consumption of fish rich in docosahexaenoic acid (DHA) is associated with prevention of cognitive decline but this association is not clear in carriers of the apolipoprotein E epsilon 4 allele (E4). Plasma and liver DHA concentrations show significant alterations in E4 carriers, in part corrected by DHA supplementation. However, whether DHA sufficiency in E4 carriers has consequences on cognition is unknown.
    \item In chronic hemodialysis, physical functioning (PF) is known to be poor. We set out to assess to what extent chronic dialysis patients are able to maintain a good physical condition over time and what the influence of age is on the trajectory of PF.
    \item Both patient characteristics and intraoperative factors have been associated with a higher risk of stroke after cardiac surgery. We hypothesized that poor systemic oxygenation in the perioperative period is associated with increased risk of stroke following cardiopulmonary bypass.
    \item Therapies to extend healthspan are poised to move from laboratory animal models to human clinical trials. Translation from mouse to human will entail challenges, among them the multifactorial heterogeneity of human aging. To inform clinical trials about this heterogeneity, we report how humans' pace of biological aging relates to personal-history characteristics.
  \end{enumerate}

  \item \textbf{Rearrange the sentences in correct order.}

  \begin{enumerate}
    \item \uline{\hspace*{3em}}
    \begin{itemize}
      \item[A] Regular use of sunbed exposure has been reported to increase 25-hydroxyvitamin-D3 [\ce{25(OH)D}] serum levels.
      \item[B] We investigated the impact of standard sunbed use compliant with the European Union standard on \ce{25(OH)D} serum modulation and well-being.
      \item[C] However, the influence of sunbeds compliant with the recent European Union standard EN-60335-2-27 on \ce{25(OH)D} serum levels is unknown.
    \end{itemize}

    \item \uline{\hspace*{3em}}
    \begin{itemize}
      \item[A] The Indian Health Service provides health care to eligible American Indians and Alaskan Natives.
      \item[B] We seek to determine the characteristics and capabilities of Indian Health Service emergency departments (EDs).
      \item[C] No published data exist on emergency services offered by this unique health care system.
    \end{itemize}

    \item \uline{\hspace*{3em}}
    \begin{itemize}
      \item[A] This trial compared immediate posttreatment effects of family-focused treatment for childhood depression (FFT-CD) with those of individual supportive psychotherapy (IP) for children 7 to 14 years old with depressive disorders 
      \item[B] Integrating family in treatment could have particularly salutary effects during this developmental period.
      \item[C] Despite the morbidity and negative outcomes associated with early-onset depression, few studies have examined the efficacy of psychosocial treatment for depressive disorders during childhood.
    \end{itemize}

    \item \uline{\hspace*{3em}}
    \begin{itemize}
      \item[A] We tested the hypothesis that aquatic treadmill exercise would augment CBF and lower HR compared with land-based treadmill exercise.
      \item[B] However, their effect on cerebral blood flow (CBF) responses has not been examined.
      \item[C] Aquatic treadmills are used as a rehabilitation method for conditions such as spinal cord injury, osteoarthritis, and stroke, and can facilitate an earlier return to exercise training for athletes.
    \end{itemize}
  \end{enumerate}

  \item \textbf{Fill in the blanks with appropriate words or the correct form of the verbs given}

  \begin{enumerate}
    \item This single-centre, randomized trial \uline{\hspace*{3em}} (test) the hypothesis that the administration of FFP after CPB (late FFP group) is superior to FFP priming (early FFP group)
    in terms of postoperative bleeding and overall red blood cell (RBC) transfusion.
    \item We assessed whether addition of the antiepileptic drug levetiracetam to the benzodiazepine clonazepam \uline{\hspace*{3em}}(improve) prehospital treatment of GCSE.
    \item \uline{\hspace*{3em}} (investigate) the prevalence, injury rate, severity, nature, and economic burden of RRIs in Dutch trailrunners.
    \item We \uline{\hspace*{3em}} (aim) to assess the safety and activity of an anti-PD-1 antibody pembrolizumab in patients with locally advanced or metastatic urothelial cancer.
    \item We \uline{\hspace*{3em}} (report) 2-year overall survival data from a randomised controlled trial assessing this treatment in previously untreated advanced melanoma.
    \item Previous studies \uline{\hspace*{3em}} (propose) that faster speed improves movement quality.
    \item We \uline{\hspace*{3em}} (seek) to determine the compliance rate within the pigmented lesions clinic at our academic institution and identify demographic variables that may influence adherence.
    \item Anxiety disorders constitute a major disease and social burden worldwide; \uline{\hspace*{3em}}, many questions concerning the underlying molecular mechanisms still remain open.
    \item Oncologic outcomes for induction chemotherapy and its role in patients with advanced olfactory neuroblastoma (ONB) \uline{\hspace*{3em}} unclear.
    \item Histaminergic neurons are crucial to maintain wakefulness, but their role in cataplexy is \uline{\hspace*{3em}}.
    \item Background and Objectives: Although many studies \uline{\hspace*{3em}} (find) no difference between thoracic epidural block and unilateral thoracic paravertebral block after thoracotomy, no previous studies \uline{\hspace*{3em}} (compare) epidural block with bilateral thoracic paravertebral block (bTPVB) in patients undergoing open liver resection. We \uline{\hspace*{3em}} (aim) investigate whether there was a significant analgesic advantage of thoracic epidural over bTPVB after liver resection.
    \item Background: Chronic tendinopathy \uline{\hspace*{3em}} (be) a commonly occurring clinical problem that affects both athletes and inactive middle-aged patients. Although some studies \uline{\hspace*{3em}} (show) that different platelet-rich plasma (PRP) preparations could exert various therapeutic effects in vitro, the role of leukocytes in PRP \uline{\hspace*{3em}} (be) defined under tendinopathy conditions in vivo.
    
    \hspace*{2em}Purpose: This study \uline{\hspace*{3em}} (compare) the effects of the intratendon delivery of leukocyte-poor PRP (Lp-PRP) versus leukocyte-rich PRP (Lr-PRP) in a rabbit chronic tendinopathy model in vivo.
    \item Background: Eye trackers \uline{\hspace*{3em}} (be) widely used among people with amyotrophic lateral sclerosis, and their benefits to quality of life \uline{\hspace*{3em}} (show). On the contrary, Brain-computer interfaces (BCIs) \uline{\hspace*{3em}} (be) still quite a novel technology, which also serves as an access technology for people with severe motor impairment.
    
    \hspace*{2em}Objective: \uline{\hspace*{3em}} (compare) a visual P300-based BCI and an eye tracker in terms of information transfer rate (ITR), usability, and cognitive workload in users with motor impairments. 
  \end{enumerate}

  \item \textbf{Discussion: Read the following examples and discuss with your partners about the reasons why modal auxiliary verbs and words indicating likelihood are used?}

  \textbf{Example 1}

  \hspace*{2em}In the treatment of anxiety disorders, attention bias modification therapy (ABMT) and cognitive-behavioral therapy (CBT) may have complementary effects by targeting different aspects of perturbed threat responses and behaviors. ABMT may target rapid, implicit threat reactions, whereas CBT may target slowly deployed threat responses. The authors used amygdala-based connectivity during a threat-attention task and a randomized controlled trialdesign to evaluate potential complementary features of these treatments in pediatric anxiety disorders.

  \begin{flushright}
    ---Complementary Features of Attention Bias Modification Therapy and Cognitive-Behavioral Therapy in Pediatric Anxiety Disorders. \emph{The American Journal of Psychiatry (2017)}
  \end{flushright}

  \textbf{Example 2}

  \hspace*{2em}Ticagrelor is an effective antiplatelet therapy for patients with coronary atherosclerotic disease and might be more effective than aspirin in preventing recurrent stroke and cardiovascular events in patients with acute cerebral ischaemia of atherosclerotic origin. Our aim was to test for a treatment-by-ipsilateral atherosclerotic stenosis interaction in a subgroup analysis of patients in the Acute Stroke or Transient Ischaemic Attack Treated with Aspirin or Ticagrelor and Patient Outcomes (SOCRATES) trial.

  \begin{flushright}
    ---Efficacy and safety of ticagrelor versus aspirin in acute stroke or transient ischaemic attack of atherosclerotic origin: a subgroup analysis of SOCRATES, a randomised, double-blind, controlled trial. \emph{The Lancet. Neurology (2017)}
  \end{flushright}

  \textbf{Example 3}

  \hspace*{2em}Psychosocial disorders have been reported in adults who stutter, especially social anxiety disorder. Social anxiety has been linked to childhood victimization. It is possible that recalled childhood victimization could be linked to psychosocial problems reported in some adults who stutter.

  \begin{flushright}
    ---Long-term Consequences of Childhood Bullying in Adults Who Stutter: Social Anxiety, Fear of Negative Evaluation, Self-esteem, and Satisfaction with Life. \emph{Journal of Fluency Disorders (2016)}
  \end{flushright}


\end{problemset}


\chapter{语步二}\label{chapter4}

语步2（M2）描述了研究过程。在这一语步中，作者陈述了研究问题的方式。这可能包括研究对象和遵循的方法，以及其他事项。在$4\sim 5$段的摘要中，M2通常是一个段落；在5段以上的摘要中，M2被分成若干段落，每段前面有一个标题，这样更清晰，也更具体。

\section{M2中的语阶}

无论M2是哪种形式，它都可以进一步分为几个语阶。根据不同的研究类型，M2的语阶可能包括：

2S1 报告医学伦理审查情况

2S2 简要解释研究设计

2S3 描述受试者或数据及其选择标准

2S4 描述实验程序，如干预措施、检查等

2S5 描述主要结果及其观察指标

2S6 描述数据分析方法

2S7 报告注册信息 

在大多数情况下，这些语阶遵循上述顺序，以实现作者的交流目的（Sample \ref{sam:M2_1} \& \ref{sam:M2_2}）。有时，一些语阶可能与其他语阶嵌合在一起（Sample \ref{sam:M2_1} \& \ref{sam:M2_3}）。例如，2S2可以嵌入2S3、2S4和2S1；2S3可以嵌入2S2、2S4和2S5，以此类推。这些语阶的顺序可能与上述模式不同。在大多数摘要中，M2包括2S2、2S3、2S4和2S5。有些语阶在摘要中是可选的，特别是2S1、2S6和2S7（Sample \ref{sam:M2_1} \& \ref{sam:M2_2}）。

\begin{sample}[label={sam:M2_1}]{\heiti}

  \textbf{METHOD} 
  
  Randomized controlled trial for adolescents (12--18 years of age) with recent (past 3 months) suicide attempts or other self-harm. Youth were randomized either to SAFETY or to treatment as usual enhanced by parent education and support accessing community treatment (E-TAU). Outcomes were evaluated at baseline, 3 months, or end of treatment period, and were followed up through 6 to 12 months. The primary outcome was youth-reported incident suicide attempts through the 3-month follow-up.
  
  \begin{flushright}
    ---Cognitive-Behavioral Family Treatment for Suicide Attempt Prevention: A Randomized Controlled Trial. \emph{Journal of the American Academy of Child and Adolescent Psychiatry (2017)}
  \end{flushright}

  \tcblower

  \noindent \textbf{STEP IDENTIFICATION}

  \vspace*{10pt}
  {\small\noindent
  \begin{tblr}{colspec={X[1,c]X[5,l]},cell{even}{1,2} = {azure9}}
    \toprule
    \textbf{Step} & \textbf{Sample} \\ 
    \midrule
    
    2S3 with 2S2 embedded & Randomized controlled trial for adolescents (12-18 years of age) with recent (past 3 months) suicideattempts or other self-harm. \\
    2S4 & Youth were randomized either to SAFETY or to treatment as usual enhanced by parent education and support accessing community treatment (E-TAU). \\
    2S5 & Outcomes were evaluated at baseline, 3 months, or end of treatment period, and were followed up through 6 to 12 months. The primary outcome was youth-reported incident suicide attempts through the 3-month follow-up. \\

    \bottomrule
  \end{tblr}
  }

\end{sample}

\begin{sample}[label={sam:M2_2}]{\heiti}
  
  \textbf{DESIGN} 
  
  Cross-sectional observational study.

  \textbf{PARTICIPANTS} 
  
  One hundred nineteen children 2 to 16 years of age (mean age, 9.4 years; standard deviation [SD], 4.56 years) with glaucoma and their parents.

  \textbf{METHODS} 
  
  Completion of 3 validated instruments for children to assess (1) functional visual ability (FVA) with the Cardiff Visual Ability Questionnaire for Children (CVAQC), (2) VR QoL with the Impact of Vision Impairment for Children (IVI-C), and (3) HR QoL with the Pediatric Quality of Life Inventory (PedsQL) version 4.0.

  \textbf{MAIN OUTCOME MEASURES} 
  
  Cardiff Visual Ability Questionnaire for Children, IVI-C, and PedsQL scores.

  \begin{flushright}
    ---Quality of Life and Functional Vision in Children with Glaucoma. \emph{Ophthalmology (2017)}
  \end{flushright}

  \tcblower

  \noindent \textbf{STEP IDENTIFICATION}

  \vspace*{10pt}
  {\small\noindent
  \begin{tblr}{colspec={X[1,c]X[5,l]},cell{even}{1,2} = {azure9}}
    \toprule
    \textbf{Step} & \textbf{Sample} \\ 
    \midrule
    
    2S2 & DESIGN: Cross-sectional observational study. \\
    2S3 & PARTICIPANTS: One hundred nineteen children 2 to 16 years of age (mean age, 9.4 years; standard deviation [SD], 4.56 years) with glaucoma and their parents. \\
    2S4 & METHODS: Completion of 3 validated instruments for children to assess (1) functional visual ability (FVA) with the Cardiff Visual Ability Questionnaire for Children (CVAQC), (2) VR QoL with the Impact of Vision Impairment for Children (IVI-C), and (3) HR QoL with the Pediatric Quality of Life Inventory (PedsQL) version 4.0. \\
    2S5 & MAIN OUTCOME MEASURES: Cardiff Visual Ability Questionnaire for Children, IVI-C, and PedsQL scores. \\
     
    \bottomrule
  \end{tblr}
  }

\end{sample}

\begin{sample}[label={sam:M2_3}]{\heiti}
  
  \textbf{MATERIALS AND METHODS} 
  
  The institutional review board approved this retrospective study and waived the informed consent requirement. Seventy-four patients with surgically confirmed PNETs and 82 patients with PDACs who underwent gadobutrol-enhanced MR imaging were included. Two radiologists independently evaluated the morphologic characteristics and temporal enhancement patterns of each tumor. Quantitative analysis, including measurement of tumor size, maximal upstream parenchymal thickness (MUPT), contrast-to-noise ratio, and apparent diffusion coefficient values, was performed. Uni-and multivariate logistic regression analyses were performed to identify relevant features to differentiate between PNETs and PDACs.

  \begin{flushright}
    ---Nonhypervascular Pancreatic Neuroendocrine Tumors: Differential Diagnosis from Pancreatic Ductal Adenocarcinomas at MR Imaging-Retrospective Cross-sectional Study. \emph{Radiology (2017)}
  \end{flushright}

  \tcblower

  \noindent \textbf{STEP IDENTIFICATION}

  \vspace*{10pt}
  {\small\noindent
  \begin{tblr}{colspec={X[1,c]X[5,l]},cell{even}{1,2} = {azure9}}
    \toprule
    \textbf{Step} & \textbf{Sample} \\ 
    \midrule
    
     2S1 with 2S2 embedded & The institutional review board approved this retrospective study and waived the informed consent requirement. \\
     2S3 & Seventy-four patients with surgically confirmed PNETs and 82 patients with PDACs who underwent gadobutrol-enhanced MR imaging were included. \\
     2S4 & Two radiologists independently evaluated the morphologic characteristics and temporal enhancement patterns of each tumor. \\
     2S5 & Quantitative analysis, including measurement of tumor size, maximal upstream parenchymal thickness (MUPT), contrast-to-noise ratio, and apparent diffusion coefficient values, was performed. \\
     2S6 & Uni-and multivariate logistic regression analyses were performed to identify relevant features to differentiate between PNETs and PDACs. \\

    \bottomrule
  \end{tblr}
  }

\end{sample}

\section{每个语阶的语言特征}

M2中的大多数语阶都是过去式，这表明学习或研究已经进行了。主动语态和被动语态都可以在这一语步中使用，“we”作为主动语态的主语（如E.g.~\ref{eg:M2-1}），其他表示研究过程的名词作为被动语态的主语（如E.g.~\ref{eg:M2-2}）。

\begin{eg}[label={eg:M2-1}]{}
  \uuline{We} \uline{conducted} a multicentre, double-blind, randomised, placebo-controlled trial at three hospitals in Australia.
\end{eg}

\begin{eg}[label={eg:M2-2}]{}
  A correlation \uuline{analysis} \uline{was used} to examine the relationship between age and weight gain.
\end{eg}

  \subsection{语阶1(2S1) 报告医学伦理审查情况}

    \subsubsection{语阶分析}

    医学伦理审查已成为大多数以人类或动物为试验对象的医学研究中的必要条件。尽管它越来越重要，但在医学摘要中却很少提及。

    \subsubsection{语言实现}

    这一语阶通常采用过去式和被动语态，以“study”或“approval”等作为主语。下面是发现的极少数2S1的例子之一。请注意是过去式。
    
    \begin{eg}
      This \uuline{study} \uline{was} approved by our institutional review board, and informed consent \uline{was} waived due to its retrospective design.
    \end{eg}

    \subsubsection{词组}

    \begin{enumerate}
      \item study was approved by the
      \begin{eg}
        This prospective \uline{study was approved by the} institutional ethics committee, and written informed consent was obtained from all participants.
      \end{eg}
    \end{enumerate}

  \subsection{语阶 2(2S2) 简要解释研究设计}

    \subsubsection{语阶分析}

    这里的研究设计狭义上指的是研究类型。在大多数摘要中，这一可选语阶通常出现在M2的开头。

    \subsubsection{语言实现}

    2S2通常在简单的句子或名词短语中作为一个单独的段落（如E.g.~\ref{eg:2s2-1}）或与其他语阶嵌入，通常与2S3（如E.g.~\ref{eg:2s2-2}），偶尔与2S4或2S1。描述性词语经常被用来修饰“study”、“trial”等。“队列(cohort)”、“回顾性(retrospective)”、“前瞻性(prospective)”、“观察性(observational)”、“横断面(cross-sectional)”等词往往与“study”搭配在一起。而“随机(randomised)”、“对照(controlled)”、“双盲(double-blind)”等词则倾向于与“trial”搭配（例如，E.g.~\ref{eg:2s2-1} \& \ref{eg:2s2-2}）。

    \begin{eg}[label={eg:2s2-1}]{}
      DESIGN: \uline{Retrospective} \uuline{study}.
    \end{eg}

    \begin{eg}[label={eg:2s2-2}]{}
      In this \uline{randomised, regimen-controlled, double-blind, phase 2} trial \textbf{(2S2)}, we enrolled adult patients with multiple basal-cell carcinomas, including those with basal-cell nevus syndrome, who had one or more histopathologically confirmed and at least six clinically evident basal-cell carcinomas \textbf{(2S3)}.
    \end{eg}

    \subsubsection{词组}

    \begin{enumerate}
      \item a retrospective cohort study (of)
      \begin{eg}{}
        This was \uline{a retrospective cohort study of} singleton pregnancies delivered between 24 0/7 and 39 6/7 weeks, using 2005 through 2006 US national linked birth and death certificate data
      \end{eg}
      
      \begin{eg}{}
        \uline{A retrospective cohort study} was conducted in 2 distinct cohorts of female members of Kaiser Permanente Southern California, which is a large integrated healthcare delivery system.
      \end{eg}

      \item we conducted a retrospective
      \begin{eg}{}
        We conducted a retrospective collaborative study involving centers from 11 countries and 11 US institutions analyzing 102 ASNs by IMS.
      \end{eg}

      \item a cross-sectional study
      \begin{eg}{}
        This was \uline{a cross-sectional study of} 83 patients enrolled in the Morphea in Adults and Children cohort.
      \end{eg}

      \item a secondary analysis of
      \begin{eg}{}
        This was \uline{a secondary analysis of} a prospective cohort study conducted at eight Canadian hospitals.
      \end{eg}
    \end{enumerate}

    \begin{task}{\heiti 语料库任务}
      当需要用一句话描述研究设计时，可以使用何种动词时态？
    \end{task}

  \subsection{语阶3(2S3) 描述受试者或数据及其选择标准}

    \subsubsection{语阶分析}

    2S3是介绍受试者或参与者的语阶，包括研究的地点、受试者的特征、数量和选择标准等。对于大多数摘要来说，2S3是常规的。

    \subsubsection{语言实现}

    这个语阶可以是一个单独的段落，在名词短语（如E.g. \ref{eg:2s3-1}）或句子中。当它在句子中时，它可以与其他语阶一起嵌入，通常与2S2（如E.g. \ref{eg:2s3-2}）或2S4（如E.g. \ref{eg:2s3-3}）在句子中。如果参与者或选择标准被用作句子主语，被动语态优于主动语态（如E.g. \ref{eg:2s3-4}）。总是使用一般过去时。

    \begin{eg}[label={eg:2s3-1}]{}
      PARTICIPANTS: A total of 152 patients (152 eyes) with DME.
    \end{eg}

    \begin{eg}[label={eg:2s3-2}]{}
      This was a retrospective, multicenter cross-sectional analysis \textbf{(2S2)} of children (<19 years old) presenting to 16 pediatric EDs (2004--2008) \textbf{(2S3)}.
    \end{eg}

    \begin{eg}[label={eg:2s3-3}]{}
      From November 1994 through January 2002, we randomly assigned 731 men with localized prostate cancer \textbf{(2S3)} to radical prostatectomy or observation \textbf{(2S4)}.
    \end{eg}

    \begin{eg}[label={eg:2s3-4}]{}
      All Danish \uuline{patients} $\geqslant 18$ years on January 1,2012 with AD diagnosed by a hospital dermatologist \uline{were included}. \uuline{Patients} \uline{were age-and sex-matched} in a l : 4 ratio with general population controls.
    \end{eg}

    \subsubsection{词组}

    \begin{enumerate}
      \item were included in the/this
      \begin{eg}{}
        All patients who received at least one dose of nivolumab \uline{were included in the} primary and safety analyses.
      \end{eg}
        
      \begin{eg}{}
        Patients diagnosed from 1995 to 2014 \uline{were included in this} study.
      \end{eg}

      \item years or older with
      \begin{eg}{}
        Eligible patients were aged 18 \uline{years or older with} histologically or cytologically confirmed recurrent stage Illb or stage IV, chemotherapy-naive NSCLC.
      \end{eg}

      \item data were collected from
      \begin{eg}{}
        \uline{Data were collected from} the China Health and Nutrition Survey, a prospective open cohort and an ongoing nationwide health and nutrition survey, consisting of 3 199 apparently healthy Chinese girls aged 6 to 18 years at entry from 1991 to 2011.
      \end{eg}
    \end{enumerate}

  \subsection{语阶4(2S4) 描述实验程序，如干预措施、检查等。}

    \subsubsection{语阶分析}

    这里的实验程序可能包括干预、检查等。当研究以实验为基础或研究中有干预措施时，2S4通常包括在摘要中。它是M2的常规语阶。

    \subsubsection{语言实现}

    在这一语阶中，被动语态比主动语态更受欢迎，因为研究对象或方法经常被用作句子的主语。总是使用一般过去式。时间指标是本语阶的特点。

    \begin{eg}{}
      The asthma \uuline{intervention} \uline{was tailored} to the participant's allergen sensitivity and exposure, and it comprised 4 visits \uline{over the course of 1 year}.
    \end{eg}

    \subsubsection{词组}

    \begin{enumerate}
      \item (patients/participants) were randomly assigned (to) (receive)
      \begin{eg}{}
        From April 3,2014, to Dec 4,2015,667 \uline{patients were randomly assigned to receive} placebo (n=286), erenumab 70 mg (n=191), or erenumab 140 mg (n=190).
      \end{eg}

      \begin{eg}{}
        Dogs \uline{were randomly assigned to receive} intravenous OA to induce ALI (n=7 for each OA group) or saline as an OA control (n=6 for each control).
      \end{eg}

      \begin{eg}{}
        \uline{Patients were randomly assigned to} a choice of 100 or 200 ug ITM or no choice.
      \end{eg}

      \begin{eg}{}
        To induce PAH, Sprague-Dawley rats \uline{were randomly assigned to} treatment with monocrotaline or normal saline.
      \end{eg}

      \begin{eg}{}
        \uline{Participants were randomly assigned} (1: 1) centrally by an interactive voice response system, to receive either ipilimumab 10 mg/kg or placebo every 3 weeks for four doses, then every 3 months for up to 3 years.
      \end{eg}
      
      \item (patients) were randomized to (receive)
      \begin{eg}{}
        Eyes \uline{were randomized to receive} intravitreal injection of bevacizumab (1. 25 mg; n= 182) or aflibercept (2.0 mg; n=180) every 4 weeks through month 6.
      \end{eg}
      \begin{eg}{}
        \uline{Patients were randomized to} nasally inhaled isopropyl alcohol versus nasallyd inhaled normal saline solution.
      \end{eg}

      \item (were) masked to treatment (allocation/assignment)
      \begin{eg}{}
        Patients, study investigators, and study sponsor personnel \uline{were masked to treatment assignment}.
      \end{eg}
      \begin{eg}{}
        Patients were not \uline{masked to treatment allocation}.
      \end{eg}
      \begin{eg}{}
        The study was open label and no-one was \uline{masked to treatment assignment}.
      \end{eg}

      \item we randomly assigned patients
      \begin{eg}{}
        \uline{We randomly assigned patients} with advanced heart failure to receive either the new centrifugal continuous-flow pump or a commercially available axial continuous-flow pump.
      \end{eg}
    \end{enumerate}

  \subsection{语阶5(2S5) 描述主要结果及其观察指标}

    \subsubsection{语阶分析}

    世界卫生组织将观察指标定义为可归因于某项干预措施或一系列干预措施的个人、群体或人口的健康变化。选择观察指标是为了评估干预措施的影响。在临床试验中，观察指标可能包括死亡率(mortality)、治愈率(cure)、临床恶化(clinical worsening)、再入院率(readmission)等。有时研究者指定的其他指标也可能包括在这个语阶中。

    \subsubsection{语言实现}

    2S5可以用名词短语表述（如E.g.~\ref{eg:2s5-1}），也可以用通常是主动语态的句子，以“outcomes”、“outcome measure”等作为句子主语（如E.g.~\ref{eg:2s5-2}）。总是使用一般过去时。

    \begin{eg}[label={eg:2s5-1}]{}
      MAIN OUTCOME MEASURES: Cardiff Visual Ability Questionnaire for Children, IVI-C, and PedsQL scores.
    \end{eg}

    \begin{eg}[label={eg:2s5-2}]{}
      \uuline{The primary outcome} \uline{was} change from baseline to week 6 in the amount of urine leakage, measured by the 1-hour pad test. \uuline{Secondary outcomes} \uline{included} mean 72-hour urinary incontinence episodes measured by a 72-hour bladder diary (72-hour incontinence episodes).
    \end{eg}

    \subsubsection{词组}

    \begin{enumerate}
      \item the primary end point was (the)
      \begin{eg}{}
        \uline{The primary end point was the} duration of treatment for symptoms of neonatal opioid withdrawal.        
      \end{eg}

      \begin{eg}{}
        \uline{The primary end point was} a composite score of anaesthetists' non-technical skills (ANTS) assessed by two blinded evaluators.   
      \end{eg}
      
      \item the primary endpoint was (the)
      \begin{eg}{}
        \uline{The primary endpoint was the} rate of uncomplicated perineal wound healing defined as a Southampton wound score of less than 2 at 30 days postoperatively.   
      \end{eg}

      \begin{eg}{}
        \uline{The primary endpoint was} percentage reduction from baseline in the number of clinically evident basal-cell carcinomas at week 73.   
      \end{eg}

      \item (the primary) outcome (measure) was (the)
      \begin{eg}{}
        \uline{The primary outcome was} the mean difference in the angle of horizontal and vertical deviations after dilation in prism diopters.   
      \end{eg}
      \begin{eg}{}
        \uline{The primary outcome was} expired tidal volume.   
      \end{eg}
      \begin{eg}{}
        \uline{The primary outcome measure was} a three-level outcome-survival without neurodevelopmental impairment, survival with neurodevelopmental impairment, or death. 
      \end{eg}
      \begin{eg}{}
        The secondary \uline{outcome measure was the} impact of postoperative prophylaxis on donor tissue-associated infections.   
      \end{eg}

      \item was the proportion of patients
      \begin{eg}{}
        The primary endpoint, which has been reported previously, \uline{was the proportion of patients} with BRAFV600 wild-type melanoma achieving an investigator-assessed objective response.
      \end{eg}

      \item in all patients who
      \begin{eg}{}
        We assessed safety \uline{in all patients who} received at least one dose of study drug. 
      \end{eg}

      \item secondary end points included
      \begin{eg}{}
        \uline{Secondary end points included} overall survival, objective response rate, duration of response, effects on disease-related symptoms, safety, and tolerability.
      \end{eg}
    \end{enumerate}

  \subsection{语阶6(2S6) 描述数据分析方法}

    \subsubsection{语阶分析}

    当摘要中需要一些特定的数据提取和分析时，2S6才会被包括在内。它通常是可选的。

    \subsubsection{语言实现}

    这一语阶的句子通常以被动语态呈现，数据提取和分析方法作为句子的主语。总是使用一般过去式。

    \begin{eg}{}
      \uuline{Multivariable logistic regression analysis} \uline{was used} to estimate the risks of these complications among obese pregnancies compared with normal-weight pregnancies.    
    \end{eg}

    \subsubsection{词组}

    \begin{enumerate}
      \item logistic regression was used to
      \begin{eg}{}
        Multivariable logistic regression was used to investigate the association of the MDS score and AMD, taking account of potential confounders and the multicenter study design.   
      \end{eg}
      \item models were used to
      \begin{eg}{}
        Multivariable Cox and logistic regression models were used to examine associations between vaccination history and screening initiation and interval adherence. 
      \end{eg}
      \item regression analysis was used
      \begin{eg}{}
        Logistic regression analysis was used to estimate factors predicting IUBT failure.   
      \end{eg}
    \end{enumerate}

  \subsection{语阶7(2S7) 报告注册信息}

    \subsubsection{语阶分析}

    根据WHO的规定，临床试验，特别是那些前瞻性的临床试验，应该进行登记，所以一些临床试验在M2的最后会包括这一语阶。尽管它很重要，但在摘要写作中却很少见到。

    \subsubsection{语言实现}

    在这一语阶中，研究或试验通常被当作主语，所以被动语态的使用更为频繁。一般来说，一般现在时比一般过去时的使用频率高得多。以下是发现的少数2S7例子中的两个。

    \begin{eg}{}
      This trial is registered with ClinicalTrials. gov, number NCTO1815840, and the study is ongoing.
    \end{eg}

    \begin{eg}{}
      The KEYNOTE-001 trial was registered with ClinicalTrials. gov, number NCTOI295827.  
    \end{eg}

    \subsubsection{词组}

    \begin{enumerate}
      \item This trial/study is registered with
      \begin{eg}{}
        This trial is registered with ClinicalTrials. gov as NCTO0569127.
      \end{eg}

      \begin{eg}{}
        This study is registered with Current Controlled Trials, number ISRCTN82857232.
      \end{eg}
    \end{enumerate}

    \begin{task}{语料库任务}
      在您的语料库中，哪些词可以与“clinical”、“experiment”、“research”、“study”或“trial”搭配？请按出现频率列出。
    \end{task}

\section{材料阅读}

这一节提供了3段M2材料，并对其中语阶的语言特点进行了详细分析。

\begin{sample}[label={myautocounter}]{\heiti}
  
  \textbf{METHOD} 
  
  Randomized controlled trial for adolescents (12--18 years of age) with recent (past 3 months) suicide attempts or other self-harm. Youth were randomized either to SAFETY or to treatment as usual enhanced by parent education and support accessing community treatment (E-TAU). Outcomes were evaluated at baseline, 3 months, or end of treatment period, and were followed up through 6 to 12 months. The primary outcome was youth-reported incident suicide attempts through the 3-month follow-up.

  \begin{flushright}
    ---Cognitive-Behavioral Family Treatment for Suicide Attempt Prevention: A Randomized Controlled Trial. \emph{Journal of the American Academy of Child and Adolescent Psychiatry (2017)}
  \end{flushright}

  \tcblower

  \noindent \textbf{STEP IDENTIFICATION}

  \vspace*{10pt}
  {\small\noindent
  \begin{tblr}{colspec={X[1,c]X[5,l]},cell{even}{1,2} = {azure9}}
    \toprule
    \textbf{Step} & \textbf{Sample} \\ 
    \midrule
    
    2S3 with 2S2 embedded & Randomized controlled trial for adolescents (12--18 years of age) with recent (past 3 months) suicide attempts or other self-harm. \\
    2S4 & Youth were randomized either to SAFETY or to treatment as usual enhanced by parent education and support accessing community treatment (E-TAU). \\
    2S5 & Outcomes were evaluated at baseline, 3 months, or end of treatment period, and were followed up through 6 to 12 months. The primary outcome was youth-reported incident suicide attempts through the 3-month follow-up. \\
      
    \bottomrule
  \end{tblr}
  }

  \noindent \textbf{ANALYSIS}

  This is a typical one-paragraphed M2. The first language structure is a noun phrase with the research design (2S2) ``randomized controlled trial'' embedded with the description of the participants and selection criteria (2S3) ``adolescents$\dots$''. The process description (2S4) and the outcome measurement (2S5) are in passive voice with ``youth'', ``outcomes'' and ``the primary outcome'' as the sentence subjects. 2S1, 2S6 and 2S7 are missing, as they are not conventional in M2. Common lexical chunks of M2 are also seen in this paragraph, such as ``randomized controlled trial'', ``$\dots$ were randomized to $\dots$'', ``Outcomes were evaluatedd $\dots$'', ``The primary outcome was $\dots$'', etc.

\end{sample}

\begin{sample}[label={myautocounter}]{\heiti}
  
  \textbf{MATERIALS AND METHODS} 
  
  The institutional review board approved this retrospective study and waived the informed consent requirement. Seventy-four patients with surgically confirmed PNETs and 82 patients with PDACs who underwent gadobutrol-enhanced MR imaging were included. Two radiologists independently evaluated the morphologic characteristics and temporal enhancement patterns of each tumor. Quantitative analysis, including measurement of tumor size, maximal upstream parenchymal thickness (MUPT), contrast-to-noise ratio, and apparent diffusion coefficient values, was performed. Uni-and multivariate logistic regression analyses were performed to identify relevant features to differentiate between PNETs and PDACs.

  \begin{flushright}
    ---Nonhypervascular Pancreatic Neuroendocrine Tumors: Differential Diagnosis from Pancreatic Ductal Adenocarcinomas at MR Imaging-Retrospective Cross-sectional Study. \emph{Radiology (2017)}
  \end{flushright}

  \tcblower

  \noindent \textbf{STEP IDENTIFICATION}

  \vspace*{10pt}
  {\small\noindent
  \begin{tblr}{colspec={X[1,c]X[5,l]},cell{even}{1,2} = {azure9}}
    \toprule
    \textbf{Step} & \textbf{Sample} \\ 
    \midrule
    
    2S1 with 2S2 embedded  & The institutional review board approved this retrospective study and waived the informed consent requirement.\\
    2S3  & Seventy-four patients with surgically confirmed PNETs and 82 patients with PDACs who underwent gadobutrol-enhanced MR imaging were included.\\
    2S4  & Two radiologists independently evaluated the morphologic characteristics and temporal enhancement patterns of each tumor.\\
    2S5  & Quantitative analysis, including measurement of tumor size, maximal upstream parenchymal thickness (MUPT), contrast-to-noise ratio, and apparent diffusion coefficient values, was performed. \\
    2S6  & Uni-and multivariate logistic regression analyses were performed to identify relevant features to differentiate between PNETs and PDACs. \\
      
    \bottomrule
  \end{tblr}
  }

  \noindent \textbf{ANALYSIS}

  Although it is in a single one paragraph, this M2 includes almost every step listed above (except for 2S7). The first sentence offers the ethical investigation information(2S1) embeddedwith the research design ``retrospective study''(2S2), followed by the participants and selection criteria in passive voice ``seventy-four patients $\dots$ were included''(2S3) and the research process in active voice ``Two radiologists $\dots$ evaluated $\dots$''(2S4). Several specific measurements ``$\dots$ of tumor size, maximal upstream parenchymal thickness,$\dots$''(2S5) and data analysis ``uni-and multivariate logistic regression analysis $\dots$''(2S6) are mentioned next. The passive voice is quite prominent in this sample, indicating the objectivity of the research process. Signals for M2 covered in this example include ``retrospective study'', ``informed consent'', ``$\dots$ patients $\dots$ were included'', ``analysis, $\dots$ was performed'', ``logistic regression analyses were performed'', etc.

\end{sample}


\begin{sample}[label={myautocounter}]{\heiti}
  
  \textbf{DESIGN} 
  
  Cross-sectional observational study.

  \textbf{PARTICIPANTS }
  
  One hundred nineteen children 2 to 16 years of age (mean age, 9.4 years; standard deviation [SD], 4.56 years) with glaucoma and their parents.

  \textbf{METHODS }
  
  Completion of 3 validated instruments for children to assess (1) functional visual ability (FVA) with the Cardiff Visual Ability Questionnaire for Children (CVAQC), (2) VR QoL with the Impact of Vision Impairment for Children (IVI-C), and (3) HR QoL with the Pediatric Quality of Life Inventory (PedsQL) version 4.0.

  \textbf{MAIN OUTCOME MEASURES} 
  
  Cardiff Visual Ability Questionnaire for Children, IVI-C, and PedsQL scores.

  \begin{flushright}
    ---Quality of Life and Functional Vision in Children with Glaucoma.
\emph{Ophthalmology (2017)}
  \end{flushright}

  \tcblower

  \noindent \textbf{STEP IDENTIFICATION}

  \vspace*{10pt}
  {\small\noindent
  \begin{tblr}{colspec={X[1,c]X[5,l]},cell{even}{1,2} = {azure9}}
    \toprule
    \textbf{Step} & \textbf{Sample} \\ 
    \midrule
    
    2S2 & {DESIGN: \\Cross-sectional observational study.} \\
    2S3 & {PARTICIPANTS: \\One hundred nineteen children 2 to 16 years of age (mean age, 9.4 years; standard deviation [ SD], 4.56 years) with glaucoma and their parents.} \\
    2S4 & {METHODS: \\Completion of 3 validated instruments for children to assess (1) functional visual ability (FVA) with the Cardiff Visual Ability Questionnaire for Children (CVAQC), (2) VR QoL with the Impact of Vision Impairment for Children (IVI-C), and (3) HR QoL with the Pediatric Quality of Life Inventory (PedsQL) version 4.0.} \\
    2S5 & {MAIN OUTCOME MEASURES: \\Cardiff Visual Ability Questionnaire for Children, IVI-C, and PedsQL scores.} \\
      
    \bottomrule
  \end{tblr}
  }


  \noindent \textbf{ANALYSIS}

  This is a typical multi-paragraphed M2. With subtitles, the steps are quite clear. In such a structured M2, noun phrases are very common. In the "DESIGN"step (2S2), some descriptive medical terms ``cross-sectional'' and ``observational'' are used to modify ``study''. In the ``PARTICIPANTS'' step (2S3), number and selection criteria of the participants are presented in noun or prepositional phrases. In the ``METHODS'' step (2S4), the examinations or assessments are described. In the ``MAIN OUTCOME MEASURES'' step (2S5), the tool used to observe is mentioned.

\end{sample}

\section{词汇表}

{\small
\begin{longtblr}[
    caption = {第四章词汇表},
    label = {tab:Glossary of Chapter 4},
    % note{a} = {英文论文中指代当前文献中的差距、问题或缺陷。即现有研究尚未解决的部分。},
]{
    width = \textwidth,
    colspec = {X[1,l,h]  X[1,l,h]  X[3,l,h]},
    rowhead = 1, rowfoot = 0, % 每个分页里表头表尾的数量
    % row{odd} = {blue8}, 
    row{even} = {azure9},
}
    
\toprule
\textbf{WORDS} & \textbf{MEANING} & \textbf{MEANING or EXAMPLE}\\
\midrule

{\textbf{collocate}\\/\textipa{"k6l@""keIt}/} & \emph{v.} 【语】词语的组合；排列 & to group or place together in some system or order \\
\textbf{embed}/\textipa{Im"beEd}/ & \emph{v.} 嵌入 & to set or fix firmly in a surrounding mass\\
{\textbf{extraction}\\/\textipa{Ik"str\ae kS@n}/} & \emph{n.} 提取；开采；提炼；拔出 & the act or process of removing or obtaining sth from sth else\\
{\textbf{intervention}\\/\textipa{Int@"benS@n}/} & \emph{n.} 干预；介入；调解 & the act of intervening\\
{\textbf{registration}\\/\textipa{""rE\textdyoghlig I"streIS@n}/} & \emph{n.} 登记 & the act of making an official record of sth/sb \\
{\textbf{regression}\\/\textipa{rI"grES@n}/} & \emph{n.} 回归；退化；倒退 & the process of going back to an earlier or less advanced form or state \\
{\textbf{tentative}\\/\textipa{"tEnt@tIv}/} & \emph{n.} 试探；尝试；实验 & not definite or certain because you may want to change it later \\

\bottomrule

\end{longtblr}
}

\begin{problemset}
  \item \textbf{With the ``Steps in M2'' offered, try to identify the steps in the following 2 examples of M2 and write down the corresponding numbers in the blanks in front of each part.}
  
  \hspace*{2em}Steps in M2:

  \hspace*{2em}2S1 Reporting on medical ethics review
  
  \hspace*{2em}2S2 Explaining briefly research design

  \hspace*{2em}2S3 Describing subjects or data and their selection criteria

  \hspace*{2em}2S4 Describing experimental procedure, such as interventions, examinations, etc.

  \hspace*{2em}2S5 Describing main outcomes and their measures

  \hspace*{2em}2S6 Describing data analysis methods

  \hspace*{2em}2S7 Reporting on registration information

  \textbf{Example 1}

  \hspace*{2em}\uline{\hspace*{3em}}Design: Retrospective case-control study.

  \hspace*{2em}\uline{\hspace*{3em}}Setting: Nijmegen, the Netherlands.

  \hspace*{2em}\uline{\hspace*{3em}}Patients: Thirty consecutive patients with SSD.
  
  \hspace*{2em}\uline{\hspace*{3em}}Interventions: Patients received a trial with a BCD headband as part of the regular workup for SSD. The patients were divided into 2 groups according to their decision to opt for a BCD (BCD+) or not (BCD-).
  
  \hspace*{2em}\uline{\hspace*{3em}}Main outcome measures: Patients completed a questionnaire on satisfaction with the BCD headband, patient-and BCD-related factors, and benefit in listening situations.
  
  \textbf{Example 2}
  
  \hspace*{2em}Methods: 
  
  \hspace*{2em}\uline{\hspace*{3em}} Eighteen adults with first-ever chronic monohemispheric subcortical stroke participated in \uline{\hspace*{3em}} this randomized, controlled, triple-blinded trial. \uline{\hspace*{3em}} Intervention consisted of priming with real or sham iTBS to the ipsilesional primary motor cortex immediately before 45 minutes of upper limb physical therapy, daily for 10 days. \uline{\hspace*{3em}} Changes in upper limb function (Action Research Arm Test [ARAT]), upper limb impairment (Fugl-Meyer Scale), and corticomotor excitability, were assessed before, during, and immediately, 1 month and 3 months after the intervention. Functional magnetic resonance images were acquired before and at one month after the intervention.

  \item \textbf{Fill in the blanks with the correct form of the verbs given. Notice the tense and voice.}
  
  \begin{enumerate}
    \item Methods: A population-based prospective cohort study\uline{~~~1)~~~}(include) 3,504 male and female Koreans aged 40 to 69 years from the Korean Genome Epidemiology Study. At the beginning of follow-up, all individuals \uline{~~~2)~~~}(be) free of metabolic syndrome and known cardiovascular disease. Each participant\uline{~~~3)~~~}(complete) a food frequency questionnaire. Incident cases of metabolic syndrome \uline{~~~4)~~~}(identify) by biennial health examinations during a follow-up period between April 17, 2003, and November 17, 2006.
    Pooled logistic regression analysis \uline{~~~5)~~~} (apply) to obtain an odds ratio (OR) of metabolic syndrome with its 95\% confidence interval for fish or n-3 fatty acid intake.
    
    \item Methods: This\uline{~~~1)~~~}(be) a retrospective matched cohort study of women with liver cirrhosis between January 2005 and January 2016 in a university hospital. Women in a case group\uline{~~~2)~~~}(match) to women in a control group according to year of delivery, age, body mass index, and parity in a 1:4 ratio. Bivariable and multivariable analyses \uline{~~~3)~~~}(perform) \uline{~~~4)~~~}(compare) the prevalence of the primary composite outcome, which \uline{~~~5)~~~}(include) any one of the following: fetal or neonatal demise, placental abruption, preeclampsia, preterm delivery at less than 37 weeks of gestation, and small-for-gestational age neonate between women in the case group and those in the control group.
    
  \end{enumerate}

  \item The following steps in M2 need improvement. You may have to change a word, add a word or delete a word. Mark out the mistakes and put the corrections in the blank provided. Ifyou change a word, underline it and write the correct word in the corresponding blank. If youd add a word, put an insertion mark ($\wedge$) in the right place and write the added word in the blank. If you delete a word, cross it out and put a slash (/) in the blank. If you consider the sentence as correct, put a tick ($\surd$).
  
  \hspace*{2em}Design, setting and participants: A randomize clinical trial of patients in\sidenote{(1)\uline{\hspace*{3em}}} persistent hypercapnia (Pa\ce{CO2}$>53$ mmHg) 2 weeks to 4 weeks after resolution\sidenote{(2)\uline{\hspace*{3em}}} of respiratory acidemia, who recruited from 13 UK centers between 2010 and\sidenote{(3)\uline{\hspace*{3em}}} 2015. Exclusion criterium included obesity (body mass index $[BMI] >35$),\sidenote{(4)\uline{\hspace*{3em}}} obstructive sleep apnea syndrome, or other causes of respiratory failure. Of 2021 patients screening, 124 were eligible.\sidenote{(5)\uline{\hspace*{3em}}}

  \hspace*{2em}Interventions: There were 59 patients randomized as to home oxygen alone\sidenote{(6)\uline{\hspace*{3em}}} (medial oxygen flow rate, 1.0 L/min [ interquartile range (IQR), 0.5--2.0\sidenote{(7)\uline{\hspace*{3em}}} L/min]) and 57 patients to home oxygen plus home NIV (median oxygen flow rate, 1. 0 L/min [IQR, 0.5--1.5 L/min]). The median home\sidenote{(8)\uline{\hspace*{3em}}} ventilator settings are an inspiratory positive airway pressure of 24 (IQR, 22--\sidenote{(9)\uline{\hspace*{3em}}}26) cm H2O, expiratory positive airway pressure of 4 (IQR, 4--5) cm H2O, and a backup rate of 14 (IQR, 14--16) breaths/minute.

  \hspace*{2em}Main outcomes and measures: Time to readmission or death within 12 months adjusted for the number of previous COPD admissions, previous use of longterm\sidenote{(10)\uline{\hspace*{3em}}} oxygen, age, and BMI.
  

\end{problemset}



\chapter{语步三}\label{chapter5}

语步3（M3）是摘要的第三部分。这一部分的交流目的是总结研究的结果或成果。它是摘要中最长的部分。不过，M3在结构上并不具有层级性，这与摘要中的其他三个语步不同。

\section{M3中的语阶}

M3中发现的语阶包括：

3S1 提供有效样本的信息

3S2 说明整体观察或主要结果 

除了极少数个例外，这两个语阶均按照上面列出的顺序书写，以实现作者的交流目的（Sample~\ref{sam:M3-1}）。在这两个语阶中，3S1是可选的，而3S2是常规的（Sample~\ref{sam:M3-2}）。尽管3S2很长，但很难对其进一步分层。

\begin{sample}[label={sam:M3-1}]{\heiti}

  \textbf{RESULTS}

  60 individuals were recruited over a 9-month period. Three withdrew, leaving 29 intervention and 28 controls participants in the final analysis. 32\% of patients with FMS met the inclusion criteria, of which 90\% enrolled. Acceptability of the intervention was high and there were no adverse events. At 6 months, 72\% of the intervention group rated their symptoms as improved, compared to 18\% in the control group. There was a moderate to large treatment effect across a range of outcomes, including three of eight Short Form 36 (SF36) domains $(d=0.46-0.79$). The SF36 Physical function was found to be a suitable primary outcome measure for a future trial; adjusted mean difference 19.8(95\% CI 10.2 to 29.5). The additional quality adjusted life vears (OALY) with intervention was 0.08(95\% CI 0. 03 to 0. 13), the mean incremental cost per QALY gained was \pounds{}12 087.

  
  \begin{flushright}
    ---Randomised feasibility study of physiotherapy for patients with functional motor symptoms. \emph{Journal of Neurology, Neurosurgery, and Psychiatry (2016)}
  \end{flushright}

  \tcblower

  \noindent \textbf{STEP IDENTIFICATION}

  \vspace*{10pt}
  {\small\noindent
  \begin{tblr}{colspec={X[1,c]X[5,l]},cell{even}{1,2} = {azure9}}
    \toprule
    \textbf{Step} & \textbf{Sample} \\ 
    \midrule
    
    3S1 & 60 individuals were recruited over a 9-month period. Three withdrew, leaving 29 intervention and 28 controls participants in the final analysis. 32\% of patients with FMS met the inclusion criteria, of which 90\% enrolled.\\
    3S2 & Acceptability of the intervention was high and there were no adverse events. At 6 months, 72\% of the intervention group rated their symptoms as improved, compared to 18\% in the control group. There was a moderate to large treatment effect across a range of outcomes, including three of eight Short Form 36 (SF36) domains $(d=0.46-0.79$). The SF36 Physical function was found to be a suitable primary outcome measure for a future trial; adjusted mean difference 19.8(95\% CI 10.2 to 29.5). The additional quality adjusted life vears (OALY) with intervention was 0.08(95\% CI 0. 03 to 0. 13), the mean incremental cost per QALY gained was \pounds{}12 087.\\
      
    \bottomrule
  \end{tblr}
  }


\end{sample}

\begin{sample}[label={sam:M3-2}]{\heiti}
  
  \textbf{RESULTS }
  
  Scores for FVA, VR QoL, and HR QoL were reduced in children with glaucoma: median CVAQC score, -1. 24 (interquartile range [ IQR],-2.2 to -0.11; range, -3. 00 higher visual ability to +2. 80 lower visual ability); mean IVI-C score, 67.3 (SD, 14.4; normal VR QoL, 96); median PedsQL self-report, 78. 8 (IQR, 67.4--90. 2); parent report, 71.2 (IQR, 55.7--85.8); and family impact score, 74.3 (IQR, 56.9-88.5; normal HR QoL, 100). Psychosocial subscores were lower than physical subscores on the PedsQL. Older children reported less impairment on CVAQC, IVI-C, and PedsQL than younger children. Parents reported greater impact on their child's HR QoL than children reported themselves.


  \begin{flushright}
    ---Quality of Life and Functional Vision in Children with Glaucoma. \emph{Ophthalmology (2017)}
  \end{flushright}

  \tcblower

  \noindent \textbf{STEP IDENTIFICATION}

  \vspace*{10pt}
  {\small\noindent
  \begin{tblr}{colspec={X[1,c]X[5,l]},cell{even}{1,2} = {azure9}}
    \toprule
    \textbf{Step} & \textbf{Sample} \\
    \midrule
    
    3S2 & Scores for FVA, VR QoL, and HR QoL were reduced in children with glaucoma: median CVAQC score, -1. 24 (interquartile range [ IQR],-2.2 to -0.11; range, -3. 00 higher visual ability to +2. 80 lower visual ability); mean IVI-C score, 67.3 (SD, 14.4; normal VR QoL, 96); median PedsQL self-report, 78. 8 (IQR, 67.4--90. 2); parent report, 71.2 (IQR, 55.7--85.8); and family impact score, 74.3 (IQR, 56.9-88.5; normal HR QoL, 100). Psychosocial subscores were lower than physical subscores on the PedsQL. Older children reported less impairment on CVAQC, IVI-C, and PedsQL than younger children. Parents reported greater impact on their child's HR QoL than children reported themselves.\\
      
    \bottomrule
  \end{tblr}
  }

\end{sample}

\section{每个语阶的语言特征}

与M2中的大多数语阶类似，M3使用过去式，并且同时使用主动语态和被动语态。

  \subsection{语阶1(3S1) 提供有效样本的信息}
    \subsubsection{语阶分析}

    在基于实验的研究中，可以补充关于有效样本的额外信息，如数量、年龄和分组等。

    \subsubsection{语言实现}

    有时，3S1使用过去式和被动语态，样本作为句子主语。有时，它则采用主动语态。

    \begin{eg}{}
      Of 1072 neonates screened, 523 \uline{were assigned} to hydrocortisone (n= 256) or placebo (n=267) and 406 \uline{survived} to 2 years of age.  
    \end{eg}

    \subsubsection{词组}

    \begin{enumerate}
      \item (patients) were randomly assigned to (receive)
      \begin{eg}{}
        Of 381 patients screened, 142 eligible \uline{patients were randomly assigned to} treatment (72 to the vandetanib group and 70 to the placebo group).
      \end{eg}

      \begin{eg}{}
        During this time, 239 infants were assessed and 181 eligible infants \uline{were randomly assigned to receive} an LMA (n=85) or an endotracheal tube (n=95).    
      \end{eg}

      \item (patients) were included in the (study)
      \begin{eg}{}
        In all, 198 \uline{patients were included in the} final model.
      \end{eg}

      \begin{eg}{}
        A total of 702 women \uline{were included in the study}, with 29 (4.1\%) maternal deaths, and a mortality ratio of 56. 98 deaths per 100,000 live births.  
      \end{eg}

      \begin{eg}{}
        These variables \uline{were included in the} RSS and assigned scores ranging from 0 to 6. 
      \end{eg}

      \item were included in the analysis
      \begin{eg}{}
        A total of 4011 eyes of 2057 subjects with T2DM \uline{were included in the analysis}.
      \end{eg}

      \item patients were enrolled and
      \begin{eg}{}
        Between Jan 8,2013, and Jan 31,2014,495 eligible adult \uline{patients were enrolled and} randomly assigned to the cobimetinib plus vemurafenib group (n=247) or placebo plus vemurafenib group (n=248).   
      \end{eg}
    \end{enumerate}

  \subsection{语阶2(3S2) 说明整体观察或主要结果}
    \subsubsection{语阶分析}

    作为M3的一个重要语阶，3S2要阐述研究的主要发现或结果。一般来说，研究结果的概述会放在在细节之前，我们可以从下面的例子中看到。

    \begin{eg}{}
      \par 

      \vspace*{10pt}
      {\small\noindent
      \begin{tblr}{colspec={X[1,c]X[5,l]},cell{even}{1,2} = {azure9},cell{odd}{1,2} = {white}}
        \toprule
        \textbf{Step} & \textbf{Sample} \\ 
        \midrule
        
        1S3a & The authors investigated whether parecoxib-supplemented IV morphine analgesia could decrease the incidence of delirium in elderly patients after total hip or knee replacement surgery.\\
        & $\dots$\\
        3S2 & The incidence of delirium was significantly reduced from 11.0\% (34/310) with placebo to 6. 2\% (19/310) with parecoxib (relative risk 0.56, 95\% confidence interval $0.33-0.96$, $P=.031$). The severity of pain and the cumulative consumptions of morphine at 24, 48, and 72 hours after surgery were significantly lower with parecoxib than with placebo (all $P<. 001$), although the differences were small. There was no difference in the incidence of postoperative complications between the 2 groups (12.3\%[38/310] with placebo versus
        11.6\% [36/310] with parecoxib; P=.80).\\
          
        \bottomrule
      \end{tblr}
      }
    \end{eg}

    本例首先回答了1S3a中的研究目的和研究问题“帕瑞昔布辅助静脉吗啡镇痛能否降低老年患者全髋关节或膝关节置换手术后谵妄的发生率”，然后从“疼痛的严重程度”、“累计消耗量”和“术后并发症”这几个方面阐述研究结果的意义。

    \subsubsection{语言实现}

    过去式是3S2最明显的语言特征。此外，被动语态比主动语态更受青睐，这样可以表明研究结果的客观性。我们发现，用于显示结果的动词也是这语阶的信号词，如“found”、“observed”等。不过，还是有使用主动语态的情况存在。例如，可以采用“analysis showed”的结构，以中立的、听起来很科学的语气提升研究分析的可信度。

    \begin{eg}{}
      No change in measures of daily PA was observed at 4-wk compared with baseline($P>0.05$)   
    \end{eg}

    \begin{eg}{}
      Multivariated \uuline{analysis} \uline{showed} that nodules less than 1. 5 cm, ultrasonographic findings suggestive of malignancy and more than 2 results of atypia from repeated FNAs were significant risk factors for malignancy ($P<0.001$).  
    \end{eg}

    医学研究中经常涉及相关性和组间差异等，因此关系性信号词构建了各要素之间的联系，其中包括结果性信号词 (e.g. \emph{these findings suggest that})、对比性信号词(e.g. \emph{no significant difference in}) 和推断性信号词 (e. g. \emph{was found to be})。

    此外，表示认识论模式的形容词，如“likely”，表达了可能性的程度，也是该语阶的特点。它们显示了非个人的认识论立场，以表达研究结果中的可能性程度。

    \begin{eg}{}
      Recidivists were more \uline{likely} to be male (P<0.0001), Black (P<0.0001), have a blood alcohol content above 80 mg/dL (P<0.0001) compared with nonrecidivists.  
    \end{eg}

    \subsubsection{词组}

    \begin{enumerate}
      \item (there was) no (significant) difference (in/between)
      \begin{eg}{}
        However, \uline{there was no significant difference in} the per cent of change in those parameters. 
      \end{eg}

      \begin{eg}{}
        After matching, \uline{there was no difference in} the 30-day rate of AMI between testing modalities.  
      \end{eg}

      \begin{eg}{}
        Model 1 showed \uline{no significant difference between} predicted and observed events (risk ratio (RR)=0. 87,95\% Cl0. 16--4.62).
      \end{eg}

      \item (there were) no (significant) differences (in/between)
      \begin{eg}{}
        \uline{There were no significant differences in} the duration of the surgical procedure ($P=0.12$), weight of the surgical specimen ($P=0.54$) or the patients' pain perception ($P=0.28$).
      \end{eg}

      \begin{eg}{}
        \uline{There were no significant differences between} men and women. 
      \end{eg}

      \begin{eg}{}
        The interventions did not raise any safety concerns and \uline{there were no differences between} groups in serious or other adverse events.  
      \end{eg}

      \begin{eg}{}
        \uline{There were no differences in} secondary outcomes except for clinician satisfaction with ease of administration.
      \end{eg}

      \item there was/were no significant
      \begin{eg}{}
        \uline{There was no significant} interaction between BNP and the 2 independent variables (P=.60, and P=.90), respectively.        
      \end{eg}

      \begin{eg}{}
        \uline{There were no significant} changes in number of total and regional nevi count and in the dermoscopic features of nevi between biological and conventional treatment groups.
      \end{eg}

      \item (was/were) not significantly different (between)
      \begin{eg}{}
        Results showed that compared to PWNS, PWS \uline{were not significantly different} in matching either the phase (timing) or the amplitude of the target in both jaw and hand tracking of predictable and unpredictable targets.
      \end{eg}

      \begin{eg}{}
        The allele frequency of this polymorphism in individuals originating from two locations with different malaria endemicity in the past \uline{was not significantly different}.  
      \end{eg}

      \begin{eg}{}
        ... and these rates were \uline{not significantly different between} groups A and B ($P=0. 346$ and $P=0. 370$, respectively).          
      \end{eg}

      \item did not differ (significantly) between (the)
      \begin{eg}{}
        Intraclass correlation coefficients (ICCs) for interreader and intrareader agreement \uline{did not differ significantly between} measurements for FA and MD. 
      \end{eg}

      \begin{eg}{}
        No serious adverse reactions were recorded and other safety measures \uline{did not differ between the} groups, after allowing for missing data. 
      \end{eg}

      \begin{eg}{}
        Complications \uline{did not differ between} devices.
      \end{eg}

      \item was/were (not) (significantly/independently) associated with (a/an)
      \begin{eg}{}
        Additionally, in subgroup analyses, rescue breathing \uline{was not associated with} neurological outcome regardless of the type of rescuer [family member: adjusted OR $0. 83$ (95\% CI 0.39--1.70); or non-family member; adjusted OR 1.91 (95\% CI 0.79--5.35)].  
      \end{eg}

      \begin{eg}{}
        AD \uline{was significantly associated with} 11 of 22 examined autoimmune diseases.   
      \end{eg}

      \begin{eg}{}
        Hypotension \uline{was independently associated with} mortality, acute kidney injury and hospital admission.    
      \end{eg}

      \item (serious) adverse events were reported (in)
      \begin{eg}{}
        Treatment-emergent \uline{adverse events were reported in} 103 patients (97\%), a majority of which were grade 1 to 2 in severity.
      \end{eg}

      \begin{eg}{}
        No \uline{serious adverse events were reported}.    
      \end{eg}

      \item were more/less likely to (have)
      \begin{eg}{}
        Survivors of cancer \uline{were more likely to have} chronic conditions and MCCs compared with adults without a history of cancer. 
      \end{eg}

      \begin{eg}{}
        Men and higher-educated people \uline{were more likely to} be highly sedentary, while women and lower-educated people \uline{were more likely to} be inactive.
      \end{eg}

      \item (was) (significantly) higher/lower in (the)$\dots$
      \begin{eg}{}
        Percent MVIC retained \uline{was significantly higher in} ACB patients at 6 ($P<0. 000 1$) and 24 hours ($P<0. 000 1$).         
      \end{eg}

      \begin{eg}{}
        PCT levels were \uline{significantly higher in the} population with Gram-negative rod (GNR) infections than in the population with Gram-positive coccal (GPC) infections.         
      \end{eg}

      \item was found to be
      \begin{eg}{}
        The SF36 Physical function \uline{was found to be} a suitable primary outcome measure for a future trial; adjusted mean difference 19. 8(95\% Cl 10.2 to 29.5).   
      \end{eg}

      \begin{eg}{}
        No statistically significant difference in overall survival was observed \uline{between the two groups} (102.7 vs 115.7 months, respectively). 
      \end{eg}
    \end{enumerate}

\section{材料阅读}

这一小节提供了2段M3材料，并对其中语阶的语言特点进行了详细分析。

\begin{sample}[label={myautocounter}]{\heiti}
  
  RESULTS

  60 individuals were recruited over a 9-month period. Three withdrew, leaving 29 intervention and 28 controls participants in the final analysis. 32\% of patients with FMS met the inclusion criteria, of which 90\% enrolled. Acceptability of the intervention was high and there were no adverse events. At 6 months, 72\% of the intervention group rated their symptoms as improved, compared to 18\% in the control group. There was a moderate to large treatment effect across a range of outcomes, including three of eight Short Form 36 (SF36) domains ($d=0.46-0.79$). The SF36 Physical function was found to be a suitable primary outcome measure for a future trial; adjusted mean difference 19. 8 (95\% CI 10.2 to 29.5). The additional quality adjusted life years (QALY) with intervention was 0.08 (95\% CI 0.03 to 0.13), the mean incremental cost per QALY gained was \pounds{}12 087.


  \begin{flushright}
    ---Randomised feasibility study of physiotherapy for patients with functional motor symptoms. \emph{Journal of Neurology, Neurosurgery, and Psychiatry (2016)}
  \end{flushright}

  \tcblower

  \noindent \textbf{STEP IDENTIFICATION}

  \vspace*{10pt}
  {\small\noindent
  \begin{tblr}{colspec={X[1,c]X[5,l]},cell{even}{1,2} = {azure9}}
    \toprule
    \textbf{Step} & \textbf{Sample} \\ 
    \midrule
  
    3S1 & 60 individuals were recruited over a 9-month period. Three withdrew, leaving 29 intervention and 28 controls participants in the final analysis. 32\% of patients with FMS met the inclusion criteria, of which 90\% enrolled.\\
    3S2 & Acceptability of the intervention was high and there were no adverse events. At 6 months, 72\% of the intervention group rated their symptoms as improved, compared to 18\% in the control group. There was a moderate to large treatment effect across a range of outcomes, including three of eight Short Form 36 (SF36) domains ($d=0.46-0.79$). The SF36 Physical function was found to be a suitable primary outcome measure for a future trial; adjusted mean difference 19. 8 (95\% CI 10.2 to 29.5). The additional quality adjusted life years (QALY) with intervention was 0.08 (95\% CI 0.03 to 0.13), the mean incremental cost per QALY gained was \pounds{}12 087.\\
    
    \bottomrule
  \end{tblr}
  }

  \noindent \textbf{ANALYSIS}

  This is an M3 sample with two steps in past tense. The first step 3S1 offers extra information about the valid samples, that is, 90\% of those with FMS who met the inclusion criteria. Passive voice is employed in this step with ``60 individuals'' as the subject. The second step 3S2 contains the major findings beginning with an overview, that is, ``acceptability of the intervention was high and there were no adverse events.'' More data of the outcome are offered in the following part, with some signals and lexical chunks commonly seen in 3S2, such as ``compared to'', ``in the control group'', ``was found to be'', ``a suitable primary outcome measure'', and ``mean difference'', etc. Both active voice (e.g. ``72\% of the intervention group rated their symptoms as$\dots$'') and passive voice (e.g. ``the SF36 Physical function was found to be$\dots$'') are used.

\end{sample}

\begin{sample}[label={myautocounter}]{\heiti}

  \textbf{RESULTS}
  
  Scores for FVA, VR QoL, and HR QoL were reduced in children with glaucoma: median CVAQC score,-1.24 (interquartile range [IQR], -2.2 to -0.11; range, -3.00 higher visual ability to +2.80 lower visual ability); mean IVI-C score, 67.3 (SD, 14.4; normal VR QoL, 96); median PedsQL self-report, 78.8 (IQR, 67.4-90.2); parent report, 71.2 (IQR, 55.7--85.8); and family impact score, 74.3 (IQR, 56.9--88.5; normal HR QoL, 100). Psychosocial subscores were lower than physical subscores on the PedsQL. Older children reported less impairment on CVAQC, IVI-C, and PedsQL than younger children. Parents reported greater impact on their child's HR QoL than children reported themselves.

  
  \begin{flushright}
    ---Quality of Life and Functional Vision in Children with Glaucoma. \emph{Ophthalmology (2017)}
  \end{flushright}

  \tcblower

  \noindent \textbf{STEP IDENTIFICATION}

  \vspace*{10pt}
  {\small\noindent
  \begin{tblr}{colspec={X[1,c]X[5,l]},cell{even}{1,2} = {azure9}}
    \toprule
    \textbf{Step} & \textbf{Sample} \\ 
    \midrule
  
    3S2 & Scores for FVA, VR QoL, and HR QoL were reduced in children with glaucoma: median CVAQC score,-1.24 (interquartile range [IQR], -2.2 to -0.11; range, -3.00 higher visual ability to +2.80 lower visual ability); mean IVI-C score, 67.3 (SD, 14.4; normal VR QoL, 96); median PedsQL self-report, 78.8 (IQR, 67.4-90.2); parent report, 71.2 (IQR, 55.7--85.8); and family impact score, 74.3 (IQR, 56.9--88.5; normal HR QoL, 100). Psychosocial subscores were lower than physical subscores on the PedsQL. Older children reported less impairment on CVAQC, IVI-C, and PedsQL than younger children. Parents reported greater impact on their child's HR QoL than children reported themselves.\\
    
    \bottomrule
  \end{tblr}
  }

  \noindent \textbf{ANALYSIS}

  This M3 is made up of only one step (3S2). An overview comes first, that is, ``Scores for FVA, VR QoL, and HR QoL were reduced in children with glaucoma'', followed by a presentation of a lot of data and their significance. Comparisons of subscores and of different groups (older children vs. younger children, parents vs. children) are made in the latter half of this step. As for the language realization, it is obvious that this step is in the past tense, and when the subject of the research, that is, ``scores'' is taken as the sentence subject, the passive voice is preferred to indicate the objectivity of the findings, but the active voice is also used, as we can see from ``older children reported$\dots$'' and ``parents reported$\dots$''.
  
\end{sample}

\section{词汇表}

{\small
\begin{longtblr}[
    caption = {第五章词汇表},
    label = {tab:Glossary of Chapter 5},
]{
    width = \textwidth,
    colspec = {X[1,l,h]  X[1,l,h]  X[3,l,h]},
    rowhead = 1, rowfoot = 0, % 每个分页里表头表尾的数量
    % row{odd} = {blue8}, 
    row{even} = {azure9},
}
    
\toprule
\textbf{WORDS} & \textbf{MEANING} & \textbf{MEANING or EXAMPLE}\\
\midrule

{\textbf{adverse}\\/\textipa{\ae d"v3rs}/} & \emph{adj.} 不利的；有害的；反面的 & negative and unpleasant; not likely to produce a good result \\
{\textbf{correlation}\\/\textipa{""k0r@"leIS@n}/} & \emph{n.} 相关；关联；相互关系 & a connection between two things in which one thing changes as the other does \\
{\textbf{epistemic}\\/\textipa{""EpI"stimIk}/} & \emph{adj.} 认识的；知识的 & of or pertaining to knowledge or the conditions for acquiring it \\
{\textbf{hierarchical}\\/\textipa{""haI@r"arkIk@l}/} & \emph{adj.} 阶层式的；分层的；分等级的；层级式的 & classified according to various criteria into successive levels or layers \\
{\textbf{stratify}\\/\textipa{"str\ae t@""faI}/} & \emph{v.} （使）分层 & to arrange sth in layers \\

\bottomrule

\end{longtblr}
}

\begin{problemset}
  \item \textbf{Fill the blanks with the correct forms of the given words. Notice the tense and voice.}
  \begin{enumerate}
    \item Results: The utilization of regional anesthetic techniques \uline{~~~1)~~~}(not differ) by OSA status and overall $<25\%$ and $15\%$ \uline{~~~2)~~~}(receive) neuraxial anesthesia and peripheral nerve blocks, respectively. Trend analysis \uline{~~~3)~~~}(show) a significant increase in peripheral nerve block use by $>50\%$ and a concurrent decrease in opioid prescription. Interestingly, while the absolute number of patients with OSA receiving perioperative oximetry, supplemental oxygen, and positive airway pressure therapy significantly \uline{~~~4)~~~}(increase) over time, the proportional use significantly \uline{~~~5)~~~}(decrease) by approximately 28\%, 36\%, and 14\%, respectively. Ashift from utilization of intensive care to telemetry and stepdown units \uline{~~~6)~~~}(see).
    
    \item Results: After controlling for potential cardiovascular risk factors, multivariate OR for metabolic syndrome \uline{~~~1)~~~}(be) 0.43(95\% CI 0.23 to 0.83) for men who \uline{~~~2)~~~}(eat) fish daily when compared with those eating fish less than once a week. Similarly, metabolic syndrome risk \uline{~~~3)~~~}(halve) for men in the top decile of n-3 fatty acid intake when compared with those in the bottom decile (OR 0.53,95\% CI 0.28 to 0.99). In particular, fish intake \uline{~~~4)~~~}(associate) significantly with triglyceride level and high-density lipoprotein cholesterol level among the metabolic syndrome components. For women, apparent associations \uline{~~~5)~~~}(not observe) between fish intake or n-3 fatty acid intake and metabolic syndrome risk.
  \end{enumerate}

  \item \textbf{``Significance'' and ``difference'' are two of the key words in M3, so we should be familiar with their different forms. Please fill in the blanks with the proper forms of the words given.}
  \begin{enumerate}
    \item \fbox{significance, significant, significantly}
    \begin{enumerate}
      \item Overall survival with pembrolizumab was \uline{\hspace*{3em}} longer in patients who previously received any radiotherapy than in patients without previous radiotherapy.
      \item This benefit was consistent across all subgroups examined (all $P<0.05$), and no \uline{\hspace*{3em}} heterogeneity of treatment effect was observed (all $P>0.05$).
      \item Different trajectories and \uline{\hspace*{3em}} of B-type natriuretic peptide, congestion and acute kidney injury in patients with heart failure.
      \item Nonetheless, \uline{\hspace*{3em}} deficits in knowledge, particularly for the diagnostic criteria for delirium, remained.
      \item Both the amount of ICG injected ($P<0.001$) and the experimental temperature ($P<0.001$) \uline{\hspace*{3em}} affected the measurements.
      \item Further, there were no \uline{\hspace*{3em}} between-group differences in motor practice effects for either jaw or hand tracking.
    \end{enumerate}

    \item \fbox{differ, difference(s), different, differently}
    \begin{enumerate}
      \item These rates were not significantly \uline{\hspace*{3em}} between groups A and B (P=0.346 and P=0.370, respectively).
      \item \uline{\hspace*{3em}} in health insurance benefit phase, drug choice, brand name, and coverage type were the greatest determinants of patient cost (P<0.001).
      \item US medical examiners and coroners apply variable practices to classify and investigate SUID, and thus, they certify the same deaths \uline{\hspace*{3em}}.
      \item Total costs of a 30-day supply from a specialist \uline{\hspace*{3em}} from family and internal medicine physicians by \${}7.36--\${}14.57.
      \item Additionally, three mutations in three \uline{\hspace*{3em}} genes were found.
      \item The two groups of patients \uline{\hspace*{3em}} for comorbidities and treatment profile.
      \item We examined whether piRNAs are expressed \uline{\hspace*{3em}} between AD cases and controls and explored the potential regulatory effects of risk SNPs on piRNA expression levels.
      \item No serious adverse reactions were recorded and other safety measures did not \uline{\hspace*{3em}} between the groups, after allowing for missing data.
    \end{enumerate}
  \end{enumerate}
\end{problemset}

\chapter{语步四}\label{chapter6}

语步4（M4）是医学学术论文摘要的最后一步，在这里可以强调和解释关键性的结果，指出研究的意义并给出建议。也可以对未来的工作或后续研究提出建议。在以假设为导向的研究中，也需要在这一语步中明确说明假设是否得到了结果的支持。

\section{M4中的语阶}

M4涉及四个语阶，顺序相对固定。

4S1 重申关键性的结果

4S2 指出局限性

4S3 说明结果的重要性

4S4 预测未来的研究

\begin{sample}[label={myautocounter}]{\heiti}
  
  \textbf{INTERPRETATION} 
  
  In patients with chronic migraine, erenumab 70 mg and 140 mg reduced the number of monthly migraine days with a safety profile similar to placebo, providing evidence that erenumab could be a potential therapy for migraine prevention. Further research is needed to understand longterm efficacy and safety of erenumab, and the applicability of this study to real-world settings.


  \begin{flushright}
    ---Safety and efficacy of erenumab for preventive treatment of chronic migraine: a randomised, double-blind, placebo-controlled phase 2 trial. \emph{The Lancet Neurology (2017)}
  \end{flushright}

  \tcblower

  \noindent \textbf{STEP IDENTIFICATION}

  \vspace*{10pt}
  {\small\noindent
  \begin{tblr}{colspec={X[1,c]X[5,l]},cell{even}{1,2} = {azure9}}
    \toprule
    \textbf{Step} & \textbf{Sample} \\ 
    \midrule
  
    4S1 & In patients with chronic migraine, erenumab 70 mg and 140 mg reduced the number of monthly migraine days with a safety profile similar to placebo, \\
    4S3 & providing evidence that erenumab could be a potential therapy for migraine prevention. \\
    4S4 & Further research is needed to understand longterm efficacy and safety of erenumab, and the applicability of this study to real-world settings. \\
    
    \bottomrule
  \end{tblr}
  }

\end{sample}

\begin{sample}[label={myautocounter}]{\heiti}
  
  \textbf{CONCLUSION} 
  
  Including information from \ce{^{18}F}-fluciclovine PET in the treatment-planning process led to significant differences in the defined target volume, with higher doses to the penile bulb but no significant differences in rectal or bladder dose or in acute genitourinary or gastrointestinal toxicity. Longer follow-up is needed to determine the impact of 18F-fluciclovine PET on cancer control and late toxicity endpoints.

  \begin{flushright}
    ---Impact of \ce{^{18}F}-fluciclovine PET on Target Volume Definition for Postprostatectomy Salvage Radiotherapy: Initial Findings from a Randomized Trial. \emph{The Journal of Nuclear Medicine (2017)}
  \end{flushright}

  \tcblower

  \noindent \textbf{STEP IDENTIFICATION}

  \vspace*{10pt}
  {\small\noindent
  \begin{tblr}{colspec={X[1,c]X[5,l]},cell{even}{1,2} = {azure9}}
    \toprule
    \textbf{Step} & \textbf{Sample} \\ 
    \midrule
  
    4S3 & CONCLUSION: Including information from \ce{^{18}F}-fluciclovine PET in the treatment-planning process led to significant differences in the defined target volume, with higher doses to the penile bulb but no significant differences in rectal or bladder dose or in acute genitourinary or gastrointestinal toxicity.\\
    4S4 & Longer follow-up is needed to determine the impact of 18F-fluciclovine PET on cancer control and late toxicity endpoints.\\
    
    \bottomrule
  \end{tblr}
  }
  
\end{sample}

\begin{sample}[label={myautocounter}]{\heiti}
  
  \textbf{CONCLUSION }
  
  Our approach results in a response rate of 40\% or more, with acceptable toxicity. \ce{^{18}F}-FMISO uptake in NSCLC patients is strongly associated with poor prognosis features that could not be reversed by radiotherapy doses up to 86 Gy.


  \begin{flushright}
    ---Phase II study of a radiotherapy total dose increase in hypoxic lesions identified by F-miso PET/CT in patients with non-small cell lung carcinoma. \emph{The Journal of Nuclear Medicine (2017)}
  \end{flushright}

  \tcblower

  \noindent \textbf{STEP IDENTIFICATION}

  \vspace*{10pt}
  {\small\noindent
  \begin{tblr}{colspec={X[1,c]X[5,l]},cell{even}{1,2} = {azure9}}
    \toprule
    \textbf{Step} & \textbf{Sample} \\ 
    \midrule
  
    4S1 & CONCLUSION Our approach results in a response rate of 40\% or more, with acceptable toxicity.\\
    4S3 & \ce{^{18}F}-FMISO uptake in NSCLC patients is strongly associated with poor prognosis features that could not be reversed by radiotherapy doses up to 86 Gy.\\
    
    \bottomrule
  \end{tblr}
  }

\end{sample}

\begin{sample}[label={myautocounter}]{\heiti}
  
  \textbf{CONCLUSION}
  
  These findings support the hypothesis and endorse ARIP as a safer APD for alleviating behavioral disturbances after TBI.


  \begin{flushright}
    ---Relative to typical antipsychotic drugs, aripiprazole is a safer alternative for alleviating behavioral disturbances after experimental brain trauma. \emph{Neurorehabilitation and Neural Repair (2016)}
  \end{flushright}

  \tcblower

  \noindent \textbf{STEP IDENTIFICATION}

  \vspace*{10pt}
  {\small\noindent
  \begin{tblr}{colspec={X[1,c]X[5,l]},cell{even}{1,2} = {azure9}}
    \toprule
    \textbf{Step} & \textbf{Sample} \\ 
    \midrule
  
    4S3 & CONCLUSION These findings support the hypothesis and endorse ARIP as a safer APD for alleviating behavioral disturbances after TBI.\\

    \bottomrule
  \end{tblr}
  }
  
\end{sample}

从上面列出的样本可以看出，在许多情况下，4S3是M4的常规语阶，而其他语阶是可选的。通过分析研究的主要发现，M4指出了研究的意义或影响，并得出了结论。

\section{每个语阶的语言特征}
  \subsection{语阶1(4S1) 重申关键性的结果}
    \subsubsection{语阶分析}

    通常情况下，正如M3中罗列的那样，研究取得的结果往往不止一个，但在这一语阶中只总结和重申最关键的结果。

    \begin{eg}{}
      \par
      \vspace*{10pt}
      {\small\noindent
      \begin{tblr}{colspec={X[1,c]X[5,l]},cell{even}{1,2} = {azure9},cell{odd}{1,2} = {white}}
        \toprule
        \textbf{Step} & \textbf{Sample} \\ 
        \midrule
        
        3S2 & {RESULTS:\\ \uline{N-back task accuracy (N2 and N3) improved} after real-rTMS (and not after sham-rTMS) compared with baseline ($p=0.029$ and $p=0.015$, respectively), only in patients. \uline{At baseline, patients with MS, compared with HCs, showed higher task-related frontal activation (left DLPFC, N2>NO), which disappeared after real-rTMS.} Task-related (N1 $>$ NO) functional connectivity between the right DLPFC and the right caudate nucleus and bilateral (para) cingulate gyrus increased in patients after real-rTMS when compared with sham stimulation.}\\
        4S1 & {CONCLUSIONS:\\ In patients with MS, N-back accuracy improved while frontal hyperactivation (seen at baseline relative to HCs) disappeared after real-rTMS.}\\
          
        \bottomrule
      \end{tblr}
      }  
    \end{eg}

    \begin{eg}{}
      \par
       \vspace*{10pt}
       {\small\noindent
       \begin{tblr}{colspec={X[1,c]X[5,l]},cell{even}{1,2} = {azure9},cell{odd}{1,2} = {white}}
         \toprule
         \textbf{Step} & \textbf{Sample} \\ 
         \midrule
         
         3S1 & RESULTS: Seventy-nine patients were pre-included, 54 were included, and 34 were F-miso positive, 24 of whom received escalated doses of up to 86 Gy.\\
         3S2 & \uline{The response rate at 3 months was 31/54} (57\%, 95\% confidence interval [43\%--71\%]). DFS and OS at 1 year were 0.86[0.77--0.96] and 0.63[0.49--0.74], respectively. DFS was longer in the F-miso negative patients ($P=0.004$). The RT dose was not associated with DFS when adjusting for the F-miso status. \uline{One toxic death (66 Gy) and 1 case of grade 4 pneumonitis ($>66$ Gy) were reported.}\\
         4S1 & Our approach results in a response rate $>40\%$ with acceptable toxicity.\\
           
         \bottomrule
       \end{tblr}
       }   
    \end{eg}

    \subsubsection{语言实现}

    同3S2一样，这里通常采用一般过去时来呈现研究结果，表明这在局部是正确的。

    \begin{eg}{}
      Conclusions and Relevance Individuals with ADHD \uline{had} higher scores on the higher education entrance tests during periods they were taking ADHD medication vs nonmedicated periods.
    \end{eg}

    \begin{eg}{}
      Conclusion combinations of eye diseases \uline{were} frequent at old age. 
    \end{eg}

    \subsubsection{词组}

    \begin{enumerate}
      \item was associated with a
      \begin{eg}{}
        Depression severity \uline{was associated with a} decrease in measures of body composition in older adolescents over a mean of 1. 5 years, whereas SSRI treatment was positively associated with these outcomes, with differential effects across treatment groups.
      \end{eg}
      \item was not associated with
      \begin{eg}{}
        In this large population-based cohort study, HG \uline{was not associated with} an increased risk of long-term all-cause mortality.
      \end{eg}
    \end{enumerate}

  \subsection{语阶2(4S2) 指出局限性}
    \subsubsection{语阶分析}

    虽然这一语阶经常被列在研究文章中以增加可信度，但在摘要中却很少见到。主要的局限性可能包括研究的范围；参加者的年龄、种族或性别；未知因素，例如现有的医疗条件；还有研究者的偏见。

    \begin{eg}{}
      Limitations Observational studies do not establish cause and effect.
    \end{eg}

    \subsubsection{语言实现}

    表示对比或否定意义的单词，如“not”或“lack”，通常被视为这一语阶的信号词。

    \begin{eg}{}
      In this trial, we provide results for 3 years of treatment, with the limitation that withdrawn individuals were \uline{not} followed up after discontinuation.   
    \end{eg}

    \begin{eg}{}
      A major limitation of this study is the \uline{lack} of histopathologic proof in most patients.     
    \end{eg}

  \subsection{语阶3(4S3) 说明结果的重要性}
    \subsubsection{语阶分析}

    这是M4中必不可少的一步，向读者提供研究可能对未来或相关领域的影响，对研究结果的可能应用提出建议，说明初始假设是否被结果证明或否定，将结果与以往的研究结果进行比较，或表明本研究取得的突破。

    \begin{eg}{}
      These findings suggest that ADHD medications may help ameliorate educationally relevant outcomes in individuals with ADHD.
    \end{eg}

    \subsubsection{语言实现}

    一般过去时和一般现在时都可以在这一语阶中使用，后者使用得更频繁，由此显得结论似乎更具普适性。

    在这一步中经常看到的“suggest”一词表明作者正在解释研究结果。


    \begin{eg}{}
      CONCLUSIONS: These results supported the hypothesis and \uline{suggest} that DSW is an effective exercise intervention for elderly obese women to improve their health and fitness.
    \end{eg}
    
    在这一步中经常使用“this is The first”来表达对研究的肯定。

    \begin{eg}{}
      To our knowledge, \uline{this is the first} observational quality improvement initiative in otolaryngology to study the operative flow of a specific procedure and provide insight into areas of patient risk and opportunities for improvement in efficiency.  
    \end{eg}
    
    此外，情态动词“may”在4S3被广泛使用，这在某种程度上可能会削弱作者主张的权威性，为读者提供更多的讨论空间。反过来，这将保护作者得出的结论不受质疑或否定。

    \begin{eg}{}
      These results \uline{may} inform future childhood cancer treatment protocols and SMN surveillance guidelines for CCSs.
    \end{eg}

    \begin{eg}{}
      CONCLUSIONS: PCATS serves as a useful, and valid, predictor of ASA PS classification. Thus, it \uline{may} also serve as a tool to triage patients to an appropriate venue for preoperative assessment that can be utilized by nonclinical schedulers. Using a simple tool such as PCATS may help streamline the presurgical patient experience and improve clinic staff utilization.
    \end{eg}

    \begin{task}{\heiti 语料库任务}
      在这语步中还使用了哪些情态动词？
    \end{task}

    \subsubsection{词组}

    \begin{enumerate}
      \item the results of this study
      \begin{eg}{}
        \uline{The results of this study} confirm previous research findings that there is gadolinium deposition in wider distribution throughout the brain.
      \end{eg}

      \item this is the first
      \begin{eg}{}
        To the best of our knowledge, \uline{this is the first} study that investigates the effect of laser acupuncture in SAIS.  
      \end{eg}

      \item these findings suggest that
      \begin{eg}{}
        \uline{These findings suggest that} adaptive WMT may be an effective adjunctive therapy for WM deficits in HIV participants.  
      \end{eg}

      \item our/these results suggest that
      \begin{eg}{}
        \uline{Our results suggest that} parental pesticide exposure before or during pregnancy may play a role in the development of childhood retinoblastoma. 
      \end{eg}

      \item can be used to
      \begin{eg}{}
        Exhaled breath \uline{can be used to} detect recent cannabis exposure.
      \end{eg}

      \item for the treatment of
      \begin{eg}{}
        Combined RF/PDL technology is promising \uline{for the treatment of} recalcitrant PWS. 
      \end{eg}

      \item has the potential to
      \begin{eg}{}
        Reducing indoor tanning \uline{has the potential to} reduce melanoma incidence, mortality, and treatment costs. 
      \end{eg}

    \end{enumerate}

  \subsection{语阶4(4S4) 预测未来的研究}
    \subsubsection{语阶分析}

    这语阶预测了未来的研究方向，指出了有待解决的问题。作者可能会试图向读者解释还可以做些什么，或者下一步要采取什么行动。

    \subsubsection{语言实现}

    义务情态的动词，用以表达一定程度的义务性和必要性，如“should”或“need”，在这一语阶中经常被使用，以强调作者希望在未来进行进一步的研究。

    \begin{eg}{}
      Further research \uline{is needed} to assess the clinical importance of these differences and measure longer-term associations.  
    \end{eg}

    \begin{eg}{}
      Larger trials \uline{should} further elucidate the effect of changes in epidermal AQP-3 expression in development of vitiligo, and that might pave the road for discovering new therapeutic modalities for the disease. 
    \end{eg}

    \subsubsection{词组}

    \begin{enumerate}
      \item further research is needed to
      \begin{eg}{}
        \uline{Further research is needed to} assess generalizability and cost-effectiveness of this intervention and to understand which components may have contributed most to the outcome. 
      \end{eg}

      \item studies are needed to
      \begin{eg}{}
        Further prospective \uline{studies are needed to} establish which patients benefit from imaging.  
      \end{eg}

    \end{enumerate}

\section{材料阅读}

这一小节提供了2段M4材料，并对其中语阶的语言特点进行了详细分析。

\begin{sample}[label={myautocounter}]{\heiti}

  \textbf{CONCLUSIONS} 
  
  In patients with MS, N-back accuracy improved while frontal hyperactivation (seen at baseline relative to HCs) disappeared after real-rTMS. Together with the changes in functional connectivity after real-rTMS in patients, these findings may represent an rTMS-induced change in network efficiency in patients with MS, shifting patients' brain function towards the healthy situation. This implicates a potentially relevant role for rTMS in cognitive rehabilitation in MS.

  
  \begin{flushright}
    ---rTMS affects working memory performance, brain activation and functional connectivity in patients with multiple sclerosis. \emph{Journal of Neurology, Neurosurgery, and Psychiatry (2017)}
  \end{flushright}

  \tcblower

  \noindent \textbf{STEP IDENTIFICATION}

  \vspace*{10pt}
  {\small\noindent
  \begin{tblr}{colspec={X[1,c]X[5,l]},cell{even}{1,2} = {azure9}}
    \toprule
    \textbf{Step} & \textbf{Sample} \\ 
    \midrule
  
    4S1 & CONCLUSIONS: In patients with MS, N-back accuracy improved while frontal hyperactivation (seen at baseline relative to HCs) disappeared after real-rTMS. \\
    4S3 & Together with the changes in functional connectivity after real-rTMS in patients, these findings may represent an rTMS-induced change in network efficiency in patients with MS, shifting patients' brain function towards the healthy situation. This implicates a potentially relevant role for rTMS in cognitive rehabilitation in MS. \\
    
    \bottomrule
  \end{tblr}
  }

  \noindent \textbf{ANALYSIS}

  The sample above first restates the pivotal findings of the research in simple past tense to indicate that the findings are locally true and then points out its significance in simple present tense to indicate that the significance is universally true. In 4S3, the expression ``$\dots$ may represent $\dots$'' reveals the possible effect of real-rTMS which ``implicates a potentially relevant role'' for it. The word ``implicate'' means to show or to claim.
  
\end{sample}

\begin{sample}[label={myautocounter}]{\heiti}
  
  \textbf{CONCLUSIONS} 
  
  At 12 months, both FES and AFOs continue to demonstrate equivalent gains in gait speed. Results suggest that long-term FES use may lead to additional improvements in walking endurance and functional ambulation; further research is needed to confirm these findings.


  \begin{flushright}
    ---Long-Term Follow-up to a Randomized Controlled Trial Comparing Peroneal Nerve Functional Electrical Stimulation to an Ankle Foot Orthosis for Patients With Chronic Stroke. \emph{Neurorehabilitation \& Neural Repair (2015)}
  \end{flushright}

  \tcblower

  \noindent \textbf{STEP IDENTIFICATION}

  \vspace*{10pt}
  {\small\noindent
  \begin{tblr}{colspec={X[1,c]X[5,l]},cell{even}{1,2} = {azure9}}
    \toprule
    \textbf{Step} & \textbf{Sample} \\ 
    \midrule
  
    4S1 & CONCLUSIONS: At 12 months, both FES and AFOs continue to demonstrate equivalent gains in gait speed. \\
    4S3 & Results suggest that long-term FES use may lead to additional improvements in walking endurance and functional ambulation; \\
    4S4 & further research is needed to confirm these findings. \\
    
    \bottomrule
  \end{tblr}
  }

  \noindent \textbf{ANALYSIS}

  The first sentence restates the most important result of the research. Different from other examples, simple present tense is used in this case to indicate the ongoing situation. The first half of the second sentence presents the significance of the results in simple present tense, with the words ``suggest'' and ``may'', which are very often used in M4 to interpret the findings. At the end of the sample, prediction of the future studies is made with the use of the chunk ``further research is needed to $\dots$''.
  
\end{sample}

\section{词汇表}

{\small
\begin{longtblr}[
    caption = {第六章词汇表},
    label = {tab:Glossary of Chapter 6},
    % note{a} = {英文论文中指代当前文献中的差距、问题或缺陷。即现有研究尚未解决的部分。},
]{
    width = \textwidth,
    colspec = {X[1,l,h]  X[1,l,h]  X[3,l,h]},
    rowhead = 1, rowfoot = 0, % 每个分页里表头表尾的数量
    % row{odd} = {blue8}, 
    row{even} = {azure9},
}
    
\toprule
\textbf{WORDS} & \textbf{MEANING} & \textbf{MEANING or EXAMPLE}\\
\midrule

{\textbf{bias}\\/\textipa{"baI@s}/} & \emph{n.} 偏倚 & mental tendency or inclination, esp an irrational preference or prejudice \\
{\textbf{deontic modality}\\/\textipa{di:"6ntIkm@U"d\ae lItI}/} & \emph{n.} 义务情态；道义情态；性情态 & Deontic modality refers to the instances when the speaker orders, promises or places an obligation to someone. \\

\bottomrule

\end{longtblr}
}

\begin{problemset}
  \item \textbf{Find out how many steps are involved in the following examples.}
  
  \hspace*{2em}4S1 Reiterating pivotal results

  \hspace*{2em}4S2 Indicating limitations 
  
  \hspace*{2em}4S3 Stating the significance of the results

  \hspace*{2em}4S4 Predicting future studies

  \begin{enumerate}
    \item Conclusions: The Short model performed well during continuous infusion up to 545 min. This model might be preferable for target-controlled infusion for long-duration anaesthesia in children.
    \item Conclusions: Using an entrustment scoring system, where supervisors score trainees on the level of supervision required, mini-CEX scores demonstrated moderate reliability within a feasible number of assessments, and evidence of validity. When scores were adjusted against an expected standard, underperforming trainees could be identified, and reliability much improved. Taken together with other evidence on trainee ability, the mini-CEX is of sufficient reliability for inclusion in high stakes decisions on trainee progression towards independent specialist practice.
    \item Conclusions: This is the first study, to our knowledge, that demonstrates that patients with advanced laryngeal cancer with 1 or more comorbidities are more likely to receive surgery than chemoradiation compared with patients without any comorbidity, independent of numerous clinical and nonclinical variables among a large national cohort. A limitation of this study is the use of comorbidity data from the National Cancer Database, which gathers its information from hospital discharge face sheets. We recognize that the National Cancer Database may be an imperfect system for the collection of comorbidity data and encourage discussion on different methods to improve the system, including incorporating comorbidity data from the Surveillance, Epidemiology, and End Results Medicare Database and medical chart-based comorbidity data collection by cancer registrars.
    \item Interpretation: Home-based CIMT can enhance the perceived use of the stroke-affected arm in daily activities more effectively than conventional therapy, but was not superior with respect to motor function. Further research is needed to confirm whether home CIMT leads to clinically significant improvements and if so to identify patients that are most likely to benefit.
  \end{enumerate}

  \item \textbf{Rearrange the sentences in correct order.}
  \begin{enumerate}
    \item \uline{\hspace*{3em}}
    
    Conclusions: 
    \begin{itemize}
      \item[A.] These findings suggest vitamin D supplementation may have a therapeutic role in the treatment of MS.
      \item[B.] However, there is uncertainty with regards to the most appropriate dose, with higher dose potentially being associated with worse outcomes.
      \item[C.] There remains the need for a well performed randomised, dose-comparison, placebo-controlled trial of vitamin D in MS.
    \end{itemize}

    \item \uline{\hspace*{3em}}
    
    Conclusions: 
    \begin{itemize}
      \item[A.] Our study was not focused on outcomes or efficacy therefore further research is needed on whether tele-epilepsy provides equal outcomes in care.
      \item[B.] Despite significantly higher face-to-face patient satisfaction, tele-epilepsy patients reported high levels of satisfaction with the tele-epilepsy model of care.
      \item[C.] As such, tele-epilepsy is a viable model for the delivery of epileptic care within regional settings.
    \end{itemize}

    \item \uline{\hspace*{3em}}
    
    Conclusions: 
    \begin{itemize}
      \item[A.] The low seroconversion rates seen after the first year of exposure suggests that less frequent JCV testing of natalizumab treated patients may be safe.
      \item[B.] In contrast to recently published data from colleagues in France and Germany, in our longitudinal cohort study rates of durable positive JCV seroconversion correlate closely with expected background rates of JCV seroconversion in the general population.
      \item[C.] In addition, JCV seroconversion and positive serostatus is not increased in patients with prior exposure to immunosuppressive agents.
    \end{itemize}

    \item \uline{\hspace*{3em}}
    
    Conclusions: 
    \begin{itemize}
      \item[A.] We suggest that controlling the inflammation, even using higher doses of systemic and topical corticosteroids, is of importance in preventing ocular complications, such as cataract.
      \item[B.] In this study, we found that development of cataract is common among pediatric eyes with uveitis and is most strongly related to the extent of inflammation recurrences and ocular complications.
    \end{itemize}

  \end{enumerate}
  
  \item \textbf{Fill in the blanks with appropriate words or the correct form of the verbs given.}
  \begin{enumerate}
    \item Partial least squares regression showed that alcohol-induced changes in mood/drug effects \uline{\hspace*{3em}}(associate) with changes in thalamic FCD in both groups.
    \item A major limitation of this study is the \uline{\hspace*{3em}} of histopathologic proof in most patients.
    \item Our findings also \uline{\hspace*{3em}}(suggest) that the association between social engagement and mental health varies by type of engagement and initial depression level.
    \item The \uline{\hspace*{3em}} of this study indicate that activation of latent HPV infections may contribute to the increased risk of HPV-related (pre) malignant lesions in female RTRs.
    \item These study findings indicate that masitinib is an effective and well tolerated agent \uline{\hspace*{3em}} the treatment of severely symptomatic indolent or smouldering systemic mastocytosis.
    \item Further research is \uline{\hspace*{3em}} to assess the clinical importance of these differences and measure longer-term associations.
  \end{enumerate}
\end{problemset}

\restoregeometry

\chapter*{课后练习答案}

\paragraph*{Chapter 1}\par
\begin{enumerate}
  \item \textbf{Identify whether the following abstracts are structured or unstructured and tell the reasons.}
  
  \textbf{Abstract 1}
  
  \hspace*{2em}structured abstract 
  
  \textbf{Abstract 2}

  \hspace*{2em}unstructured abstract 
  
  \textbf{Abstract 3}

  \hspace*{2em}structured abstract
\end{enumerate}

\paragraph*{Chapter 2}\par
\begin{enumerate}
  \item \textbf{Identify the four moves in the abstracts.}
  
  \textbf{Abstract 1}

  \hspace*{2em}M1: Objective 

  \hspace*{2em}M2: Design, Setting, and Methods 

  \hspace*{2em}M3: Results M4: Conclusion 
  
  \textbf{Abstract 2}

  \hspace*{2em}M1: Skeletal muscle... walking speed.

  \hspace*{2em}M2: Muscle mitochondrial.. respectively.

  \hspace*{2em}M3: In multivariate linear regression analyses,... four walking tasks.
  
  \hspace*{2em}M4: This is the first.. through this mechanism.
  
  \textbf{Abstract 3}
  
  \hspace*{2em}M1: Propacetamol... in ED patients.
  
  \hspace*{2em}M2: This was a... vasopressor administration.
  
  \hspace*{2em}M3: Postinfusion SBP and DBP... hemodynamic change.
  
  \hspace*{2em}M4: Administration of propacetamol.. propacetamol.
  
  \item \textbf{Reorder the four moves in the abstracts.}

  \textbf{Abstract 1}: 2-1-4-3
  
  \textbf{Abstract 2}: 3-1-2-4

  \textbf{Abstract 3}: 4-2-1-3

\end{enumerate}

\paragraph*{Chapter 3}\par
\begin{enumerate}
  \item \textbf{Find out how many steps are involved in the following examples.}
  \begin{enumerate}
    \item 1S1+1S2
    \item 1S2+1S3a
    \item 1S1+1S3b
    \item 1S1+1S2+1S3a
  \end{enumerate}

  \item \textbf{Rearrange the sentences in correct order.}
  \begin{enumerate}
    \item ACB
    \item ACB
    \item CBA
    \item CBA
  \end{enumerate}

  \item \textbf{Fill in the blanks with appropriate words or the correct form of the verbs given.}
  \begin{enumerate}
    \item tested
    \item would improve
    \item To investigate
    \item aimed
    \item report
    \item have proposed
    \item sought
    \item however
    \item remain
    \item unknown/unclear/not known
    \item have found; have compared; aimed to
    \item is; have shown; has not (yet) been; compared
    \item are; have been (previously) shown; are;To compare
  \end{enumerate}

  \item \textbf{Discussion：Read and the following examples and discuss with your partners about the reasons why modal auxiliary verbs and words indicating likelihood are used?}
  
  略
\end{enumerate}

\paragraph*{Chapter 4}\par
\begin{enumerate}
  \item \textbf{With the ``Steps in M2'' offered, try to identify the steps in the following 2 examples of M2 and write down the corresponding numbers in the blanks.}
  
  \textbf{Example 1}\par
  \vspace*{10pt}
  {\small\noindent
  \begin{tblr}{colspec={X[1,c]X[5,l]},cell{even}{1,2} = {azure9}}
    \toprule
    \textbf{Step} & \textbf{Sample} \\ 
    \midrule
    
    2S2 & DESIGN: Retrospective case-control study.\\
    2S3 & SETTING: Nijmegen, the Netherlands.\\
    2S3 & PATIENTS: Thirty consecutive patients with SSD.\\
    2S4 & INTERVENTIONS: Patients received a trial with a BCD headband as part of the regular workup for SSD. The patients were divided into 2 groups according to their decision to opt for a BCD (BCD+) or not (BCD-). \\
    2S5 & MAIN OUTCOME MEASURES: Patients completed a questionnaire on satisfaction with the BCD headband, patient-and BCD-related factors, and benefit in listening situations. \\
      
    \bottomrule
  \end{tblr}
  }

  \textbf{Example 2}\par
  \vspace*{10pt}
  {\small\noindent
  \begin{tblr}{colspec={X[1,c]X[5,l]},cell{even}{1,2} = {azure9}}
    \toprule
    \textbf{Step} & \textbf{Sample} \\ 
    \midrule
    
    2S3 with 2S2 embedded & (2S3) Eighteen adults with first-ever chronic monohemispheric subcortical stroke participated in (2S2) this randomized, controlled, triple-blinded trial. \\
    2S4 & Intervention consisted of priming with real or sham iTBS to the ipsilesional primary motor cortex immediately before 45 minutes of upper limb physical therapy, daily for 10 days. \\
    2S5 & Changes in upper limb function (Action Research Arm Test [ARAT]), upper limb impairment (Fugl-Meyer Scale), and corticomotor excitability, were assessed before, during, and immediately, 1 month and 3 months after the intervention. Functional magnetic resonance images were acquired before and at one month after the intervention. \\
      
    \bottomrule
  \end{tblr}
  }

  \item \textbf{Fill in the blanks with the correct form of the verbs given. Notice the tense and voice.}
  \begin{enumerate}
    \item ~
    \begin{enumerate}
      \item included
      \item were
      \item completed
      \item were identified
      \item was applied
    \end{enumerate}
    
    \item ~
    \begin{enumerate}
      \item is/was
      \item were matched
      \item were performed
      \item to compare
      \item included
    \end{enumerate}
  \end{enumerate}

  \item \textbf{Error Correction}
  \begin{enumerate}
    \item randomized
    \item with
    \item were ($\wedge$)
    \item criteria
    \item screened
    \item as (/)
    \item median
    \item were
    \item an
    \item $\surd$
  \end{enumerate}
\end{enumerate}

\paragraph*{Chapter 5}\par
\begin{enumerate}
  \item \textbf{Fill the blanks with the correct forms of the given words. Notice the tense and voice.}
  \begin{enumerate}
    \item ~
    \begin{enumerate}
      \item did not differ
      \item received
      \item showed
      \item increased
      \item decreased
      \item was seen
    \end{enumerate}
    
    \item ~
    \begin{enumerate}
      \item was
      \item ate
      \item was halved
      \item was associated
      \item were not observed
    \end{enumerate}
  \end{enumerate}

  \item \textbf{``Significance'' and ``difference'' are two of the key words in M3, so we should be familiar with their different forms. Please fill in the blanks with the proper forms of the words given.}
  \begin{enumerate}
    \item ~
    \begin{enumerate}
      \item significantly
      \item significant
      \item significance
      \item significant
      \item significantly
      \item significant
    \end{enumerate}

    \item ~
    \begin{enumerate}
      \item different
      \item Differences
      \item differently
      \item differed
      \item different
      \item differed
      \item differently
      \item differ
    \end{enumerate}
  \end{enumerate}
\end{enumerate}

\paragraph*{Chapter 6}\par
\begin{enumerate}
  \item \textbf{Find out how many steps are involved in the following examples.}
  \begin{enumerate}
    \item 4S1+4S3
    \item 4S1+4S3
    \item 4S3+4S2+424
    \item 4S3+4S4
  \end{enumerate}

  \item \textbf{Rearrange the sentences in correct order.}
  \begin{enumerate}
    \item ABC
    \item BCA
    \item BCA
    \item BA
  \end{enumerate}

  \item \textbf{Fill in the blanks with appropriate words or the correct form of the verbs given.}
  \begin{enumerate}
    \item were associated
    \item lack
    \item suggest
    \item results/findings
    \item for
    \item needed
  \end{enumerate}
\end{enumerate}


\appendix

\chapter{语步中的四词及以上的词组}

{\small
\begin{longtblr}[
    caption = {语步中的四词及以上的词组},
    label = {tab:Four-word and Longer Lexical Bundles in Moves},
    remark{Note 1} = {structures (Str), functions (Fun), noun based bundles (N), preposition based bundles (P), verb based bundles (V), otherbundles (O), research based bundles (R), text based bundles (T), stance based bundles (S)},
    remark{Note 2} = {Lexical bundles that appear more than once are in bold.}
]{
    width = \textwidth,
    colspec = {X[12,l,h]  X[7,c,h]  X[7,c,h] X[1,c,h] X[1,c,h]},
    rowhead = 1, rowfoot = 0, % 每个分页里表头表尾的数量
    % row{odd} = {blue8}, 
    row{even} = {azure9},
}
    
\toprule
\textbf{Bundles} & \textbf{Tokens} & \textbf{Dispersion} & \textbf{Str} & \textbf{Fun} \\ 
\midrule
\SetCell[c=5]{c,white}\textbf{M1 Bundles (60 types)} & & \\ 
\midrule
the aim of this study was to evaluate the & 20 & 20 & V & R \\
the aim of this sudy was to investigate & 31 & 31 & V & R \\
the aim of this study was to assess & 23 & 23 & V & R \\
the aim of this study was to & 160(86   only this bundle) & 160(86   only this bundle) & V & R \\
the purpose of this study was to & 106 & 106 & V & R \\
the objective of this study was to & 46 & 46 & V & R \\
of this study was to detemine the & 26 & 26 & V & R \\
of this study was to investigate & 54(23   only this bundle) & 54(23   only this bundle) & V & R \\
of this study was to determine & 48(22   only this bundle) & 48(22   only this bundle) & V & R \\
of this study was to evaluate & 44(22   only this bundle) & 44(22   only this bundle) & V & R \\
of this study was to compare & 31 & 31 & V & R \\
of this study was to examine & 25 & 25 & V & R \\
aim of the study was to & 20 & 20 & V & R \\
of this study was to & 352(49   only this bundle) & 352(49   only this bundle) & V & R \\
study was to evaluate the & 47 & 47 & V & R \\
study was to invstigate the & 44 & 44 & V & R \\
of the study was to & 41(21   only this bundle) & 41(21   only this bundle) & V & R \\
the efficacy and safety of & 37 & 37 & N & R \\
study was to asses the & 36 & 36 & V & R \\
。 the safety and efficacy of & 32 & 32 & N & R \\
study was to determine the & 29 & 29 & V & R \\
of this study were to & 28 & 28 & V & R \\
we aimed to assess the & 27 & 27 & V & R \\
of this study is to & 26 & 26 & V & R \\
little is known about the & 23 & 23 & V & T \\
study was to compare the & 22 & 22 & V & R \\
to detemine the effect of & 21 & 21 & O & R \\
to examine the association between & 20 & 20 & O & R \\
the aim of this & 205(45   only this bundle) & 205(45   only this bundle) & N & R \\
aim of this study & 189(29   only this bundle) & 189(29   only this bundle) & N & R \\
the purpose of this & 127(21   only this bunde) & 127(21   only this bundle) & N & R \\
study was to invesigate & 71(27   only this bundle) & 71(27   only this bundle) & V & R \\
study was to evaluate & 67(20   only this bundle) & 67(20   only this bundle) & V & R \\
this sudy aimed to & 67 & 67 & V & R \\
study was to assess & 53(30   only this bundle) & 53(30   only this bundle) & V & R \\
we aimed to assess & 53(26   only this bundle) & 53(26   only this bundle) & V & R \\
little is known about & 49(26   only this bundle) & 49(26   only this bundle) & V & T \\
of the study was & 41(20   oly this bundle) & 41(20   only this bundle) & V & R \\
the study was to & 41(20   only this bundle) & 41(20   only this bundle) & V & R \\
study was to compare & 40 & 40 & V & R \\
the aim of the & 40 & 40 & N & R \\
for the treatment of & 34 & 34 & P & R \\
in patients with advanced & 28 & 26 & P & R \\
in the emergency department & 27 & 27 & P & R \\
our objective was to & 27 & 27 & V & R \\
in the treatment of & 26 & 26 & P & R \\
has been associated with & 25 & 25 & V & T \\
has been shown to & 25 & 25 & V & T \\
our aim was to & 25 & 25 & V & R \\
we sought to determine & 24 & 24 & V & R \\
aimed to determine the & 23 & 23 & V & R \\
aimed to investigate the & 23 & 23 & V & R \\
was to determine whether & 23 & 23 & V & R \\
aimed to evaluate the & 21 & 21 & V & R \\
we aimed to evaluate & 21 & 21 & V & R \\
is one of the & 20 & 20 & V & T \\
study aimed to determine & 20 & 20 & V & R \\
tested the hypothesis that & 20 & 20 & V & R \\
the objective of the & 20 & 20 & N & R \\
was to examine the & 20 & 20 & V & R \\
\midrule
\SetCell[c=5]{c,white}\textbf{M1 Bundles (50 types)} & & \\ 
\midrule
all   patiens who recived at least one dose of & 21 & 21 & O & R \\
patients were randomly assigned to receive & 26 & 26 & V & R \\
the primary end point was the & 20 & 20 & V & R \\
this trial is regisered with & 76 & 76 & V & R \\
\textbf{were randomly assigned to receive} & 65(39   only this bundle) & 65(39   only this bundle) & V & R \\
the primary end point was & 59(39   only this bundle) & 58(38   only this bundle) & V & R \\
patients were randomly assigned to & 51(25   only this bundle) & 50(24   only this bundle) & V & R \\
the primary outcome was the & 44 & 44 & V & R \\
this study is regisered with & 44 & 44 & V & R \\
the primary endpoint was the & 33 & 33 & V & R \\
the primary outcome measure was & 31 & 31 & V & R \\
a retrospective cohort study of & 26 & 26 & N & R \\
was the proportion of patients & 24 & 23 & V & R \\
logisic regression was used to & 23 & 23 & V & R \\
study was approved by the & 22 & 22 & V & R \\
the primary oucome was & 187(143   only this bundle) & 186(142   only this bundle) & V & R \\
were randomly assigned to & 145(80   only this bundle) & 141(77   only this bundle) & V & R \\
the primary endpoint was & 117(84   only this bundle) & 115(82   only this bundle) & V & R \\
a retropective cohort sudy & 54(28   only this bundie) & 54(28   only this bundle) & N & R \\
were included in the & 46 & 44 & V & R \\
at the time of & 44 & 43 & P & R \\
on the basis of & 42 & 41 & P & T \\
were randomized to receive & 41 & 37 & V & R \\
were masked to treatment & 39 & 39 & V & R \\
models were used to & 36 & 35 & V & R \\
analysis was used to & 34 & 32 & V & R \\
we conducted a retrospecive & 34 & 34 & V & R \\
was used to identify & 30 & 30 & V & R \\
were included in this & 28 & 28 & V & R \\
at the end of & 27 & 25 & P & R \\
at a dose of & 26 & 22 & P & R \\
was assessed using the & 26 & 25 & V & R \\
years or older with & 26 & 26 & O & R \\
patients were randomized to & 25 & 23 & V & R \\
was used to assess & 25 & 24 & V & R \\
were used to assess & 25 & 25 & V & R \\
we randomly assigned patients & 25 & 25 & V & R \\
masked to treatment allocation & 24 & 24 & V & R \\
participants were randomly assigned & 24 & 24 & V & R \\
analyses were performed to & 23 & 23 & V & R \\
outcome measure was the & 23 & 23 & V & R \\
were used to detemine & 23 & 23 & V & R \\
masked to treatment assignment & 22 & 21 & V & R \\
data were collected from & 21 & 21 & V & R \\
in all patients who & 21 & 20 & P & R \\
a cross sectional study & 20 & 20 & N & R \\
a secondary analysis of & 20 & 20 & N & R \\
regression analysis was used & 20 & 20 & V & R \\
secondary end points included & 20 & 20 & V & R \\
were used to evaluate & 20 & 20 & V & R \\
outcome measure was the & 23 & 23 & V & R \\
were used to detemine & 23 & 23 & V & R \\
masked to treatment asignment & 22 & 21 & V & R \\
daa were cllcted from & 21 & 21 & V & R \\
\midrule
\SetCell[c=5]{c,white}\textbf{M3 Bundles (38 types)} & & \\ 
\midrule
there was no significat dfference in & 25 & 23 & V & T \\
there were no significant differences in & 24 & 24 & V & T \\
there were no significant differences between & 22 & 21 & V & T \\
there was no difference in & 34 & 31 & V & T \\
adverse events were reported in & 33 & 28 & V & R \\
there were no differences in & 30 & 28 & V & T \\
paticnts were randomly asigned to & 30 & 25 & V & R \\
serious adverse events were reported & 29 & 28 & V & R \\
\textbf{were randomly assigned to receive} & 26 & 22 & V & R \\
were more likely to have & 25 & 21 & V & S \\
did not differ significantly between & 24 & 24 & V & T \\
patiets were included in the & 24 & 24 & V & R \\
did not differ between the & 24 & 22 & V & T \\
at a median follow-up of & 21 & 21 & P & R \\
were incuded in the study & 20 & 20 & V & R \\
in the placebo group & 162 & 76 & P & R \\
in the control group & 117 & 72 & P & R \\
\textbf{was associated with a/an} & 131 & 117 & V & T \\
\textbf{were included in the} & 102(78   only this bundle) & 94(70   only this bundle) & V & R \\
were more likely to & 82(57   only this bundle) & 68(47   only this bundle) & V & S \\
there was no significant & 76(51   only this bundle) & 70(47   only this bundle) & V & T \\
there were no significant & 68(44   only this bundle) & 67(43   only this bundle) & V & T \\
were randomly assigned to & 68(38   only this bundle) & 61(33   only this bundle) & V & R \\
a media fllow-up of & 55(34   only this bundle) & 55(34   only this bundle) & N & R \\
did not dfffe betwe & 52(28   only this bundle) & 47(25   only this buodle) & V & T \\
\textbf{was not associated with} & 52 & 47 & V & T \\
at the time of & 51 & 46 & P & R \\
was signifcantly associated with & 44 & 42 & V & T \\
was signifcantly higher in & 42 & 39 & V & T \\
area under the curve & 40 & 37 & N & R \\
during the study period & 40 & 37 & P & R \\
the most common grade & 40 & 35 & N & R \\
were less likely to & 40 & 35 & V & S \\
events were reported in & 37 & 30 & V & T \\
compared with the control & 33 & 26 & V & R \\
were not significantly dfferent & 33 & 32 & V & T \\
between the two gropgs & 31 & 27 & P & T \\
was not significantly dfferent & 31 & 29 & V & T \\
significantly higher in the & 30 & 30 & O & T \\
was found to be & 29 & 28 & V & T \\
was significantly lower in & 29 & 26 & V & T \\
incuded in the analys & 28 & 28 & V & R \\
an increased risk of & 28 & 24 & P & R \\
not significantly dfferent between & 28 & 26 & V & T \\
was the most common & 26 & 25 & V & R \\
were not associated with & 26 & 25 & V & T \\
group compared with the & 26 & 23 & O & R \\
was associated with increased & 25 & 23 & V & T \\
associated with an increased & 25 & 22 & V & T \\
no saisticallly signifcant difference & 24 & 23 & N & T \\
were significantly higbher in & 24 & 24 & V & T \\
sigpificant dffcence in the & 23 & 22 & N & T \\
no significant dffcence between & 23 & 21 & N & T \\
significant dfftences between the & 23 & 21 & N & T \\
was independenly associated with & 23 & 23 & V & T \\
patients were enrolled and & 22 & 22 & V & R \\
were associated with a & 21 & 20 & V & T \\
no sigificat dffrences were & 20 & 20 & V & T \\
\midrule
\SetCell[c=5]{c,white}\textbf{M4 Bundles (18 types)} & & \\ 
\midrule
further research is needed to & 20 & 20 & V & S \\
the results of this study & 20 & 20 & N & T \\
\textbf{was associated with a} & 44 & 42 & V & T \\
this is the first & 39 & 39 & V & S \\
\textbf{was not associated with} & 36 & 36 & V & T \\
these findings suggest that & 32 & 32 & V & T \\
studies are needed to & 28 & 28 & V & S \\
is associated with a & 27 & 26 & V & T \\
our results suggest that & 27 & 27 & V & T \\
these results suggest that & 27 & 27 & V & T \\
can be used to & 26 & 24 & V & S \\
an increased risk of & 22 & 22 & N & R \\
associated with an increased & 22 & 22 & V & T \\
for the treatment of & 20 & 20 & P & R \\
has the potential to & 20 & 20 & V & S \\
our findings suggest that & 20 & 20 & V & T \\ 

\bottomrule

\end{longtblr}
}

\chapter{语阶中的四词及以上的词组}

\begin{landscape}
  {\tiny
  \begin{longtblr}[
      caption = {语阶中的四词及以上的词组},
      label = {tab:Four-word and Longer Lexical Bundles in Steps},
  ]{
      width = \linewidth,
      colspec = {X[2,c,h]  X[8,l,h]  X[8,l,h]  X[6,l,h]},
      rowhead = 1, rowfoot = 0, % 每个分页里表头表尾的数量
      % row{odd} = {blue8}, 
      row{even} = {azure9},
  }
      
  \toprule
  \textbf{Move/ Step} & \textbf{Communicative functions} & \textbf{Specific bundles} & \textbf{Common bundles}\\
  \midrule
  M1 & Creating a research territory/ space &  & \\
  1S1 & Presenting current knowledge or relevant information established by previos studies & {has been asociated with \\
      has been shown to \\
      is one of the} & {for the teatment of \\
      in patients with advanccd \\
      in the emergency depatmen \\
      in the treatment of}\\
  1S2 & Establishing a niche/problem & {little is known about the \\
      little is known about} & \\
  1S3a & Indicating main purposs & {the aim of this study was to evaluate the \\
      the aim of this study was to invesigate \\
      the aim of this study was to asscss \\
      the aim of this study was to \\
      the pupose of this study was to \\
      the objecive of this study was to \\
      of this study was to determine the \\
      of this study was to invsigate \\
      of this sudy was to determine \\
      of this study was to evaluate \\
      of this sudy was to compare \\
      of this study was to examine \\
      aim of the study was to \\
      of this study was to \\
      study was to evaluate the \\
      study was to investigate the \\
      of the study was to \\
      the efficacy and safety of \\
      study was to assess the \\
      the safety and efficacy of \\
      study was to determine the \\
      of this study were to \\
      we aimed to asses the\\
      of this study is to\\
      study was to compare the\\
      to determine the effect of\\
      to examine the association between\\
      the aim of this\\
      aim of this study\\
      the purpose of this\\
      study was to investigate\\
      study was to evaluate\\
      this study aimed to\\
      study was to assess\\
      we aimed to assess\\
      of the study was\\
      the study was to\\
      study was to compare\\
      the aim of the\\
      our objective was to\\
      our aim was to\\
      we sought to determine\\
      aimed to determine the\\
      aimed to investigate the\\
      was to determine whether\\
      aimed to evaluate the\\
      we aimed to evaluate\\
      study aimed to determine\\
      the objective of the\\
      was to examine the} & {for the treatment of \\
      in patients with advanced \\
      in the emergency departnent \\
      in the treatment of}\\
  1S3b & Raising hypotheses & tested the hypothesis that & in patients with advanced\\
  M2 & Describing research process &  & \\
  2S1 & Reporting on medical ethics review & study was approved by the & \\
  2S2 & Explaining briefly research design & {a retrospective cohort study of\\
      a retrospective cohort study\\
      we conducted a retrospective\\
      a cross-sectional study\\
      a secondary analysis of} & \\
  2S3 & Describing subjects or data and their selection criteria & {were included in the\\
      were included in this\\
      years or older with\\
      data were collected from} & {at the time of\\
      on the basis of}\\
  2S4 & Describing experimental procedure, such as interventions, examinations,
    etc. & {patients were randomly assigned to receive\\
      were randomly assigned to receive\\
      patients were randomly assigned to\\
      were randomly assigned to\\
      were randomized to receive\\
      were masked to treatment\\
      patients were randomized to\\
      we randomly assigned patients\\
      masked to treatment allocation\\
      participants were randomly assigned\\
      masked to treatment assignment} & {all patients who received at least one dose of \\
      at the time of \\
      on the basis of \\
      at the end of\\
      at a dose of\\
      was assessed using the\\
      was used to assess\\
      were used to assess\\
      were used to evaluate}\\
  2S5 & Describing main outcomes and their measures & {the primary end point was the\\
      the primary end point was\\
      the primary outcome was the\\
      the primary endpoint was the\\
      the primary outcome measure was\\
      was the proportion of patients\\
      the primary outcome was\\
      the primary endpoint was\\
      outcome measure was the\\
      in all patients who\\
      secondary end points included} & {all patients who received at least one dose of \\
      at the time of \\
      on the basis of \\
      was used to identify\\
      at the end of \\
      at a dose of \\
      was assessed using the\\
      was used to assess \\
      were used to assess\\
      analysis was used to \\
      analyses were performed to \\
      were used to determine\\
      were used to evaluate}\\
  2S6 & Describing data analysis methods & {logistic regression was used to\\
      models were used to\\
      regression analysis was used} & {was used to identify\\
      analysis was used to\\
      analyses were performed to}\\
  2S7 & Reporting on registration information & {this trial is registered with\\
      this study is registered with} & were used to determine\\
  M3 & Summarizing results &  & \\
  3S1 & Explaining information on valid samples & {patients were randomly assigned to\\
      were randomly assigned to receive\\
      patients were included in the\\
      were included in the study\\
      were included in the\\
      included in the analysis\\
      patients were enrolled and} & {at a median follow-up of\\
      in the placebo group\\
      in the control group\\
      a median follow-up of during the study period}\\
  3S2 & Illustrating overall observation or main results & {there was no significant difference in\\
      there were no significant differences in\\
      there were no significant differences between\\
      there was no difference in\\
      adverse events were reported in\\
      there were no differences in\\
      serious adverse events were reported\\
      were more likely to have\\
      did not differ significantly between\\
      did not differ between the\\
      was associated with a/an\\
      were more likely to\\
      there was no significant\\
      there were no significant\\
      were randomly assigned to\\
      did not differ between\\
      was not associated with\\
      at the time of\\
      was significantly associated with\\
      was significantly higher in\\
      area under the curve\\
      the most common grade\\
      were less likely to\\
      events were reported in\\
      compared with the control\\
      were not significantly different between the two groups\\
      was not significantly different\\
      significantly higher in the\\
      was found to be\\
      was significantly lower in\\
      an increased risk of\\
      not significantly different between\\
      was the most common\\
      were not associated with\\
      group compared with the\\
      was associated with increased\\
      associated with an increased\\
      no statistically significant difference\\
      were significantly higher in\\
      significant difference in the\\
      no significant difference between\\
      significant differences between the\\
      was independently associated with\\
      were associated with a\\
      no significant differences were} & \\
  M4 & Drawing conclusions &  & \\
  4S1 & Reiterating pivotal results & {was associated with a\\
      was not associated with} & {is associated with a\\
      an increased risk of \\
      associated with an increased}\\
  4S2 & Indicating limitations &  & \\
  4S3 & Stating the significance of the results & {the results of this study\\
      this is the first\\
      these findings suggest that\\
      our results suggest that\\
      these results suggest that\\
      can be used to\\
      for the treatment of\\
      has the potential to\\
      our findings suggest that} & {is associated with a\\
      an increased risk of \\
      associated with an increased}\\
  4S4 & Predicting future studies & {further research is needed to\\
      studies are needed to} & \\
  
  \bottomrule
  
  \end{longtblr}

  % \begin{longtable} to \textwidth {XXXX}

  %   % \toprule
  %   % Move/ Step & Communicative functions & Specific bundles & Common bundles \\
  %   % \midrule 
  %   % \endfirsthead

  %   % Appear the table header at the top of every page
  %   \toprule
  %   Move/ Step & Communicative functions & Specific bundles & Common bundles \\
  %   \midrule 
  %   \endhead 
  
  %   % Appear \hline at the bottom of every page
  %   \bottomrule
  %   \endfoot 


  %   M1 & Creating a research territory/ space &  &  \\
  %   1S1 & Presenting current knowledge or relevant information established by previos studies & \begin{tabular}[c]{@{}l@{}}has been asociated with \\ has been shown to \\ is one of the\end{tabular} & \begin{tabular}[c]{@{}l@{}}for the teatment of \\ in patients with advanccd \\ in the emergency depatmen \\ in the treatment of\end{tabular} \\
  %   1S2 & Establishing a niche/problem & \begin{tabular}[c]{@{}l@{}}little is known about the \\ little is known about\end{tabular} &  \\
  %   1S3a & Indicating main purposs & \begin{tabular}[c]{@{}l@{}}the aim of this study was to evaluate the \\ the aim of this study was to invesigate \\ the aim of this study was to asscss \\ the aim of this study was to \\ the pupose of this study was to \\ the objecive of this study was to \\ of this study was to determine the \\ of this study was to invsigate \\ of this sudy was to determine \\ of this study was to evaluate \\ of this sudy was to compare \\ of this study was to examine \\ aim of the study was to \\ of this study was to \\ study was to evaluate the \\ study was to investigate the \\ of the study was to \\ the efficacy and safety of \\ study was to assess the \\ the safety and efficacy of \\ study was to determine the \\ of this study were to \\ we aimed to asses the\\ of this study is to\\ study was to compare the\\ to determine the effect of\\ to examine the association between\\ the aim of this\\ aim of this study\\ the purpose of this\\ study was to investigate\\ study was to evaluate\\ this study aimed to\\ study was to assess\\ we aimed to assess\\ of the study was\\ the study was to\\ study was to compare\\ the aim of the\\ our objective was to\\ our aim was to\\ we sought to determine\\ aimed to determine the\\ aimed to investigate the\\ was to determine whether\\ aimed to evaluate the\\ we aimed to evaluate\\ study aimed to determine\\ the objective of the\\ was to examine the\end{tabular} & \begin{tabular}[c]{@{}l@{}}for the treatment of \\ in patients with advanced \\ in the emergency departnent \\ in the treatment of\end{tabular} \\
  %   1S3b & Raising hypotheses & tested the hypothesis that & in patients with advanced \\
  %   M2 & Describing research process &  &  \\
  %   2S1 & Reporting on medical ethics review & study was approved by the &  \\
  %   2S2 & Explaining briefly research design & \begin{tabular}[c]{@{}l@{}}a retrospective cohort study of\\ a retrospective cohort study\\ we conducted a retrospective\\ a cross-sectional study\\ a secondary analysis of\end{tabular} &  \\
  %   2S3 & Describing subjects or data and their selection criteria & \begin{tabular}[c]{@{}l@{}}were included in the\\ were included in this\\ years or older with\\ data were collected from\end{tabular} & \begin{tabular}[c]{@{}l@{}}at the time of\\ on the basis of\end{tabular} \\
  %   2S4 & Describing experimental procedure, such as interventions, examinations, etc. & \begin{tabular}[c]{@{}l@{}}patients were randomly assigned to receive\\ were randomly assigned to receive\\ patients were randomly assigned to\\ were randomly assigned to\\ were randomized to receive\\ were masked to treatment\\ patients were randomized to\\ we randomly assigned patients\\ masked to treatment allocation\\ participants were randomly assigned\\ masked to treatment assignment\end{tabular} & \begin{tabular}[c]{@{}l@{}}all patients who received at least one dose of \\ at the time of \\ on the basis of \\ at the end of\\ at a dose of\\ was assessed using the\\ was used to assess\\ were used to assess\\ were used to evaluate\end{tabular} \\
  %   2S5 & Describing main outcomes and their measures & \begin{tabular}[c]{@{}l@{}}the primary end point was the\\ the primary end point was\\ the primary outcome was the\\ the primary endpoint was the\\ the primary outcome measure was\\ was the proportion of patients\\ the primary outcome was\\ the primary endpoint was\\ outcome measure was the\\ in all patients who\\ secondary end points included\end{tabular} & \begin{tabular}[c]{@{}l@{}}all patients who received at least one dose of \\ at the time of \\ on the basis of \\ was used to identify\\ at the end of \\ at a dose of \\ was assessed using the\\ was used to assess \\ were used to assess\\ analysis was used to \\ analyses were performed to \\ were used to determine\\ were used to evaluate\end{tabular} \\
  %   2S6 & Describing data analysis methods & \begin{tabular}[c]{@{}l@{}}logistic regression was used to\\ models were used to\\ regression analysis was used\end{tabular} & \begin{tabular}[c]{@{}l@{}}was used to identify\\ analysis was used to\\ analyses were performed to\end{tabular} \\
  %   2S7 & Reporting on registration information & \begin{tabular}[c]{@{}l@{}}this trial is registered with\\ this study is registered with\end{tabular} & were used to determine \\
  %   M3 & Summarizing results &  &  \\
  %   3S1 & Explaining information on valid samples & \begin{tabular}[c]{@{}l@{}}patients were randomly assigned to\\ were randomly assigned to receive\\ patients were included in the\\ were included in the study\\ were included in the\\ included in the analysis\\ patients were enrolled and\end{tabular} & \begin{tabular}[c]{@{}l@{}}at a median follow-up of\\ in the placebo group\\ in the control group\\ a median follow-up of during the study period\end{tabular} \\
  %   3S2 & Illustrating overall observation or main results & \begin{tabular}[c]{@{}l@{}}there was no significant difference in\\ \\ there were no significant differences in\\ there were no significant differences between\\ there was no difference in\\ adverse events were reported in\\ there were no differences in\\ serious adverse events were reported\\ were more likely to have\\ did not differ significantly between\\ did not differ between the\\ was associated with a/an\\ were more likely to\\ there was no significant\\ there were no significant\\ were randomly assigned to\\ did not differ between\\ was not associated with\\ at the time of\\ was significantly associated with\\ was significantly higher in\\ area under the curve\\ the most common grade\\ were less likely to\\ events were reported in\\ compared with the control\\ were not significantly different between the two groups\\ was not significantly different\\ significantly higher in the\\ was found to be\\ was significantly lower in\\ an increased risk of\\ not significantly different between\\ was the most common\\ were not associated with\\ group compared with the\\ was associated with increased\\ associated with an increased\\ no statistically significant difference\\ were significantly higher in\\ significant difference in the\\ no significant difference between\\ significant differences between the\\ was independently associated with\\ were associated with a\\ no significant differences were\end{tabular} &  \\
  %   M4 & Drawing conclusions &  &  \\
  %   4S1 & Reiterating pivotal results & \begin{tabular}[c]{@{}l@{}}was associated with a\\ was not associated with\end{tabular} & \begin{tabular}[c]{@{}l@{}}is associated with a\\ an increased risk of \\ associated with an increased\end{tabular} \\
  %   4S2 & Indicating limitations &  &  \\
  %   4S3 & Stating the significance of the results & \begin{tabular}[c]{@{}l@{}}the results of this study\\ this is the first\\ these findings suggest that\\ our results suggest that\\ these results suggest that\\ can be used to\\ for the treatment of\\ has the potential to\\ our findings suggest that\end{tabular} & \begin{tabular}[c]{@{}l@{}}is associated with a\\ an increased risk of \\ associated with an increased\end{tabular} \\
  %   4S4 & Predicting future studies & \begin{tabular}[c]{@{}l@{}}further research is needed to\\ studies are needed to\end{tabular} &  \\
  
  % \end{longtable}
  
  }
\end{landscape}



\backmatter



	
\end{document}